\documentclass[11pt]{amsart}

% Encoding and fonts
\usepackage[T1]{fontenc}
\usepackage[utf8]{inputenc}
\usepackage{lmodern}

% Math packages
\usepackage{amsmath}
\usepackage{amssymb}
\usepackage{amsfonts}
\usepackage{amsthm}
\usepackage{mathtools}

% Layout and utilities
\usepackage{geometry}
\usepackage{enumitem}
\usepackage{booktabs}
\geometry{margin=1in}
\usepackage[hidelinks]{hyperref}

\numberwithin{equation}{section}
\setcounter{tocdepth}{3}
\setcounter{secnumdepth}{3}

% ---------------------------------------------------------
% Theorem-like environments
% ---------------------------------------------------------
\theoremstyle{plain}
\newtheorem{theorem}{Theorem}[section]
\newtheorem{proposition}[theorem]{Proposition}
\newtheorem{lemma}[theorem]{Lemma}
\newtheorem{corollary}[theorem]{Corollary}
\newtheorem{conjecture}[theorem]{Conjecture}

\theoremstyle{definition}
\newtheorem{definition}[theorem]{Definition}
\newtheorem{experiment}[theorem]{Experiment}
\newtheorem{heuristic}[theorem]{Heuristic}

\theoremstyle{remark}
\newtheorem{remark}[theorem]{Remark}
\newtheorem{warning}[theorem]{Warning}

\newcommand{\RH}{\mathrm{RH}}
\newcommand{\Li}{\lambda}
\newcommand{\E}{\mathsf E}
\newcommand{\Omat}{\mathsf O}
\newcommand{\Kmat}{\mathsf K}
\newcommand{\J}{\mathsf J}
\newcommand{\e}{\mathrm e}
\newcommand{\dd}{\,\mathrm d}
\newcommand{\diag}{\operatorname{diag}}
\newcommand{\rank}{\operatorname{rank}}
\newcommand{\spec}{\operatorname{spec}}
\newcommand{\Span}{\operatorname{span}}
\newcommand{\Tr}{\operatorname{Tr}}

\title[Reflection Compressions of the Li--Weil Kernel]{Even--Odd Reflection Compressions of the Li--Weil Kernel:\\
The Even-Family RH Route, Inverse Dynamics, and Determinant Geometry}

\author{V\'ictor Casa\~na-Giner}
\address{Bio Logic Crop Science S.L., E-46120 Alboraya (Valencia), Spain}
\email{vcasana@biologiccropscience.com}

\begin{document}

\begin{abstract}
We study the even and odd reflection compressions
\(\mathsf E_N=(\lambda_{m+n}-\lambda_{|m-n|})\) and
\(\mathsf O_N=(\lambda_{m+n-1}-\lambda_{|m-n|})\) obtained by reflection-antisymmetrizing
Lagarias's bilateral Li--Weil Gram kernel.  The underlying Li test class and Weil scalar product are
classical; the matrices \(\mathsf E_N,\mathsf O_N\) are formed here by the explicit
antisymmetric changes of basis stated below and are then analyzed in those coordinates.  Unconditionally they inherit Toeplitz--Hankel form, discrete wave equations,
rank-two displacement, exact Schur recursions, and a closed two-generator inverse dynamics.  Under
RH they become strictly positive sine Gram matrices and admit Chebyshev moment, Hankel,
Christoffel--Darboux, Jacobi, Cauchy--Binet, and determinantal-process descriptions.

A central point of the present version is that the odd family is \emph{not needed} for the implication
to the Riemann hypothesis.  The even family alone satisfies
\[
\RH\Longleftrightarrow \mathsf E_N\succeq0\ \text{for every }N
\Longleftrightarrow \mathsf E_N\succ0\ \text{for every }N.
\]
Thus the proof problem may be formulated entirely as propagation of the even Schur pivots
\(p_N=\lambda_{2N}-b_N^T\mathsf E_{N-1}^{-1}b_N\).  The local \(LDL^T\) dynamics already reduces
that propagation to one explicit \(2\times2\) cell inequality.

We therefore separate two proof routes.  Route A works directly with \(\mathsf E_N\): the local
reflection dynamics, the prime--archimedean normal form, and an exact integration-by-parts formula
reduce positivity to a signed prime-number-theorem discrepancy against an explicit structured
polynomial test.  Route B is organized around \(\mathsf E_N^{-1}\): inversion preserves exact
displacement rank two, nesting adds a rank-one update, the inverse is the coefficient matrix of the
Christoffel--Darboux reproducing kernel, and the unique Schur residual direction produces an
adaptive prime--continuum test.  Either the local-cell inequality or the adaptive inverse inequality,
proved for every \(N\), would imply RH.

Finally, the de Bruijn--Newman heat flow supplies a falsification filter for proof mechanisms.
Rodgers--Tao proved \(\Lambda\ge0\); hence every negative heat time lies outside the all-real-zero
region.  The algebraic reflection, displacement, Schur, and inverse identities survive the heat
deformation, whereas the Euler-product/prime representation does not.  Consequently no argument
using only the heat-flow-invariant matrix structure can supply the missing sign.  Within the present
framework the required zeta-specific input must enter through the arithmetic prime side.
All finite computations are explicitly separated from theorem statements.
\end{abstract}

\maketitle
\clearpage
\renewcommand{\contentsname}{Table of Contents}
\tableofcontents
\clearpage

\section{Introduction}
\label{sec:introduction}

The Riemann hypothesis asserts that the nontrivial zeros of the Riemann zeta-function lie on the critical line $\Re s=1/2$; standard accounts include \cite{Riemann1859,Edwards1974,Titchmarsh1986,Conrey2003,Bombieri2010}. Li's criterion replaces the location of the zeros by positivity of a scalar sequence $\{\lambda_n\}_{n\ge1}$ \cite{Li1997}. Bombieri and Lagarias connected Li's coefficients with the Guinand--Weil explicit formula and with positivity of Weil's distribution \cite{Weil1952,BombieriLagarias1999}. Lagarias later placed the Li coefficients inside an exact Hermitian Weil scalar product on a rational test-function class and proved, for the special functions
\begin{equation}
G_k(s):=1-\left(1-\frac1s\right)^k,
\qquad k\in\mathbb Z,
\label{eq:intro_Lagarias_test}
\end{equation}
the identity
\begin{equation}
\langle G_r,G_s\rangle_W
=\lambda_r+\lambda_{-s}-\lambda_{r-s}.
\label{eq:intro_Lagarias_Gram}
\end{equation}
For the Riemann zeta-function $\lambda_{-n}=\lambda_n$, so \eqref{eq:intro_Lagarias_Gram} is the bilateral Li--Weil kernel \cite[Theorem~3.1]{Lagarias2007}. Bombieri studied Weil's quadratic functional directly in a different analytic framework \cite{Bombieri2000}. Suzuki has developed further Hilbert-space and moment realizations of the same positivity framework \cite{Suzuki2023Norms,Suzuki2023Screw,Suzuki2025Hilbert}.

The present paper applies a specific finite reflection-antisymmetrization to the bilateral
Li--Weil kernel \eqref{eq:intro_Lagarias_Gram}, which is taken directly from
Lagarias~\cite[Theorem~3.1]{Lagarias2007}.  More precisely, the even test vectors are
\[
F_n^E=\frac{G_n-G_{-n}}{\sqrt2},
\]
and the odd vectors are the corresponding half-shifted reflections
\(F_n^O=(G_n-G_{1-n})/\sqrt2\).  Taking their Lagarias--Weil Gram matrices produces
\begin{align}
\E_N
&:=\left(\lambda_{m+n}-\lambda_{|m-n|}\right)_{m,n=1}^N,
\label{eq:intro_E_definition}\\
\Omat_N
&:=\left(\lambda_{m+n-1}-\lambda_{|m-n|}\right)_{m,n=1}^N.
\label{eq:intro_O_definition}
\end{align}
Thus the source object---Li's coefficients, the Weil form, and Lagarias's bilateral kernel---is
classical \cite{Li1997,BombieriLagarias1999,Lagarias2007}.  The transformation to the two
reflection matrices above, and the identities obtained from their native reflection coordinates, are
the objects studied in this paper.  They should not be confused with earlier matrix reformulations of
RH, including Bombieri's finite Weil forms, the Hankel/Toeplitz eigenvalue constructions of
Matiyasevich, the multiplicative matrices of Hilberdink, or Suzuki's screw- and moment-space
realizations \cite{Bombieri2000,Matiyasevich2007a,Matiyasevich2007b,Hilberdink2006,Hilberdink2015,Suzuki2023Screw}.
Their diagonal splitting is
\begin{equation}
(\E_N)_{nn}=\lambda_{2n},
\qquad
(\Omat_N)_{nn}=\lambda_{2n-1}.
\label{eq:intro_diagonals}
\end{equation}
Thus the pair contains the complete Li sequence while retaining nontrivial mixed entries.

The first structural point is exact and does not assume the Riemann hypothesis. If
\begin{equation}
F_n^E:=\frac{G_n-G_{-n}}{\sqrt2},
\qquad
F_n^O:=\frac{G_n-G_{1-n}}{\sqrt2},
\label{eq:intro_reflected_tests}
\end{equation}
then polarization of \eqref{eq:intro_Lagarias_Gram} gives
\begin{equation}
(\E_N)_{mn}=\langle F_m^E,F_n^E\rangle_W,
\qquad
(\Omat_N)_{mn}=\langle F_m^O,F_n^O\rangle_W.
\label{eq:intro_exact_Weil_Gram}
\end{equation}
Under $\RH$, if $\rho=1/2+i\gamma$ and
\begin{equation}
1-\frac1\rho=\e^{i\theta_\gamma},
\label{eq:intro_cayley}
\end{equation}
then these exact Weil Gram matrices become the positive sine Gram matrices
\begin{align}
(\E_N)_{mn}
&=4\sum_{\gamma>0}\sin(m\theta_\gamma)\sin(n\theta_\gamma),
\label{eq:intro_gram_E}\\
(\Omat_N)_{mn}
&=4\sum_{\gamma>0}\sin\bigl((m-\tfrac12)\theta_\gamma\bigr)
\sin\bigl((n-\tfrac12)\theta_\gamma\bigr),
\label{eq:intro_gram_O}
\end{align}
where zeros are counted with multiplicity. These identities imply strict positivity of every finite truncation. Conversely, positivity of all truncations recovers the even and odd Li inequalities from the diagonals and therefore recovers Li's criterion.

The remainder of the paper develops consequences of this reflection pair. Native coordinates give Toeplitz--Hankel matrices and low-rank displacement identities. Under $\RH$, the zero-side sine representations become moment Gram matrices after Chebyshev changes of basis; the two measures are linked by a Christoffel transform at $-1$. Their determinants have positive Cauchy--Binet--Vandermonde expansions. We first prove an abstract endpoint-exponent law: for a positive discrete Gram measure whose support satisfies $1-x(t)\asymp t^{-\beta}$ and only coarse logarithmic node and separation conditions, the determinant obeys $\log D_N=-(\beta+o(1))N^2\log N$. The Li--Cayley transform has $\beta=2$, hence
\begin{equation}
\log\det\E_N=-(2+o(1))N^2\log N,
\qquad
\log\det\Omat_N=-(2+o(1))N^2\log N.
\label{eq:intro_determinant_law}
\end{equation}
The coefficient $-2$ is therefore an endpoint-geometry constant, not a record of the logarithmic refinement in the Riemann--von Mangoldt density. The generic occurrence of $N^2\log N$ decay in discrete Hankel determinants has precedents in the Dirichlet-series setting of Monien and Haynes--Zudilin \cite{Monien2009,HaynesZudilin2015}; the distinction here is the Li--Weil reflection measure and the endpoint-exponent mechanism. At this order no pair-correlation input enters. Schur complements are residual Gram energies, and inverse matrices represent Christoffel--\allowbreak Darboux kernels. Finally, the associated \emph{infinite} Jacobi operators are unitarily equivalent to multiplication by
\begin{equation}
x_\gamma:=\frac{4\gamma^2-1}{4\gamma^2+1}
\label{eq:intro_xgamma}
\end{equation}
on an atomic $L^2$ space. Their spectrum therefore contains the transformed distinct ordinates and the accumulation point $1$. This is an infinite moment reconstruction from the full Li data; the eigenvalues of a finite Jacobi truncation are Gaussian-quadrature nodes, not exact zeta-zero ordinates.


A complementary route, developed in Section~\ref{sec:arch_prime_laguerre}, separates the completed zeta function into its archimedean factor and the analytic arithmetic factor $(s-1)\zeta(s)$.  In the Cayley variable $s=1/(1-z)$ the latter has coefficients $Z_k$ with $\lambda_k=A_k+Z_k$.  The arithmetic coefficients admit an exact regularized von Mangoldt--Laguerre formula, and reflection removes the prime main term so that the prime matrix pairs a Laguerre-difference kernel directly with $\psi(t)-t$.  The associated scalar root-growth criterion $\RH\Longleftrightarrow\limsup |Z_k|^{1/k}\le1$ is useful diagnostically but belongs to the classical Keiper--Li singularity mechanism after factorization.  The extra structure specific to the reflection compression appears only before scalarization: quadratic tests become Fourier transforms of nonnegative sine squares, each prime atom has rank-two displacement/Hankel structure, and first differences expose a contractive matrix-valued cosine family together with an explicit positive decomposition of the archimedean matrix up to one scalar negative constant.

Three finite-$N$ structures specific to the reflection pair are central below.  The bordering columns $\mathsf b^{(N)}_m=\lambda_{m+N}-\lambda_{|m-N|}$ satisfy, unconditionally, the same discrete wave equation in $N$ as the rows do in $m$; together with the rank-two displacement this closes the Schur update and yields an exact local recurrence for consecutive pivot ratios in terms of the first two subdiagonals of the $LDL^T$ factor.  Independently, the Christoffel bridge of Corollary~\ref{cor:determinant_bridge} normalizes as $\det\E_N=4^NF_N\det\Omat_N$ and factorises the ratio of the even and odd pivot sequences as $\kappa_N=4F_N/F_{N-1}$.  All numerical positivity statements are based on the arithmetic reconstruction of the Li coefficients in Subsection~\ref{subsec:arithmetic_reconstruction}, which uses Stieltjes constants and no zero of $\zeta$.

A point not visible from any single Vandermonde summand is how the \emph{full} Cauchy--Binet determinant distributes its mass among zero subsets.  Section~\ref{sec:zero_leverage} resolves this at finite $N$.  After normalizing the nonnegative Cauchy--Binet terms one obtains a canonical probability measure on $N$-element subsets of positive zero ordinates.  It is a rank-$N$ projection determinantal process, and its marginal inclusion probabilities are
\begin{equation}
\ell_{N,\gamma}^{E,O}
=(a_{N,\gamma}^{E,O})^T(\E_N\text{ or }\Omat_N)^{-1}a_{N,\gamma}^{E,O}.
\label{eq:intro_zero_leverage}
\end{equation}
These quantities have a direct matrix meaning: removing the atom at $\gamma$ changes the determinant by the exact factor $1-\ell_{N,\gamma}^{E,O}$.  Thus the question of which zeros actually control the determinant is encoded by the matrix itself, without replacing the full Cauchy--Binet sum by one selected Vandermonde term.

\begin{remark}[Exact source ledger]
\label{rem:status_provenance}
The construction starts from published material, and the source is recorded explicitly rather than
absorbed into the notation.
\begin{enumerate}[label=(\alph*)]
\item The coefficients \(\lambda_n\) and their positivity criterion are Li's \cite{Li1997}; the
closely related power-series viewpoint predates Li in Keiper \cite{Keiper1992}.
\item The underlying explicit-formula positivity principle is Weil's quadratic criterion
\cite{Weil1952}.  The general multiset extension of Li's criterion, its relation with the
Guinand--Weil explicit formula, and the inverse-Mellin Li test functions are taken from
Bombieri--Lagarias \cite{BombieriLagarias1999}.
\item The rational test functions
\(G_n(s)=1-(1-1/s)^n\), the Li test class, and especially the exact identity
\[
\langle G_n,G_m\rangle_W=\lambda_n+\lambda_{-m}-\lambda_{n-m}
\]
are taken directly from Lagarias \cite[Section~3, Theorem~3.1]{Lagarias2007}.  In other words,
the parent kernel used in this paper is copied, up to notation, from Lagarias's Theorem~3.1; it is
not derived here.  This theorem is the immediate source from which the matrices in this paper are formed.
\item The even matrix is obtained by the elementary change of test basis
\[
G_n\longmapsto F_n^E:=\frac{G_n-G_{-n}}{\sqrt2}
\]
and then taking the Lagarias--Weil Gram matrix.  The odd matrix uses
\(F_n^O=(G_n-G_{1-n})/\sqrt2\).  Thus \(\E_N\) and \(\Omat_N\) are not independent replacements
for the Li--Weil theory: they are finite Gram matrices produced by these explicit linear
transformations of Lagarias's published kernel.  The present paper claims only the identities proved
for these transformed coordinates; it does not claim priority for the underlying Li/Weil kernel,
and it does not assert that the elementary antisymmetrization itself has never appeared elsewhere.
\item The associated Laguerre polynomial in the Li test function is also not new.  Lagarias records
that the inverse-Mellin functions were explicitly determined in Bombieri--Lagarias and identifies
the corresponding polynomial with the associated Laguerre family; see
\cite[Appendix~A]{Lagarias2007} and \cite{BombieriLagarias1999}.  Coffey developed an explicit
theta--Laguerre formulation of the Li/Keiper constants \cite{Coffey2007ThetaLaguerre}.
The scalar decomposition $\lambda_n=A_n+Z_n$ used below is obtained directly by splitting
$\log\xi(s)$ into the archimedean factor and $\log((s-1)\zeta(s))$ before taking Cayley
coefficients; no priority is claimed for that scalar factorization.  What is derived here is the
cancellation and matrix structure that appear after taking the reflection differences.
\item Cauchy--Binet, Hankel moment theory, Christoffel transforms, Christoffel--Darboux kernels,
Jacobi matrices, Darboux factorization, and projection determinantal processes are standard tools;
see \cite{Chihara1978,StahlTotik1992,HoughKrishnapurPeresVirag2009}.  They are not claimed as
new mechanisms here.
\end{enumerate}
Other approaches include Bombieri's work on Weil's quadratic functional and Roesler's eigenvalue
problems \cite{Bombieri2000,Roesler1986,Roesler1987}; Matiyasevich's Taylor-coefficient
Toeplitz/Hankel matrices \cite{Matiyasevich2007a,Matiyasevich2007b}; Hilberdink's multiplicative
Toeplitz matrices \cite{Hilberdink2006,Hilberdink2015}; and Suzuki's screw/model-space
constructions \cite{Suzuki2023Norms,Suzuki2023Screw}.
\end{remark}

\section{The xi-function and Li coefficients}
\label{sec:li_background}

Define the completed zeta-function
\begin{equation}
\xi(s)=\frac12 s(s-1)\pi^{-s/2}\Gamma\!\left(\frac{s}{2}\right)\zeta(s).
\label{eq:xi_definition}
\end{equation}
It is entire of order one and satisfies
\begin{equation}
\xi(s)=\xi(1-s)=\overline{\xi(\overline{s})}.
\label{eq:xi_symmetry}
\end{equation}
The standard analytic theory may be found in \cite{Edwards1974,Titchmarsh1986,Bombieri2010}.

\begin{definition}[Li coefficients]
\label{def:li_coefficients}
For $n\ge1$, define
\begin{equation}
\lambda_n=
\frac{1}{(n-1)!}
\left.
\frac{\dd^n}{\dd s^n}
\left[s^{n-1}\log\xi(s)\right]
\right|_{s=1},
\qquad \lambda_0:=0.
\label{eq:li_derivative_definition}
\end{equation}
Equivalently, with the symmetric limiting convention over nontrivial zeros,
\begin{equation}
\lambda_n=
\sum_{\rho}
\left[1-\left(1-\frac1\rho\right)^n\right].
\label{eq:li_zero_definition}
\end{equation}
\end{definition}

These numbers are also called the Keiper--Li coefficients: the power-series expansion of $\log\xi$ underlying \eqref{eq:li_derivative_definition} and its high-precision evaluation were studied by Keiper \cite{Keiper1992} before Li's positivity criterion appeared.

For the zeta xi-function the zero symmetry implies an integer-index extension
\begin{equation}
\lambda_{-n}=\lambda_n
\qquad(n\ge1),
\label{eq:li_negative_symmetry}
\end{equation}
as recorded, for example, in \cite[Eq.~(1.2)]{Lagarias2007}. This identity is unconditional and will be used whenever differences of integer indices occur.

\begin{theorem}[Li]
\label{thm:li_criterion}
The Riemann hypothesis is equivalent to
\begin{equation}
\lambda_n\ge0
\qquad\text{for every }n\ge1.
\label{eq:li_criterion}
\end{equation}
\end{theorem}

\begin{proof}
This is Li's positivity criterion \cite{Li1997}; see also \cite{BombieriLagarias1999,Lagarias2007}.
\end{proof}

Define
\begin{equation}
R_\xi:=\min\!\left\{1,\inf_{\rho}\left|1-\frac1\rho\right|\right\}.
\label{eq:li_generating_radius}
\end{equation}
The Li generating identity is the formal power-series identity
\begin{equation}
\Lambda(z):=\sum_{n\ge1}\lambda_n z^n
=
\frac{z}{(1-z)^2}
\frac{\xi'}{\xi}\!\left(\frac{1}{1-z}\right).
\label{eq:li_generating_function}
\end{equation}
Analytically, \eqref{eq:li_generating_function} holds in the disk $|z|<R_\xi$. Under $\RH$, every nontrivial zero satisfies $|1-1/\rho|=1$, and therefore $R_\xi=1$. If $\RH$ fails, a zero with real part $>1/2$ produces a singularity with $|1-1/\rho|<1$, so the unconditional radius must not be stated as $1$.

Under $\RH$ we may parameterize each zero in the upper half-plane by
\begin{equation}
\rho=\frac12+i\gamma,
\qquad
1-\frac1\rho=\e^{i\theta_\gamma},
\qquad 0<\theta_\gamma<\pi.
\label{eq:theta_parameterization}
\end{equation}
A direct calculation gives
\begin{align}
\cos\theta_\gamma
&=\frac{4\gamma^2-1}{4\gamma^2+1},
\label{eq:theta_cos}\\
\sin\theta_\gamma
&=\frac{4\gamma}{4\gamma^2+1}.
\label{eq:theta_sin}
\end{align}
Pairing conjugate zeros in \eqref{eq:li_zero_definition} yields the absolutely convergent representation
\begin{equation}
\lambda_n
=2\sum_{\gamma>0}
\left(1-\cos(n\theta_\gamma)\right).
\label{eq:li_cosine_representation}
\end{equation}
For fixed $n$, the summands are $O_n(\gamma^{-2})$.

\section{Definition of the reflection matrices}
\label{sec:definitions}

\begin{definition}[Even and odd reflection matrices]
\label{def:reflection_matrices}
For $N\ge1$, define
\begin{align}
\E_N
&:=\left(e_{mn}\right)_{m,n=1}^N,
&
e_{mn}&:=\lambda_{m+n}-\lambda_{|m-n|},
\label{eq:E_definition}\\
\Omat_N
&:=\left(o_{mn}\right)_{m,n=1}^N,
&
o_{mn}&:=\lambda_{m+n-1}-\lambda_{|m-n|}.
\label{eq:O_definition}
\end{align}
\end{definition}

Both are real symmetric. Their first $4\times4$ truncations are
\begin{equation}
\E_4=
\begin{pmatrix}
\lambda_2&\lambda_3-\lambda_1&\lambda_4-\lambda_2&\lambda_5-\lambda_3\\
\lambda_3-\lambda_1&\lambda_4&\lambda_5-\lambda_1&\lambda_6-\lambda_2\\
\lambda_4-\lambda_2&\lambda_5-\lambda_1&\lambda_6&\lambda_7-\lambda_1\\
\lambda_5-\lambda_3&\lambda_6-\lambda_2&\lambda_7-\lambda_1&\lambda_8
\end{pmatrix},
\label{eq:E4}
\end{equation}
and
\begin{equation}
\Omat_4=
\begin{pmatrix}
\lambda_1&\lambda_2-\lambda_1&\lambda_3-\lambda_2&\lambda_4-\lambda_3\\
\lambda_2-\lambda_1&\lambda_3&\lambda_4-\lambda_1&\lambda_5-\lambda_2\\
\lambda_3-\lambda_2&\lambda_4-\lambda_1&\lambda_5&\lambda_6-\lambda_1\\
\lambda_4-\lambda_3&\lambda_5-\lambda_2&\lambda_6-\lambda_1&\lambda_7
\end{pmatrix}.
\label{eq:O4}
\end{equation}

\begin{proposition}[Toeplitz--Hankel decomposition]
\label{prop:toeplitz_hankel}
Let
\begin{equation}
\mathsf T_N:=\bigl(\lambda_{|m-n|}\bigr)_{m,n=1}^N,
\label{eq:T_definition}
\end{equation}
\begin{equation}
\mathsf H_N^{(0)}:=\bigl(\lambda_{m+n}\bigr)_{m,n=1}^N,
\qquad
\mathsf H_N^{(-1)}:=\bigl(\lambda_{m+n-1}\bigr)_{m,n=1}^N.
\label{eq:H_definitions}
\end{equation}
Then
\begin{equation}
\E_N=\mathsf H_N^{(0)}-\mathsf T_N,
\qquad
\Omat_N=\mathsf H_N^{(-1)}-\mathsf T_N.
\label{eq:toeplitz_hankel_decomposition}
\end{equation}
Moreover,
\begin{equation}
\E_N-\Omat_N
=\bigl(\lambda_{m+n}-\lambda_{m+n-1}\bigr)_{m,n=1}^N,
\label{eq:E_minus_O}
\end{equation}
which is a pure Hankel matrix.
\end{proposition}

\begin{proof}
All statements follow directly from \eqref{eq:E_definition} and \eqref{eq:O_definition}.
\end{proof}

\section{Algebraic reconstruction and information content}
\label{sec:algebraic_reconstruction}

The two arrays contain substantial algebraic redundancy.  This can be described exactly without assuming $\RH$.

\begin{proposition}[Reconstruction from the infinite first rows]
\label{prop:first_row_reconstruction}
Since
\begin{equation}
o_{1j}=\lambda_j-\lambda_{j-1},
\qquad j\ge1,
\label{eq:odd_first_row_differences}
\end{equation}
one has
\begin{equation}
\lambda_n=\sum_{j=1}^{n}o_{1j}
\qquad(n\ge1),
\label{eq:lambda_from_O_first_row}
\end{equation}
and, for all $m,n\ge1$,
\begin{equation}
o_{mn}
=\sum_{j=|m-n|+1}^{m+n-1}o_{1j}.
\label{eq:O_interval_sum_reconstruction}
\end{equation}
If $r:=\min(m,n)$, then
\begin{equation}
e_{mn}
=\sum_{j=0}^{r-1}
e_{1,\,|m-n|+1+2j}.
\label{eq:E_step_two_reconstruction}
\end{equation}
Finally,
\begin{equation}
e_{mn}
=\sum_{j=|m-n|+1}^{m+n}o_{1j}.
\label{eq:E_from_O_first_row}
\end{equation}
Thus the semi-infinite odd array determines the complete Li sequence and the semi-infinite even array.
\end{proposition}

\begin{proof}
Equation \eqref{eq:lambda_from_O_first_row} is telescoping.  Substituting it into the definition of $o_{mn}$ gives \eqref{eq:O_interval_sum_reconstruction}.  Since
\begin{equation}
e_{1k}=\lambda_{k+1}-\lambda_{k-1},
\label{eq:E_first_row_step_two}
\end{equation}
the right side of \eqref{eq:E_step_two_reconstruction} telescopes from $\lambda_{|m-n|}$ to $\lambda_{m+n}$.  Equation \eqref{eq:E_from_O_first_row} is the same telescoping identity written in the unit-step differences $o_{1j}$.
\end{proof}

\begin{remark}[Finite truncations versus infinite first rows]
\label{rem:finite_first_row_warning}
Proposition \ref{prop:first_row_reconstruction} concerns the first row of the \emph{semi-infinite} arrays.  The leading block $\Omat_N$ is not determined by its own $N$ displayed first-row entries: reconstructing all of $\Omat_N$ from \eqref{eq:O_interval_sum_reconstruction} requires the prefix $o_{1j}$ through $j=2N-1$.  Likewise, reconstructing $\E_N$ from its infinite first row requires $e_{1j}$ through $j=2N-1$.  Reconstructing $\E_N$ from the odd first row by \eqref{eq:E_from_O_first_row} requires $o_{1j}$ through $j=2N$.
\end{remark}

\begin{proposition}[Injectivity of the odd array and one-dimensional kernel of the even array]
\label{prop:information_kernel}
On real sequences $(\lambda_n)_{n\ge0}$ with $\lambda_0=0$, the map
\begin{equation}
(\lambda_n)_{n\ge0}\longmapsto (o_{mn})_{m,n\ge1}
\label{eq:O_sequence_map}
\end{equation}
is injective.  The kernel of
\begin{equation}
(\lambda_n)_{n\ge0}\longmapsto (e_{mn})_{m,n\ge1}
\label{eq:E_sequence_map}
\end{equation}
consists exactly of the one-parameter family
\begin{equation}
\lambda_{2k}=0,
\qquad
\lambda_{2k+1}=c
\qquad(k\ge0),
\label{eq:E_kernel_sequences}
\end{equation}
with $c\in\mathbb R$.  Consequently $(\E,\lambda_1)$ determines the whole Li sequence and hence $\Omat$.
\end{proposition}

\begin{proof}
If the odd array vanishes, then its first row gives $\lambda_j-\lambda_{j-1}=0$ for every $j$, and $\lambda_0=0$ gives $\lambda_j=0$.  Thus the odd map is injective.

If the even array vanishes, its diagonal gives $\lambda_{2n}=0$ for all $n\ge1$.  Its entries $e_{1,2k}=0$ give $\lambda_{2k+1}=\lambda_{2k-1}$, so every odd term equals $\lambda_1$.  Conversely, a sequence of the form \eqref{eq:E_kernel_sequences} gives $e_{mn}=0$ because $m+n$ and $m-n$ have the same parity.  This proves the kernel statement.
\end{proof}

\begin{proposition}[Reflection boundary conditions]
\label{prop:reflection_boundary_conditions}
Using the even extension $\lambda_{-n}=\lambda_n$, the bilateral kernels satisfy
\begin{align}
e(-m,n)&=-e(m,n),
\label{eq:E_reflection_oddness}\\
o(1-m,n)&=-o(m,n).
\label{eq:O_reflection_oddness}
\end{align}
Hence
\begin{equation}
e(0,n)=0,
\qquad
o(0,n)=-o(1,n).
\label{eq:reflection_boundary_values}
\end{equation}
Together with symmetry in $m,n$ and the discrete wave equation of Section \ref{sec:displacement}, these are the natural reflection boundary conditions at $0$ and $1/2$.
\end{proposition}

\begin{proof}
The identities follow by direct substitution in
$e(m,n)=\lambda_{m+n}-\lambda_{m-n}$ and
$o(m,n)=\lambda_{m+n-1}-\lambda_{m-n}$, using $\lambda_{-k}=\lambda_k$.
\end{proof}

\section{Exact Weil realization and critical-line Gram representations}
\label{sec:gram_criterion}

Lagarias defines on the Li test-function class a Hermitian Weil scalar product for which the functions \eqref{eq:intro_Lagarias_test} satisfy \eqref{eq:intro_Lagarias_Gram}; see \cite[Section~3, Theorem~3.1]{Lagarias2007}. In particular, for the Riemann zeta-function,
\begin{equation}
\langle G_r,G_s\rangle_W
=\lambda_r+\lambda_s-\lambda_{|r-s|},
\qquad r,s\in\mathbb Z,
\label{eq:Lagarias_kernel_zeta}
\end{equation}
where \eqref{eq:li_negative_symmetry} has been used.

\begin{theorem}[Exact reflected Weil Gram matrices]
\label{thm:exact_Weil_reflection}
For every $m,n\ge1$, without assuming $\RH$,
\begin{align}
e_{mn}
&=\langle F_m^E,F_n^E\rangle_W,
\label{eq:E_exact_Weil_Gram}\\
o_{mn}
&=\langle F_m^O,F_n^O\rangle_W,
\label{eq:O_exact_Weil_Gram}
\end{align}
where $F_n^E,F_n^O$ are given by \eqref{eq:intro_reflected_tests}.
\end{theorem}

\begin{proof}
Using \eqref{eq:Lagarias_kernel_zeta} and bilinearity/polarization,
\begin{align}
\langle F_m^E,F_n^E\rangle_W
&=\frac12\bigl[
K(m,n)-K(m,-n)-K(-m,n)+K(-m,-n)
\bigr]
\label{eq:E_exact_polarization}\\
&=\lambda_{m+n}-\lambda_{|m-n|}=e_{mn},
\label{eq:E_exact_polarization_result}
\end{align}
where $K(r,s):=\lambda_r+\lambda_s-\lambda_{|r-s|}$. Similarly,
\begin{align}
\langle F_m^O,F_n^O\rangle_W
&=\frac12\bigl[
K(m,n)-K(m,1-n)-K(1-m,n)+K(1-m,1-n)
\bigr]
\label{eq:O_exact_polarization}\\
&=\lambda_{m+n-1}-\lambda_{|m-n|}=o_{mn}.
\label{eq:O_exact_polarization_result}
\end{align}
\end{proof}

The same reflection is transparent in the Cayley coordinate
\begin{equation}
q(s):=1-\frac1s.
\label{eq:q_coordinate}
\end{equation}
Indeed,
\begin{align}
F_n^E(s)
&=\frac{q(s)^{-n}-q(s)^n}{\sqrt2},
\label{eq:FE_q_form}\\
F_n^O(s)
&=\frac{q(s)^{1-n}-q(s)^n}{\sqrt2}.
\label{eq:FO_q_form}
\end{align}

\begin{theorem}[Integer and half-integer sine Gram representations]
\label{thm:gram_representations}
Assume $\RH$. Enumerate the positive ordinates of the nontrivial zeros with multiplicity. Then for every $m,n\ge1$,
\begin{align}
e_{mn}
&=4\sum_{\gamma>0}
\sin(m\theta_\gamma)\sin(n\theta_\gamma),
\label{eq:E_gram}\\
o_{mn}
&=4\sum_{\gamma>0}
\sin\bigl((m-\tfrac12)\theta_\gamma\bigr)
\sin\bigl((n-\tfrac12)\theta_\gamma\bigr).
\label{eq:O_gram}
\end{align}
The two series converge absolutely for fixed $m,n$.
\end{theorem}

\begin{proof}
Under $\RH$, $q(\rho)=\e^{i\theta_\gamma}$ for $\rho=1/2+i\gamma$. Equations \eqref{eq:FE_q_form}--\eqref{eq:FO_q_form} become
\begin{align}
F_n^E(\rho)
&=-\sqrt2\,i\sin(n\theta_\gamma),
\label{eq:FE_on_zero}\\
F_n^O(\rho)
&=-\sqrt2\,i\e^{i\theta_\gamma/2}
\sin\bigl((n-\tfrac12)\theta_\gamma\bigr).
\label{eq:FO_on_zero}
\end{align}
For the numerical normalization, it is safest to use the already established Li cosine formula \eqref{eq:li_cosine_representation}. Subtracting the terms with indices $m+n$ and $|m-n|$ and applying
\begin{equation}
\cos((m-n)\theta)-\cos((m+n)\theta)
=2\sin(m\theta)\sin(n\theta)
\label{eq:integer_sine_identity}
\end{equation}
gives \eqref{eq:E_gram}. Likewise,
\begin{equation}
\cos((m-n)\theta)-\cos((m+n-1)\theta)
=2\sin\bigl((m-\tfrac12)\theta\bigr)
\sin\bigl((n-\tfrac12)\theta\bigr)
\label{eq:half_integer_sine_identity}
\end{equation}
gives \eqref{eq:O_gram}. Since $\theta_\gamma=O(\gamma^{-1})$, every summand is $O_{m,n}(\gamma^{-2})$, and the zero-counting estimate gives absolute convergence.
\end{proof}

\begin{remark}[Negative definite sequences and screw functions]
\label{rem:negative_definite_context}
Representation \eqref{eq:li_cosine_representation} says that, under $\RH$, the extended sequence $\lambda$ is a \emph{negative definite} function on $\mathbb Z$ in the sense of harmonic analysis on groups: it is real, even, vanishes at $0$, and has the L\'evy--Khinchine form $\int(1-\cos n\theta)\dd\sigma(\theta)$ with $\sigma$ the atomic measure carried by the angles $\theta_\gamma$. For any such function the kernel $\lambda_r+\lambda_s-\lambda_{r-s}$ is positive definite, and the passage from that kernel to the sine Gram matrices of Theorem \ref{thm:gram_representations} is the standard consequence of the L\'evy--Khinchine representation \cite{Schoenberg1938,BergChristensenRessel1984}. Conversely, a negative definite $\lambda$ with $\lambda_0=0$ satisfies $\lambda_n\ge0$ for all $n$, so Lagarias's kernel identity \eqref{eq:Lagarias_kernel_zeta} amounts to the statement that $\RH$ is equivalent to negative definiteness of the Li sequence. The continuous counterpart of this circle of ideas is Krein's theory of screw functions, developed for the Riemann zeta-function by Suzuki \cite{Suzuki2023Screw,SuzukiScrewLine}. The positivity mechanism exploited below is therefore classical; what is specific here is the pair of reflection compressions of Definition \ref{def:reflection_matrices} and the structure theory derived from it.
\end{remark}

\begin{theorem}[Li--Weil reflection criterion]
\label{thm:main_criterion}
The following are equivalent:
\begin{enumerate}[label=(\roman*)]
\item $\RH$;
\item $\E_N\succeq0$ and $\Omat_N\succeq0$ for every $N\ge1$;
\item $\E_N\succ0$ and $\Omat_N\succ0$ for every $N\ge1$;
\item $\det\E_N>0$ and $\det\Omat_N>0$ for every $N\ge1$.
\end{enumerate}
\end{theorem}

\begin{proof}
Assume $\RH$. Theorem \ref{thm:gram_representations} gives positive semidefiniteness. If $c^T\E_Nc=0$, then
\begin{equation}
\sum_{m=1}^{N}c_m\sin(m\theta_\gamma)=0
\qquad\text{for every distinct positive zero ordinate }\gamma.
\label{eq:E_zero_polynomial}
\end{equation}
There are infinitely many distinct ordinates and $\theta_\gamma\to0$. The finite trigonometric polynomial in \eqref{eq:E_zero_polynomial} is analytic in $\theta$ and therefore must vanish identically; hence $c=0$. The half-integer system is handled in the same way. Thus (i) implies (iii), and (iii) implies (ii).

Assume (ii). For $N\ge n$ the diagonal entries satisfy
\begin{equation}
\lambda_{2n}=(\E_N)_{nn}\ge0,
\qquad
\lambda_{2n-1}=(\Omat_N)_{nn}\ge0.
\label{eq:recover_li_from_diagonal}
\end{equation}
Hence every Li coefficient is nonnegative and Theorem \ref{thm:li_criterion} gives $\RH$. Finally, because the families are nested leading principal truncations, Sylvester's criterion gives the equivalence of (iii) and (iv).
\end{proof}

\begin{theorem}[The even reflection family alone is an RH criterion]
\label{thm:even_family_alone}
The following are equivalent:
\begin{enumerate}[label=(\roman*)]
\item \(\RH\);
\item \(\E_N\succeq0\) for every \(N\ge1\);
\item \(\E_N\succ0\) for every \(N\ge1\).
\end{enumerate}
In particular, the odd family is not required for a proof of \(\RH\).
\end{theorem}

\begin{proof}
The implication \((i)\Rightarrow(iii)\Rightarrow(ii)\) is already contained in
Theorem~\ref{thm:main_criterion}.  Assume now that \(\E_N\succeq0\) for every \(N\).  Then
\[
\lambda_{2N}=(\E_N)_{NN}\ge0
\qquad(N\ge1).
\]
We reduce the converse to the general multiset theorem of Bombieri--Lagarias
\cite[Theorem~1 and its symmetric corollary]{BombieriLagarias1999}; this reduction is not a
new Li-type criterion.

For every nontrivial zero \(\rho\), put
\[
q_\rho:=1-\frac1\rho.
\]
Zeros with \(q_\rho^2=1\) are already on the critical line and contribute zero to every even Li
coefficient, so omit them from the following transformed multiset.  For the remaining zeros define
\begin{equation}
\sigma_\rho:=\frac1{1-q_\rho^2}=\frac{\rho^2}{2\rho-1}.
\label{eq:even_BL_transform}
\end{equation}
Then
\begin{equation}
1-\frac1{\sigma_\rho}=q_\rho^2,
\qquad
\lambda_{2n}
=\sum_\rho\left[1-\left(1-\frac1{\sigma_\rho}\right)^n\right].
\label{eq:even_Li_as_BL_sum}
\end{equation}
Moreover \(\sigma_\rho=\rho/2+O(1)\) as \(|\rho|\to\infty\), so the Bombieri--Lagarias
summability hypothesis is inherited from the zeta zero multiset.  The symmetries
\(\rho\mapsto\bar\rho\) and \(\rho\mapsto1-\rho\) become
\(\sigma\mapsto\bar\sigma\) and \(\sigma\mapsto1-\sigma\), respectively.  Hence the symmetric
Bombieri--Lagarias criterion applied to the transformed multiset gives
\[
\lambda_{2n}\ge0\ \text{for every }n
\quad\Longrightarrow\quad
\Re\sigma_\rho=\frac12\ \text{for every }\rho.
\]
But
\[
\Re\sigma_\rho=\frac12
\quad\Longleftrightarrow\quad
\left|1-\frac1{\sigma_\rho}\right|=1
\quad\Longleftrightarrow\quad
|q_\rho|=1
\quad\Longleftrightarrow\quad
\Re\rho=\frac12.
\]
Thus every nontrivial zero lies on the critical line and \(\RH\) follows.
\end{proof}

\begin{remark}[Source of the even-only converse]
\label{rem:even_only_source}
The converse in Theorem~\ref{thm:even_family_alone} is an application of the general
Bombieri--Lagarias multiset criterion after the elementary quadratic Cayley transform
\eqref{eq:even_BL_transform}.  The paper therefore claims no priority for the underlying
positivity criterion.  What is used later is only that the matrix diagonal exposes precisely these
even Bombieri--Lagarias sums.
\end{remark}

\begin{remark}[Role of the odd family]
\label{rem:odd_not_needed_RH}
The odd matrices remain useful because they supply the Christoffel/Darboux bridge, an independent
pivot sequence, and a second moment system.  None of those structures is logically necessary for the
implication to \(\RH\): throughout the proof routes isolated below, the target may be taken to be
\(\E_N\succ0\) for every \(N\).
\end{remark}

\begin{corollary}[Full rank]
\label{cor:full_rank}
Under $\RH$,
\begin{equation}
\rank \E_N=\rank \Omat_N=N
\qquad(N\ge1).
\label{eq:full_rank}
\end{equation}
\end{corollary}

\begin{proof}
Under \(\RH\), Theorem~\ref{thm:main_criterion} gives
\(\E_N\succ0\) and \(\Omat_N\succ0\).  A positive definite \(N\times N\)
matrix is nonsingular, hence has rank \(N\).
\end{proof}

\begin{corollary}[All principal minors]
\label{cor:all_principal_minors}
Under $\RH$, every principal minor of $\E_N$ and $\Omat_N$ is strictly positive.
\end{corollary}

\begin{proof}
A principal submatrix is the Gram matrix of a finite set of distinct integer, respectively half-integer, sine modes. The same analytic-zero argument as in Theorem \ref{thm:main_criterion} gives linear independence.
\end{proof}

\section{Reflection geometry inside the Li test class}
\label{sec:reflection_compression}

The exact matrix identities above can be phrased as reflections of Lagarias's bilateral Li-class Gram kernel. Define
\begin{equation}
\Kmat(r,s):=\langle G_r,G_s\rangle_W
=\lambda_r+\lambda_s-\lambda_{|r-s|},
\qquad r,s\in\mathbb Z.
\label{eq:bilateral_kernel}
\end{equation}
Let
\begin{equation}
R_0(k)=-k,
\qquad
R_{1/2}(k)=1-k.
\label{eq:two_reflections}
\end{equation}
Then
\begin{equation}
R_{1/2}R_0(k)=k+1,
\label{eq:reflection_translation}
\end{equation}
so the two reflections generate the infinite dihedral action on the integer label set.

If $u_m=e_m-e_{-m}$ and $v_m=e_m-e_{1-m}$ in the finitely supported bilateral coordinate space, then
\begin{align}
\frac12\,u_m^*\Kmat u_n
&=e_{mn},
\label{eq:E_reflection_compression}\\
\frac12\,v_m^*\Kmat v_n
&=o_{mn}.
\label{eq:O_reflection_compression}
\end{align}
Thus $\E_N$ and $\Omat_N$ are exact congruence compressions of the known Lagarias Gram kernel. The point of the present coordinates is the resulting even--odd diagonal split, the Toeplitz--Hankel form, and the matrix structures developed below.

\section{Discrete wave equations and displacement rank}
\label{sec:displacement}

The Toeplitz--Hankel structure has a useful equivalent formulation as a low-displacement-rank identity.

\begin{definition}[Universal tridiagonal matrices]
\label{def:universal_tridiagonal}
Let $\J_E^{(N)}$ be the $N\times N$ symmetric tridiagonal matrix with zero diagonal and $1/2$ on both adjacent diagonals:
\begin{equation}
\J_E^{(N)}=
\frac12
\begin{pmatrix}
0&1&&&\\
1&0&1&&\\
&1&0&\ddots&\\
&&\ddots&\ddots&1\\
&&&1&0
\end{pmatrix}.
\label{eq:J_E_definition}
\end{equation}
Let $\mathcal J_O^{(N)}$ be the same matrix except for the upper-left entry $-1/2$:
\begin{equation}
\mathcal J_O^{(N)}=
\begin{pmatrix}
-\tfrac12&\tfrac12&&&\\
\tfrac12&0&\tfrac12&&\\
&\tfrac12&0&\ddots&\\
&&\ddots&\ddots&\tfrac12\\
&&&\tfrac12&0
\end{pmatrix}.
\label{eq:J_O_definition}
\end{equation}
\end{definition}

\begin{proposition}[Discrete wave relation]
\label{prop:discrete_wave}
For $m,n\ge2$, both arrays satisfy
\begin{equation}
a_{m+1,n}+a_{m-1,n}
=a_{m,n+1}+a_{m,n-1},
\label{eq:discrete_wave}
\end{equation}
where $a=e$ or $a=o$. The boundary relations are exactly those encoded by $\J_E$ and $\mathcal J_O$.
\end{proposition}

\begin{proof}
For $e_{mn}=\lambda_{m+n}-\lambda_{|m-n|}$, the four terms in \eqref{eq:discrete_wave} agree after matching the sum and difference indices. The same calculation applies to $o_{mn}$ because the sum index is shifted by one. The boundary $m=1$ is described by the identities $xU_0=U_1/2$ and $xP_0=(P_1-P_0)/2$ developed in Section \ref{sec:moment_representations}.
\end{proof}

\begin{theorem}[Displacement rank at most two]
\label{thm:rank_two_displacement}
Let
\begin{align}
r_E^{(N)}&:=\bigl(e_{1,N+1},\dots,e_{N-1,N+1},0\bigr)^T,
\label{eq:rE_definition}\\
r_O^{(N)}&:=\bigl(o_{1,N+1},\dots,o_{N-1,N+1},0\bigr)^T,
\label{eq:rO_definition}
\end{align}
and let $e_N$ denote the last standard basis vector of $\mathbb R^N$. Then
\begin{align}
\J_E^{(N)}\E_N-\E_N\J_E^{(N)}
&=\frac12\left(r_E^{(N)}e_N^T-e_N(r_E^{(N)})^T\right),
\label{eq:E_displacement}\\
\mathcal J_O^{(N)}\Omat_N-\Omat_N\mathcal J_O^{(N)}
&=\frac12\left(r_O^{(N)}e_N^T-e_N(r_O^{(N)})^T\right).
\label{eq:O_displacement}
\end{align}
In particular, each displacement has rank at most two.
\end{theorem}

\begin{proof}
The semi-infinite arrays satisfy the commutator relation entrywise as a consequence of \eqref{eq:discrete_wave} and the boundary relation; no bounded-operator assertion is needed here. Finite truncation removes only the coupling from index $N$ to $N+1$. The omitted term is precisely the outer-product defect in \eqref{eq:E_displacement} and \eqref{eq:O_displacement}. The final component of each boundary vector has been set to zero because it cancels identically in the skew outer product.
\end{proof}

\begin{corollary}[Displacement of the inverses]
\label{cor:inverse_displacement}
Whenever $\E_N$ and $\Omat_N$ are invertible,
\begin{align}
\J_E^{(N)}\E_N^{-1}-\E_N^{-1}\J_E^{(N)}
&=-\E_N^{-1}
\left(\J_E^{(N)}\E_N-\E_N\J_E^{(N)}\right)
\E_N^{-1},
\label{eq:E_inverse_displacement}\\
\mathcal J_O^{(N)}\Omat_N^{-1}-\Omat_N^{-1}\mathcal J_O^{(N)}
&=-\Omat_N^{-1}
\left(\mathcal J_O^{(N)}\Omat_N-\Omat_N\mathcal J_O^{(N)}\right)
\Omat_N^{-1}.
\label{eq:O_inverse_displacement}
\end{align}
Hence the inverse matrices also have displacement rank at most two with respect to the same tridiagonal operators.  More precisely, if $A$ denotes either $\E_N$ or $\Omat_N$ and $A$ is invertible, then
\begin{equation}
\rank\bigl(\J^{(N)}A^{-1}-A^{-1}\J^{(N)}\bigr)
=
\rank\bigl(\J^{(N)}A-A\J^{(N)}\bigr),
\label{eq:inverse_displacement_rank_equality}
\end{equation}
with the appropriate choice of $\J_E^{(N)}$ or $\mathcal J_O^{(N)}$.  Thus inversion preserves the exact displacement rank, whether it is $0$ or $2$.
\end{corollary}

\begin{proof}
For an invertible matrix \(A\),
\[
\J A^{-1}-A^{-1}\J=-A^{-1}(\J A-A\J)A^{-1},
\]
by direct expansion.  Since left and right multiplication by invertible
matrices preserve rank, the rank equality follows as well.
\end{proof}

\section{Chebyshev moment representations}
\label{sec:moment_representations}

Let $U_n$ denote the Chebyshev polynomial of the second kind, normalized by
\begin{equation}
U_n(\cos\theta)=\frac{\sin((n+1)\theta)}{\sin\theta}.
\label{eq:U_definition}
\end{equation}
Define
\begin{equation}
P_n(x):=U_n(x)+U_{n-1}(x),
\qquad U_{-1}:=0.
\label{eq:P_definition}
\end{equation}
Then
\begin{equation}
P_{m-1}(\cos\theta)
=\frac{\sin((m-\tfrac12)\theta)}{\sin(\theta/2)}.
\label{eq:P_sine_identity}
\end{equation}

Under $\RH$, let $m_\gamma$ denote the multiplicity of the zero $1/2+i\gamma$ and set
\begin{equation}
x_\gamma:=\cos\theta_\gamma
=\frac{4\gamma^2-1}{4\gamma^2+1}.
\label{eq:x_gamma}
\end{equation}
Define finite positive Borel measures on $[-1,1]$ by
\begin{align}
\mu_E
&:=4\sum_{\gamma>0}^{\mathrm{distinct}}m_\gamma\sin^2\theta_\gamma\,\delta_{x_\gamma},
\label{eq:mu_E}\\
\mu_O
&:=4\sum_{\gamma>0}^{\mathrm{distinct}}m_\gamma\sin^2(\theta_\gamma/2)\,\delta_{x_\gamma}.
\label{eq:mu_O}
\end{align}
These are finite measures because their total masses are $\lambda_2$ and $\lambda_1$, respectively; equivalently the weights are $O(m_\gamma\gamma^{-2})$ and the standard zero-counting estimate gives convergence.

\begin{theorem}[Moment representation]
\label{thm:moment_representation}
Under $\RH$,
\begin{align}
e_{mn}
&=\int_{-1}^{1}U_{m-1}(x)U_{n-1}(x)\,\dd\mu_E(x),
\label{eq:E_moment_representation}\\
o_{mn}
&=\int_{-1}^{1}P_{m-1}(x)P_{n-1}(x)\,\dd\mu_O(x).
\label{eq:O_moment_representation}
\end{align}
Moreover,
\begin{equation}
\dd\mu_E(x)=2(1+x)\,\dd\mu_O(x).
\label{eq:christoffel_measure_relation}
\end{equation}
\end{theorem}

\begin{proof}
We first verify that the integrals are well defined.  For fixed $m,n$, the polynomials $U_{m-1}U_{n-1}$ and $P_{m-1}P_{n-1}$ are bounded on $[-1,1]$, while $\mu_E$ and $\mu_O$ are finite positive measures.  Thus both Gram integrals converge absolutely.

For the even family, using the atomic definition of $\mu_E$ and $U_{k-1}(\cos\theta)=\sin(k\theta)/\sin\theta$,
\begin{align}
\int U_{m-1}(x)U_{n-1}(x)\,\dd\mu_E(x)
&=4\sum_{\gamma>0}^{\mathrm{distinct}}m_\gamma\sin^2\theta_\gamma
U_{m-1}(x_\gamma)U_{n-1}(x_\gamma)\\
&=4\sum_{\gamma>0}^{\mathrm{distinct}}m_\gamma
\sin(m\theta_\gamma)\sin(n\theta_\gamma).
\end{align}
Counting each ordinate according to multiplicity turns the last sum into the right side of \eqref{eq:E_gram}, hence it equals $e_{mn}$.  For the odd family, \eqref{eq:P_sine_identity} gives
\begin{align}
\int P_{m-1}(x)P_{n-1}(x)\,\dd\mu_O(x)
&=4\sum_{\gamma>0}^{\mathrm{distinct}}m_\gamma\sin^2(\theta_\gamma/2)
P_{m-1}(x_\gamma)P_{n-1}(x_\gamma)\\
&=4\sum_{\gamma>0}^{\mathrm{distinct}}m_\gamma
\sin((m-\tfrac12)\theta_\gamma)
\sin((n-\tfrac12)\theta_\gamma),
\end{align}
which is \eqref{eq:O_gram}.

It remains to compare the atom masses.  At every support point $x_\gamma=\cos\theta_\gamma$,
\begin{equation}
\frac{4m_\gamma\sin^2\theta_\gamma}
{4m_\gamma\sin^2(\theta_\gamma/2)}
=\frac{\sin^2\theta_\gamma}{\sin^2(\theta_\gamma/2)}
=4\cos^2(\theta_\gamma/2)=2(1+x_\gamma).
\label{eq:weight_ratio}
\end{equation}
Thus the two atomic measures have the same support and their masses satisfy
$\mu_E(\{x_\gamma\})=2(1+x_\gamma)\mu_O(\{x_\gamma\})$ at every atom, which is exactly \eqref{eq:christoffel_measure_relation}.
\end{proof}

\begin{remark}[Christoffel transform]
\label{rem:christoffel_transform}
Equation \eqref{eq:christoffel_measure_relation} says that the even measure is, up to the scalar factor $2$, the Christoffel transform of the odd measure at the point $-1$. Here $-1$ is the left endpoint of the ambient interval $[-1,1]$ but lies outside the actual zeta support $\{x_\gamma\}\cup\{1\}$. Thus the two matrix families are linked at the level of measures, orthogonal polynomials, Jacobi matrices, determinants, and Cholesky pivots.
\end{remark}

The weights can be written directly in terms of zero ordinates:
\begin{align}
4\sin^2\theta_\gamma
&=\frac{64\gamma^2}{(4\gamma^2+1)^2},
\label{eq:E_weight_gamma}\\
4\sin^2(\theta_\gamma/2)
&=\frac{4}{4\gamma^2+1}.
\label{eq:O_weight_gamma}
\end{align}

\section{Changes of basis and hidden Hankel matrices}
\label{sec:changes_of_basis}

Write
\begin{equation}
X_N(x):=(1,x,\dots,x^{N-1})^T.
\label{eq:monomial_vector}
\end{equation}
There exist explicit lower-triangular matrices $B_E^{(N)}$ and $B_O^{(N)}$ satisfying
\begin{align}
\bigl(U_0(x),\dots,U_{N-1}(x)\bigr)^T
&=B_E^{(N)}X_N(x),
\label{eq:CE_basis_change}\\
\bigl(P_0(x),\dots,P_{N-1}(x)\bigr)^T
&=B_O^{(N)}X_N(x).
\label{eq:CO_basis_change}
\end{align}
The coefficients of $U_n$ are explicitly
\begin{equation}
U_n(x)=
\sum_{j=0}^{\lfloor n/2\rfloor}
(-1)^j\binom{n-j}{j}(2x)^{n-2j}.
\label{eq:U_explicit_coefficients}
\end{equation}
Thus both change-of-basis matrices have diagonal
\begin{equation}
1,2,2^2,\dots,2^{N-1},
\label{eq:C_diagonal}
\end{equation}
and hence
\begin{equation}
\det B_E^{(N)}=\det B_O^{(N)}=2^{N(N-1)/2}.
\label{eq:C_determinants}
\end{equation}

Define moments
\begin{equation}
\mu_k^E:=\int_{-1}^{1}x^k\,\dd\mu_E(x),
\qquad
\mu_k^O:=\int_{-1}^{1}x^k\,\dd\mu_O(x),
\label{eq:moment_sequences}
\end{equation}
and Hankel matrices
\begin{equation}
H_N^E:=\bigl(\mu_{i+j}^E\bigr)_{i,j=0}^{N-1},
\qquad
H_N^O:=\bigl(\mu_{i+j}^O\bigr)_{i,j=0}^{N-1}.
\label{eq:hidden_hankel_matrices}
\end{equation}

\begin{theorem}[Congruence to moment Hankel matrices]
\label{thm:hankel_congruence}
Under $\RH$,
\begin{align}
\E_N&=B_E^{(N)}H_N^E(B_E^{(N)})^T,
\label{eq:E_hankel_congruence}\\
\Omat_N&=B_O^{(N)}H_N^O(B_O^{(N)})^T.
\label{eq:O_hankel_congruence}
\end{align}
Consequently,
\begin{align}
\det\E_N&=2^{N(N-1)}\det H_N^E,
\label{eq:E_hankel_det_relation}\\
\det\Omat_N&=2^{N(N-1)}\det H_N^O.
\label{eq:O_hankel_det_relation}
\end{align}
\end{theorem}

\begin{proof}
Let $X_N(x)=(1,x,\ldots,x^{N-1})^T$.  For the even family put
$\Phi_E(x)=(U_0(x),\ldots,U_{N-1}(x))^T=B_E^{(N)}X_N(x)$.  Then the moment representation gives
\begin{align}
\E_N
&=\int \Phi_E(x)\Phi_E(x)^T\,\dd\mu_E(x)\\
&=B_E^{(N)}\left(\int X_N(x)X_N(x)^T\,\dd\mu_E(x)\right)(B_E^{(N)})^T\\
&=B_E^{(N)}H_N^E(B_E^{(N)})^T.
\end{align}
Every entry is a finite integral, so pulling the fixed finite matrices through the integral is purely algebraic.  The same argument with
$\Phi_O=(P_0,\ldots,P_{N-1})^T=B_O^{(N)}X_N$ gives the odd identity.

Taking determinants and using that the basis-change matrices are triangular with diagonal $1,2,\ldots,2^{N-1}$,
\begin{equation}
\det B_E^{(N)}=\det B_O^{(N)}=2^{N(N-1)/2}.
\end{equation}
Hence
\begin{equation}
\det\E_N=(\det B_E^{(N)})^2\det H_N^E
=2^{N(N-1)}\det H_N^E,
\end{equation}
and identically for $\Omat_N$.
\end{proof}

\begin{corollary}[Moment relation]
\label{cor:moment_relation}
For every $k\ge0$,
\begin{equation}
\mu_k^E=2\bigl(\mu_k^O+\mu_{k+1}^O\bigr).
\label{eq:moment_relation}
\end{equation}
Equivalently,
\begin{equation}
H_N^E=2\left(H_N^O+H_N^{O,(1)}\right),
\qquad
H_N^{O,(1)}:=\bigl(\mu_{i+j+1}^O\bigr)_{i,j=0}^{N-1}.
\label{eq:hankel_relation}
\end{equation}
\end{corollary}

\begin{proof}
The measure identity \eqref{eq:christoffel_measure_relation} gives, for every \(k\ge0\),
\[
\mu_k^E=\int x^k\,\dd\mu_E(x)
=2\int x^k(1+x)\,\dd\mu_O(x)
=2(\mu_k^O+\mu_{k+1}^O).
\]
Applying this entrywise to the Hankel matrices gives \eqref{eq:hankel_relation}.
\end{proof}

The first moments may be read directly from the original matrices:
\begin{align}
\mu_0^E&=\lambda_2,
&\mu_1^E&=\frac{\lambda_3-\lambda_1}{2},
&\mu_2^E&=\frac{\lambda_4}{4},
\label{eq:first_E_moments}\\
\mu_0^O&=\lambda_1,
&\mu_1^O&=\frac{\lambda_2-2\lambda_1}{2},
&\mu_2^O&=\frac{\lambda_3-2\lambda_2+3\lambda_1}{4}.
\label{eq:first_O_moments}
\end{align}

\section{Determinants and Vandermonde expansions}
\label{sec:determinants}

Enumerate the positive zero ordinates with multiplicity as $\gamma_1,\gamma_2,\ldots$, and write $\theta_j:=\theta_{\gamma_j}$ and $x_j:=x_{\gamma_j}$.

\begin{theorem}[Cauchy--Binet determinant formulas]
\label{thm:cauchy_binet_determinants}
Assume $\RH$. Then
\begin{equation}
\det\E_N
=4^N
\sum_{1\le j_1<\cdots<j_N}
\left(\det\bigl[\sin(m\theta_{j_k})\bigr]_{m,k=1}^N\right)^2,
\label{eq:E_cauchy_binet}
\end{equation}
and
\begin{equation}
\det\Omat_N
=4^N
\sum_{1\le j_1<\cdots<j_N}
\left(\det\bigl[\sin((m-\tfrac12)\theta_{j_k})\bigr]_{m,k=1}^N\right)^2.
\label{eq:O_cauchy_binet}
\end{equation}
\end{theorem}

\begin{proof}
We give the even proof; the odd one is identical after replacing the sine frequencies.  For $M\ge1$, let
\begin{equation}
A_{N,M}^E:=2\bigl[\sin(m\theta_j)\bigr]_{1\le m\le N,\ 1\le j\le M}.
\end{equation}
Then the truncated Gram matrix is
\begin{equation}
G_{N,M}^E=A_{N,M}^E(A_{N,M}^E)^T
=4\sum_{j=1}^{M}s_j s_j^T,
\end{equation}
where $s_j=(\sin\theta_j,\ldots,\sin N\theta_j)^T$.  For each entry $(m,n)$, the series
$4\sum_j\sin(m\theta_j)\sin(n\theta_j)$ converges absolutely by Theorem~\ref{thm:gram_representations}; hence $G_{N,M}^E\to\E_N$ entrywise.  Since the dimension $N$ is fixed, entrywise convergence is equivalent to convergence in any matrix norm, and the determinant, being a polynomial in the entries, satisfies
\begin{equation}
\det G_{N,M}^E\longrightarrow\det\E_N.
\end{equation}

For each finite $M$, the ordinary Cauchy--Binet formula applied to the $N\times M$ matrix $A_{N,M}^E$ gives, when $M\ge N$,
\begin{align}
\det G_{N,M}^E
&=\sum_{J\subset\{1,\ldots,M\},\ |J|=N}\det(A_{N,M}^E[:,J])^2\\
&=4^N\sum_{1\le j_1<\cdots<j_N\le M}
\left(\det[\sin(m\theta_{j_k})]_{m,k=1}^{N}\right)^2.
\end{align}
Every summand is nonnegative.  Therefore the finite subset sums form an increasing sequence in $M$ and have a limit in $[0,\infty]$.  Because their values equal $\det G_{N,M}^E$ and the latter converge to the finite number $\det\E_N$, monotone convergence of these finite sums gives exactly \eqref{eq:E_cauchy_binet}.  The same construction with columns
$2(\sin(\theta_j/2),\ldots,\sin((N-\tfrac12)\theta_j))^T$ gives \eqref{eq:O_cauchy_binet}.
\end{proof}

\begin{theorem}[Vandermonde form]
\label{thm:vandermonde_form}
Under $\RH$,
\begin{align}
\det\E_N
&=4^N2^{N(N-1)}
\sum_{1\le j_1<\cdots<j_N}
\left(\prod_{a=1}^N\sin^2\theta_{j_a}\right)
\left(\prod_{1\le a<b\le N}(x_{j_b}-x_{j_a})^2\right),
\label{eq:E_vandermonde}\\
\det\Omat_N
&=4^N2^{N(N-1)}
\sum_{1\le j_1<\cdots<j_N}
\left(\prod_{a=1}^N\sin^2(\theta_{j_a}/2)\right)
\left(\prod_{1\le a<b\le N}(x_{j_b}-x_{j_a})^2\right).
\label{eq:O_vandermonde}
\end{align}
\end{theorem}

\begin{proof}
For the even family,
\begin{equation}
\sin(m\theta)=\sin\theta\,U_{m-1}(\cos\theta).
\label{eq:sin_U_factor}
\end{equation}
The leading coefficient of $U_{m-1}$ is $2^{m-1}$, hence
\begin{equation}
\det[U_{m-1}(x_j)]_{m,j=1}^{N}
=2^{N(N-1)/2}\prod_{j<k}(x_k-x_j).
\label{eq:U_vandermonde}
\end{equation}
Squaring and inserting into \eqref{eq:E_cauchy_binet} yields \eqref{eq:E_vandermonde}. The proof for $\Omat_N$ uses \eqref{eq:P_sine_identity}; $P_{m-1}$ has the same leading coefficient $2^{m-1}$.
\end{proof}

In terms of the zero enumeration with multiplicity,
\begin{align}
\det\E_N
&=4^N2^{N(N-1)}
\sum_{1\le j_1<\cdots<j_N}
\prod_{a=1}^N\frac{16\gamma_{j_a}^2}{(4\gamma_{j_a}^2+1)^2}
\prod_{a<b}(x_{j_b}-x_{j_a})^2,
\label{eq:E_vandermonde_gamma}\\
\det\Omat_N
&=4^N2^{N(N-1)}
\sum_{1\le j_1<\cdots<j_N}
\prod_{a=1}^N\frac{1}{4\gamma_{j_a}^2+1}
\prod_{a<b}(x_{j_b}-x_{j_a})^2.
\label{eq:O_vandermonde_gamma}
\end{align}
If two selected indices correspond to the same multiple zero, the associated Vandermonde factor vanishes; selecting one copy produces the correct multiplicity factor.

\begin{corollary}[Arbitrary principal minors]
\label{cor:arbitrary_principal_minor_formula}
Let $I=\{n_1<\cdots<n_r\}\subset\{1,\dots,N\}$. Under $\RH$,
\begin{equation}
\det \E_N[I,I]
=4^r\sum_{1\le j_1<\cdots<j_r}
\left(\det\bigl[\sin(n_a\theta_{j_b})\bigr]_{a,b=1}^r\right)^2,
\label{eq:E_arbitrary_principal_minor}
\end{equation}
and
\begin{equation}
\det \Omat_N[I,I]
=4^r\sum_{1\le j_1<\cdots<j_r}
\left(\det\bigl[\sin((n_a-\tfrac12)\theta_{j_b})\bigr]_{a,b=1}^r\right)^2.
\label{eq:O_arbitrary_principal_minor}
\end{equation}
\end{corollary}

\begin{proof}
We prove the even identity.  Let $A_E$ be the full zero-column operator and let $A_E[I,:]$ denote the $r\times\infty$ operator obtained by retaining the rows indexed by $I=\{n_1,\ldots,n_r\}$.  Taking the corresponding principal block of the Gram identity $\E_N=A_EA_E^T$ gives
\begin{equation}
\E_N[I,I]=A_E[I,:]A_E[I,:]^T.
\end{equation}
For a finite truncation to the first $M$ zero indices, Cauchy--Binet gives
\begin{equation}
\det\E_{N,M}[I,I]
=4^r\sum_{1\le j_1<\cdots<j_r\le M}
\left(\det[\sin(n_a\theta_{j_b})]_{a,b=1}^r\right)^2.
\end{equation}
The entries of the truncated principal block converge absolutely to those of $\E_N[I,I]$ by Theorem~\ref{thm:gram_representations}, so its determinant converges.  The right-hand sums increase monotonically because all terms are nonnegative.  Passing to the limit $M\to\infty$ proves \eqref{eq:E_arbitrary_principal_minor}.  The odd formula follows identically from the half-integer zero columns.
\end{proof}

\section{Canonical Cauchy--Binet ensemble and zero influence}
\label{sec:zero_leverage}

The determinant formulas above are sums over all $N$-element zero subsets.  It is therefore important to distinguish the full determinant from any one Vandermonde term.  In this section the complete Cauchy--Binet sum itself is normalized and used to define the canonical zero-selection law attached to each reflection matrix.

For a positive ordinate $\gamma$, counted with multiplicity, define the column vectors
\begin{align}
a_{N,\gamma}^{E}
&:=2\bigl(\sin\theta_\gamma,\sin(2\theta_\gamma),\ldots,\sin(N\theta_\gamma)\bigr)^T,
\label{eq:E_zero_column}\\
a_{N,\gamma}^{O}
&:=2\bigl(\sin(\tfrac12\theta_\gamma),\sin(\tfrac32\theta_\gamma),\ldots,
\sin((N-\tfrac12)\theta_\gamma)\bigr)^T.
\label{eq:O_zero_column}
\end{align}
Let $A_E$ and $A_O$ denote the $N\times\infty$ matrices having these columns.  Under $\RH$ the sine Gram formulas are precisely
\begin{equation}
\E_N=A_EA_E^T,
\qquad
\Omat_N=A_OA_O^T.
\label{eq:zero_column_gram}
\end{equation}
The column operators are Hilbert--Schmidt because the traces of the two Gram matrices are finite.

\begin{theorem}[Canonical projection ensemble]
\label{thm:canonical_projection_ensemble}
Assume $\RH$ and fix $N$.  Let $M_N$ denote either $\E_N$ or $\Omat_N$, and let $A$ be the corresponding column operator in \eqref{eq:zero_column_gram}.  For each $N$-element subset $J$ of the positive zero indices define
\begin{equation}
\mathbb P_N(J):=
\frac{\det(A_J)^2}{\det M_N}.
\label{eq:CB_probability}
\end{equation}
Then $\mathbb P_N$ is a probability measure.  Moreover
\begin{equation}
K_N:=A^TM_N^{-1}A
\label{eq:projection_kernel_zero_indices}
\end{equation}
is an orthogonal projection of rank $N$ on the zero-index space, and
\begin{equation}
\mathbb P_N(J)=\det K_N[J,J],
\qquad |J|=N.
\label{eq:projection_DPP_formula}
\end{equation}
Thus the normalized Cauchy--Binet expansion is a fixed-size projection determinantal process on the positive zero indices.
\end{theorem}

\begin{proof}
For fixed $N$, regard $A$ as an operator from the zero-index Hilbert space $\ell^2$ to $\mathbb R^N$.  Its columns are $a_{N,\gamma}$.  Since
\begin{equation}
\sum_{\gamma>0}\|a_{N,\gamma}\|_2^2
=\operatorname{tr}\left(\sum_{\gamma>0}a_{N,\gamma}a_{N,\gamma}^T\right)
=\operatorname{tr}M_N<\infty,
\end{equation}
$A$ is Hilbert--Schmidt, hence bounded, and the series $AA^T=\sum a_\gamma a_\gamma^T$ converges in trace norm (in finite dimension, equivalently entrywise and in every matrix norm).  Under $\RH$, Theorem~\ref{thm:main_criterion} gives $M_N\succ0$, so $M_N^{-1/2}$ exists.

The infinite Cauchy--Binet identity is obtained by truncation.  Let $A^{(M)}$ contain the first $M$ columns.  Then
$A^{(M)}(A^{(M)})^T\to M_N$ in norm.  For $M\ge N$, finite Cauchy--Binet gives
\begin{equation}
\det\bigl(A^{(M)}(A^{(M)})^T\bigr)
=\sum_{J\subset\{1,\ldots,M\},\ |J|=N}\det(A_J)^2.
\end{equation}
The left side tends to $\det M_N$ by continuity of the determinant.  The right side is monotone increasing in $M$ because every summand is nonnegative; therefore
\begin{equation}
\det M_N=\sum_{|J|=N}\det(A_J)^2.
\label{eq:CB_normalization_again}
\end{equation}
Since $\det M_N>0$, \eqref{eq:CB_probability} is nonnegative and has total mass one.

Now define $B:=M_N^{-1/2}A$.  As an operator $\ell^2\to\mathbb R^N$ it satisfies
\begin{equation}
BB^*=M_N^{-1/2}AA^TM_N^{-1/2}=I_{\mathbb R^N}.
\end{equation}
Thus $B$ is a coisometry.  Its adjoint product
$K_N:=B^*B=A^TM_N^{-1}A$ is bounded and self-adjoint, and
\begin{equation}
K_N^2=B^*(BB^*)B=B^*B=K_N.
\end{equation}
Hence $K_N$ is the orthogonal projection onto $\operatorname{Ran}B^*$; because $BB^*=I_N$, this range has dimension $N$.

For a finite $N$-element index set $J$, the principal block is
$K_N[J,J]=B_J^TB_J$.  Since $B_J=M_N^{-1/2}A_J$ is square,
\begin{align}
\det K_N[J,J]
&=\det(B_J^TB_J)=\det(B_J)^2\\
&=\det(M_N^{-1/2})^2\det(A_J)^2
=\frac{\det(A_J)^2}{\det M_N},
\end{align}
which proves \eqref{eq:projection_DPP_formula}.
\end{proof}

\begin{definition}[Zero leverage]
\label{def:zero_leverage}
For either reflection family define
\begin{equation}
\ell_{N,\gamma}:=(a_{N,\gamma})^TM_N^{-1}a_{N,\gamma}.
\label{eq:zero_leverage}
\end{equation}
If a zero has multiplicity greater than one, each copy is regarded as a separate zero index with the same column vector.
\end{definition}

\begin{corollary}[Exact one- and two-zero marginals]
\label{cor:zero_marginals}
Under $\RH$,
\begin{align}
\mathbb P_N(\gamma\in J)&=\ell_{N,\gamma},
\label{eq:one_zero_marginal}\\
\mathbb P_N(\gamma,\eta\in J)
&=\ell_{N,\gamma}\ell_{N,\eta}-K_N(\gamma,\eta)^2,
\qquad \gamma\ne\eta,
\label{eq:two_zero_marginal}\\
\sum_{\gamma>0}\ell_{N,\gamma}&=N.
\label{eq:total_leverage_N}
\end{align}
In particular, the full determinant assigns an intrinsic, normalized influence profile to the zero ordinates.
\end{corollary}

\begin{proof}
We derive the marginals directly.  Fix one index $\gamma$ and remove its column from $A$.  The Cauchy--Binet normalization for the reduced column set gives
\begin{equation}
\mathbb P_N(\gamma\notin J)
=\frac{\det(M_N-a_\gamma a_\gamma^T)}{\det M_N}.
\end{equation}
Because $M_N$ is invertible,
\begin{align}
\det(M_N-a_\gamma a_\gamma^T)
&=\det M_N\,\det(I-M_N^{-1}a_\gamma a_\gamma^T)\\
&=\det M_N\,(1-a_\gamma^TM_N^{-1}a_\gamma),
\end{align}
where the last equality follows from $\det(I-uv^T)=1-v^Tu$ (choose a basis whose first vector is $u$ when $u\ne0$, or use the characteristic polynomial of the rank-one operator).  Hence
\begin{equation}
\mathbb P_N(\gamma\in J)=\ell_{N,\gamma}.
\end{equation}

For two distinct indices $\gamma,\eta$, put $A_S=(a_\gamma,a_\eta)$ and
$K_S=A_S^TM_N^{-1}A_S$.  The finite-rank determinant identity
\begin{equation}
\det(M_N-A_SA_S^T)=\det M_N\,\det(I_2-K_S)
\end{equation}
follows by factoring out $M_N$ and using $\det(I-UV)=\det(I-VU)$ for the finite matrices $U=M_N^{-1}A_S$ and $V=A_S^T$; the latter identity follows by taking the determinant of the two Schur complements of
$\begin{psmallmatrix}I&U\\V&I\end{psmallmatrix}$.  Thus
\begin{equation}
\mathbb P_N(\gamma,\eta\notin J)
=(1-\ell_\gamma)(1-\ell_\eta)-K_N(\gamma,\eta)^2.
\end{equation}
Inclusion--exclusion, together with the one-point formula, yields
\begin{align}
\mathbb P_N(\gamma,\eta\in J)
&=1-\mathbb P_N(\gamma\notin J)-\mathbb P_N(\eta\notin J)
+\mathbb P_N(\gamma,\eta\notin J)\\
&=\ell_\gamma\ell_\eta-K_N(\gamma,\eta)^2.
\end{align}
Finally, absolute convergence of the leverage sum follows from positivity and the trace computation
\begin{align}
\sum_{\gamma>0}\ell_{N,\gamma}
&=\sum_{\gamma>0}\operatorname{tr}(M_N^{-1}a_\gamma a_\gamma^T)\\
&=\operatorname{tr}\left(M_N^{-1}\sum_{\gamma>0}a_\gamma a_\gamma^T\right)
=\operatorname{tr}I_N=N,
\end{align}
where interchange of sum and trace is valid because the summands are nonnegative and the ambient matrix space is finite dimensional.
\end{proof}

The leverage has an equivalent deterministic interpretation that does not require probabilistic terminology.

\begin{theorem}[Deletion law]
\label{thm:zero_deletion_law}
Let $M_N^{(-\gamma)}$ be the Gram matrix obtained after removing one copy of the positive zero ordinate $\gamma$ from the zero-side representation.  Then
\begin{equation}
\boxed{
\frac{\det M_N^{(-\gamma)}}{\det M_N}=1-\ell_{N,\gamma}.}
\label{eq:single_zero_deletion_ratio}
\end{equation}
More generally, if $S$ is a finite set of zero indices and $K_N[S,S]$ is the corresponding principal block of \eqref{eq:projection_kernel_zero_indices}, then
\begin{equation}
\frac{\det M_N^{(-S)}}{\det M_N}
=\det\bigl(I-K_N[S,S]\bigr).
\label{eq:multi_zero_deletion_ratio}
\end{equation}
\end{theorem}

\begin{proof}
Under $\RH$, $M_N\succ0$ and is therefore invertible.  Removing one indexed copy of the zero column $a=a_{N,\gamma}$ changes the Gram sum by exactly
\begin{equation}
M_N^{(-\gamma)}=M_N-aa^T.
\end{equation}
Factoring out $M_N$ gives
\begin{equation}
\det(M_N-aa^T)=\det M_N\,\det(I-M_N^{-1}aa^T).
\end{equation}
The rank-one matrix $M_N^{-1}aa^T$ has at most one nonzero eigenvalue, equal to its trace $a^TM_N^{-1}a=\ell_{N,\gamma}$; hence
\begin{equation}
\det(I-M_N^{-1}aa^T)=1-\ell_{N,\gamma},
\end{equation}
proving \eqref{eq:single_zero_deletion_ratio}.  Notice that $M_N^{(-\gamma)}$ is only positive semidefinite in general, so the ratio is allowed to be zero; no invertibility of the deleted matrix is assumed.

For a finite indexed set $S$ of size $r$, let $A_S$ be the $N\times r$ matrix formed from the corresponding columns.  Then
\begin{equation}
M_N^{(-S)}=M_N-A_SA_S^T
\end{equation}
and therefore
\begin{align}
\frac{\det M_N^{(-S)}}{\det M_N}
&=\det(I_N-M_N^{-1}A_SA_S^T)\\
&=\det(I_r-A_S^TM_N^{-1}A_S).
\end{align}
The second equality is Sylvester's finite-dimensional identity $\det(I_N+UV)=\det(I_r+VU)$, which follows by evaluating the determinant of
$\begin{psmallmatrix}I_N&U\\-V&I_r\end{psmallmatrix}$ using either diagonal block as the Schur pivot.  By definition the $r\times r$ matrix $A_S^TM_N^{-1}A_S$ is exactly $K_N[S,S]$, proving \eqref{eq:multi_zero_deletion_ratio}.
\end{proof}

\subsection{Transport of a single zero}

The same framework gives an exact answer to how one moved zero is represented in the Li sequence and how it changes the matrix determinant.  This is a sensitivity identity for the zero-side Gram representation; the perturbed zero configuration need not itself arise from an arithmetic $L$-function.

Let one conjugate pair $1/2\pm i\gamma$ be moved to $1/2\pm i\gamma'$, leaving all other zeros fixed, and write
\begin{equation}
\theta=2\arctan\frac{1}{2\gamma},
\qquad
\theta'=2\arctan\frac{1}{2\gamma'}.
\label{eq:theta_transport}
\end{equation}

\begin{proposition}[Single-zero transport through the Li coefficients]
\label{prop:single_zero_transport_Li}
The contribution of the moved conjugate pair to the change of the Li coefficient is
\begin{equation}
\boxed{
\delta\lambda_n
=2\bigl(\cos(n\theta)-\cos(n\theta')\bigr).}
\label{eq:delta_lambda_single_zero}
\end{equation}
Consequently the reflection matrices change by
\begin{align}
\delta\E_N
&=a_{N,\gamma'}^E(a_{N,\gamma'}^E)^T
-a_{N,\gamma}^E(a_{N,\gamma}^E)^T,
\label{eq:E_single_zero_rank_two}\\
\delta\Omat_N
&=a_{N,\gamma'}^O(a_{N,\gamma'}^O)^T
-a_{N,\gamma}^O(a_{N,\gamma}^O)^T.
\label{eq:O_single_zero_rank_two}
\end{align}
Thus moving one conjugate zero pair is an update of rank at most two in either reflection family.
\end{proposition}

\begin{proof}
Under $\RH$, the pair $1/2\pm i\gamma$ contributes
\begin{equation}
2\bigl(1-\cos(n\theta)\bigr)
\end{equation}
to $\lambda_n$.  Subtraction before and after the move gives \eqref{eq:delta_lambda_single_zero}.  Using
\begin{equation}
\cos((m+n)\theta)-\cos((m-n)\theta)
=-2\sin(m\theta)\sin(n\theta)
\end{equation}
in the definitions of $\E_N$ and $\Omat_N$ yields \eqref{eq:E_single_zero_rank_two}--\eqref{eq:O_single_zero_rank_two}.
\end{proof}

\begin{theorem}[Exact determinant response to a moved zero]
\label{thm:moved_zero_determinant_ratio}
Let $M_N$ denote either reflection matrix, let $a=a_{N,\gamma}$ and $a'=a_{N,\gamma'}$ be the corresponding old and new zero columns, and define with respect to the \emph{unperturbed} inverse
\begin{equation}
\ell=a^TM_N^{-1}a,
\qquad
\ell'=(a')^TM_N^{-1}a',
\qquad
\kappa=(a')^TM_N^{-1}a.
\label{eq:moved_zero_leverages}
\end{equation}
Then
\begin{equation}
\boxed{
\frac{\det M_N'}{\det M_N}
=(1-\ell)(1+\ell')+\kappa^2.}
\label{eq:exact_moved_zero_ratio}
\end{equation}
\end{theorem}

\begin{proof}
Under $\RH$, the unperturbed matrix $M_N$ is positive definite and hence invertible.  Proposition~\ref{prop:single_zero_transport_Li} gives the exact finite-rank update
\begin{equation}
M_N'=M_N+a'(a')^T-aa^T=M_N+UCU^T,
\end{equation}
where $U=(a',a)$ is $N\times2$ and $C=\operatorname{diag}(1,-1)$.  Factoring out $M_N$,
\begin{equation}
\frac{\det M_N'}{\det M_N}
=\det(I_N+M_N^{-1}UCU^T).
\end{equation}
Apply Sylvester's identity with $X=M_N^{-1}U$ and $Y=CU^T$:
\begin{equation}
\det(I_N+XY)=\det(I_2+YX)
=\det(I_2+CU^TM_N^{-1}U).
\end{equation}
No invertibility of $M_N'$ is required for this determinant identity; if the moved configuration makes $M_N'$ singular, both sides simply vanish.

By the definitions of $\ell,\ell',\kappa$ and symmetry of $M_N^{-1}$,
\begin{equation}
U^TM_N^{-1}U
=\begin{pmatrix}\ell'&\kappa\\\kappa&\ell\end{pmatrix}.
\end{equation}
Hence
\begin{equation}
I_2+CU^TM_N^{-1}U
=\begin{pmatrix}1+\ell'&\kappa\\-\kappa&1-\ell\end{pmatrix},
\end{equation}
and direct evaluation gives
\begin{equation}
\frac{\det M_N'}{\det M_N}
=(1-\ell)(1+\ell')+\kappa^2.
\end{equation}
\end{proof}

\begin{corollary}[Infinitesimal zero sensitivity]
\label{cor:infinitesimal_zero_sensitivity}
For the even family let
\begin{equation}
s_E(\theta):=(\sin\theta,\sin2\theta,\ldots,\sin N\theta)^T,
\qquad
\dot s_E(\theta):=(\cos\theta,2\cos2\theta,\ldots,N\cos N\theta)^T.
\end{equation}
Then a motion of the conjugate pair at ordinate $\gamma$ satisfies
\begin{equation}
\boxed{
\frac{\partial}{\partial\gamma}\log\det\E_N
=-\frac{32}{4\gamma^2+1}\,
\dot s_E(\theta_\gamma)^T\E_N^{-1}s_E(\theta_\gamma).}
\label{eq:E_zero_logdet_derivative}
\end{equation}
For the odd family the identical formula holds with
\begin{align}
s_O(\theta)&:=\bigl(\sin(\tfrac12\theta),\sin(\tfrac32\theta),\ldots,
\sin((N-\tfrac12)\theta)\bigr)^T,\\
\dot s_O(\theta)&:=\bigl(\tfrac12\cos(\tfrac12\theta),\tfrac32\cos(\tfrac32\theta),\ldots,
(N-\tfrac12)\cos((N-\tfrac12)\theta)\bigr)^T.
\end{align}
\end{corollary}

\begin{proof}
Fix $N$ and one positive ordinate $\gamma_0>0$ in the unperturbed zero configuration.  The other zero columns are held fixed.  For $t>0$ define
\begin{equation}
\theta(t):=2\arctan\frac{1}{2t},
\qquad
a(t):=2s(\theta(t)),
\end{equation}
with $s=s_E$ or $s=s_O$ according to the reflection family, and define the one-parameter matrix
\begin{equation}
\widetilde M_N(t):=M_N-a(\gamma_0)a(\gamma_0)^T+a(t)a(t)^T.
\label{eq:local_moving_zero_matrix}
\end{equation}
This definition avoids any differentiation of the infinite zero sum: only one finite-rank term depends on $t$.  Since $\theta$ is $C^\infty$ on $(0,\infty)$ and $s$ is a finite vector of trigonometric functions, $t\mapsto\widetilde M_N(t)$ is $C^\infty$ as a map into the real symmetric $N\times N$ matrices.

At $t=\gamma_0$ one has $\widetilde M_N(\gamma_0)=M_N\succ0$ under $\RH$.  The smallest eigenvalue is a continuous function of the matrix entries; hence there exists $\varepsilon>0$, with $(\gamma_0-\varepsilon,\gamma_0+\varepsilon)\subset(0,\infty)$, such that $\widetilde M_N(t)\succ0$ throughout this interval.  Consequently $\det\widetilde M_N(t)>0$ there and $\log\det\widetilde M_N(t)$ is well defined and differentiable.

For any differentiable curve of invertible matrices $A(t)$, Jacobi's formula follows from
$\det(A+H)=\det A\,\det(I+A^{-1}H)$ and the first-order expansion
$\det(I+\varepsilon B)=1+\varepsilon\operatorname{tr}B+o(\varepsilon)$; hence
\begin{equation}
\frac{\dd}{\dd t}\log\det A(t)
=\operatorname{tr}\bigl(A(t)^{-1}A'(t)\bigr).
\end{equation}
Applying this at $t=\gamma_0$ to \eqref{eq:local_moving_zero_matrix},
\begin{equation}
\widetilde M_N'(\gamma_0)
=a'(\gamma_0)a(\gamma_0)^T+a(\gamma_0)a'(\gamma_0)^T,
\end{equation}
and therefore, using symmetry of $M_N^{-1}$ and cyclicity of the trace,
\begin{align}
\left.\frac{\dd}{\dd t}\log\det\widetilde M_N(t)\right|_{t=\gamma_0}
&=\operatorname{tr}\!\left(M_N^{-1}
[a'a^T+aa'^T]\right)\notag\\
&=2(a')^TM_N^{-1}a,
\end{align}
where $a=a(\gamma_0)$ and $a'=a'(\gamma_0)$.  Finally,
\begin{equation}
\theta'(t)
=2\frac{1}{1+(2t)^{-2}}\left(-\frac{1}{2t^2}\right)
=-\frac{4}{4t^2+1},
\label{eq:theta_gamma_derivative}
\end{equation}
so $a'(t)=2\theta'(t)\dot s(\theta(t))$.  Substitution gives
\begin{equation}
\left.\frac{\dd}{\dd t}\log\det\widetilde M_N(t)\right|_{t=\gamma_0}
=-\frac{32}{4\gamma_0^2+1}
\dot s(\theta_{\gamma_0})^TM_N^{-1}s(\theta_{\gamma_0}),
\end{equation}
which is \eqref{eq:E_zero_logdet_derivative} and its odd analogue.  The derivative is thus a local derivative on the open domain $t>0$ on which the perturbed matrix remains invertible; no statement is made at the boundary $t=0$.
\end{proof}

\begin{theorem}[High-zero leverage decay at fixed matrix order]
\label{thm:high_zero_leverage_decay}
Assume $\RH$ and fix $N$.  Put
\begin{align}
v_E&:=(1,2,\ldots,N)^T,
&c_N^E&:=v_E^T\E_N^{-1}v_E,
\label{eq:cNE_definition}\\
v_O&:=(\tfrac12,\tfrac32,\ldots,N-\tfrac12)^T,
&c_N^O&:=v_O^T\Omat_N^{-1}v_O.
\label{eq:cNO_definition}
\end{align}
Then, as $\gamma\to\infty$ through positive zero ordinates,
\begin{align}
\ell_{N,\gamma}^E
&=\frac{4c_N^E}{\gamma^2}+O_N(\gamma^{-4}),
\label{eq:E_leverage_high_zero}\\
\ell_{N,\gamma}^O
&=\frac{4c_N^O}{\gamma^2}+O_N(\gamma^{-4}),
\label{eq:O_leverage_high_zero}
\end{align}
and
\begin{align}
\frac{\partial}{\partial\gamma}\log\det\E_N
&=-\frac{8c_N^E}{\gamma^3}+O_N(\gamma^{-5}),
\label{eq:E_logdet_high_zero}\\
\frac{\partial}{\partial\gamma}\log\det\Omat_N
&=-\frac{8c_N^O}{\gamma^3}+O_N(\gamma^{-5}).
\label{eq:O_logdet_high_zero}
\end{align}
Consequently, if the positive ordinates are enumerated with multiplicity,
\begin{equation}
\ell_{N,\gamma_j}^{E,O}
=O_N\left(\frac{(\log j)^2}{j^2}\right),
\qquad
\sum_{j>J}\ell_{N,\gamma_j}^{E,O}
=O_N\left(\frac{(\log J)^2}{J}\right).
\label{eq:leverage_tail_bound}
\end{equation}
\end{theorem}

\begin{proof}
As $\gamma\to\infty$,
\begin{equation}
\theta_\gamma
=2\arctan\frac1{2\gamma}
=\frac1\gamma+O(\gamma^{-3}).
\end{equation}
For fixed $N$ this gives
\begin{align}
s_E(\theta_\gamma)&=\frac{v_E}{\gamma}+O_N(\gamma^{-3}),
&\dot s_E(\theta_\gamma)&=v_E+O_N(\gamma^{-2}),\\
s_O(\theta_\gamma)&=\frac{v_O}{\gamma}+O_N(\gamma^{-3}),
&\dot s_O(\theta_\gamma)&=v_O+O_N(\gamma^{-2}).
\end{align}
Substitution into \eqref{eq:zero_leverage} gives \eqref{eq:E_leverage_high_zero}--\eqref{eq:O_leverage_high_zero}; substitution into Corollary~\ref{cor:infinitesimal_zero_sensitivity} gives \eqref{eq:E_logdet_high_zero}--\eqref{eq:O_logdet_high_zero}.  Finally the Riemann--von Mangoldt formula gives $\gamma_j\asymp j/\log j$, and summation of $(\log j)^2/j^2$ gives \eqref{eq:leverage_tail_bound}.
\end{proof}

\begin{remark}[What the leverage theorem does and does not say]
\label{rem:leverage_scope}
The decay in Theorem~\ref{thm:high_zero_leverage_decay} is a statement about a \emph{fixed} matrix order $N$.  It proves that an individual sufficiently high zero has vanishing influence on that fixed determinant and that the normalized Cauchy--Binet mass has a summable high-zero tail.  It does not provide a uniform estimate when $N$ and the zero height grow together, and therefore it does not replace the large-$N$ determinant analysis of Section~\ref{sec:determinant_decay}.
\end{remark}

\section{Universal endpoint decay and the Li--Weil determinant law}
\label{sec:determinant_decay}

The determinant asymptotic separates naturally into a geometric leading term and lower-order information. The leading coefficient does not require the precise Riemann--von Mangoldt law. It is fixed by the rate at which the discrete support approaches its unique endpoint. The following theorem isolates that mechanism before specializing to the zeta zero measures.

\subsection{An endpoint-exponent theorem for discrete Gram determinants}

Let
\begin{equation}
\mu=\sum_{j\ge1}w_j\,\delta_{x_j},
\qquad
0<w_j,
\qquad
x_j=x(t_j)\uparrow1,
\label{eq:abstract_discrete_measure}
\end{equation}
be a finite positive measure, with distinct $t_1<t_2<\cdots\to\infty$. Write
\begin{equation}
D_N[\mu]
:=\det\left(\int x^{r+s}\,\dd\mu(x)\right)_{r,s=0}^{N-1}
\label{eq:abstract_hankel_determinant}
\end{equation}
for its monomial Hankel determinant.

\begin{theorem}[Endpoint-exponent law]
\label{thm:endpoint_exponent_law}
Assume that there is a number $\beta>0$ for which the following conditions hold.
\begin{enumerate}[label=(\roman*)]
\item There are constants $c,C,t_0>0$ such that, for $t\ge t_0$,
\begin{equation}
c\,t^{-\beta}\le1-x(t)\le C\,t^{-\beta},
\qquad
c\,t^{-\beta-1}\le x'(t)\le C\,t^{-\beta-1}.
\label{eq:endpoint_power_law}
\end{equation}
\item The full node sequence is not superlinearly dense on the logarithmic scale:
\begin{equation}
\log t_j\ge (1-o(1))\log j.
\label{eq:abstract_node_lower_scale}
\end{equation}
\item For every sufficiently large $N$ there are $N$ support parameters
$s_{N,1}<\cdots<s_{N,N}$ and a number $\eta_N\to0$ such that
\begin{equation}
s_{N,N}\le N^{1+\eta_N},
\qquad
s_{N,b}-s_{N,a}\ge(b-a)N^{-\eta_N}
\quad(1\le a<b\le N).
\label{eq:abstract_separated_configuration}
\end{equation}
\item The weights on these selected atoms are subleading at the determinant scale:
\begin{equation}
\sum_{a=1}^N \log w(s_{N,a})
\ge-o(N^2\log N).
\label{eq:abstract_weight_condition}
\end{equation}
\end{enumerate}
Then
\begin{equation}
\boxed{
\lim_{N\to\infty}
\frac{\log D_N[\mu]}{N^2\log N}
=-\beta.}
\label{eq:endpoint_exponent_limit}
\end{equation}
More generally, the same limit holds for the Gram determinant in any polynomial basis $p_0,p_1,\ldots$ with $\deg p_n=n$ and leading coefficients $c_n$ satisfying
\begin{equation}
\sum_{n=0}^{N-1}\log|c_n|=O(N^2).
\label{eq:basis_subleading_condition}
\end{equation}
\end{theorem}

\begin{proof}
We divide the proof into the upper bound, the lower bound, and the change-of-basis statement.

\emph{Upper bound.}
Because $\mu$ is a finite positive measure with infinite support, the monomials $1,x,\ldots,x^n$ are linearly independent in $L^2(\mu)$ for every $n$: a nonzero polynomial has only finitely many zeros, whereas the support contains infinitely many distinct points.  Hence the monic orthogonal polynomial $\pi_n$ of degree $n$ exists uniquely, and its squared norm
\begin{equation}
h_n:=\int |\pi_n(x)|^2\,\dd\mu(x)>0
\end{equation}
is the minimum of $\int|q|^2\dd\mu$ among all monic polynomials $q$ of degree $n$.  Gram--Schmidt, or equivalently triangular change of basis from monomials to the monic orthogonal polynomials, gives the exact product identity
\begin{equation}
D_N[\mu]=\prod_{n=0}^{N-1}h_n.
\label{eq:abstract_D_product_h}
\end{equation}

For $n\ge1$ take the monic competitor
\begin{equation}
q_n(x):=\prod_{j=1}^{n}(x-x_j).
\end{equation}
It vanishes at the first $n$ atoms.  Since $x_j\uparrow1$, for $j>n$ and every $1\le r\le n$,
\begin{equation}
0<x_j-x_r\le1-x_r.
\end{equation}
Therefore
\begin{equation}
|q_n(x_j)|\le\prod_{r=1}^{n}(1-x_r),\qquad j>n.
\end{equation}
Writing $M:=\mu(\mathbb R)<\infty$ and using the extremal property of $\pi_n$,
\begin{align}
h_n
&\le\int |q_n(x)|^2\,\dd\mu(x)\\
&=\sum_{j>n}w_j|q_n(x_j)|^2\\
&\le M\prod_{r=1}^{n}(1-x_r)^2.
\label{eq:abstract_monic_upper}
\end{align}

By hypothesis (i), after increasing a fixed initial index if necessary,
\begin{equation}
\log(1-x_j)\le \log C-\beta\log t_j.
\end{equation}
Hypothesis (ii) means that there exists a sequence $\varepsilon_j\to0$ such that
$\log t_j\ge(1-\varepsilon_j)\log j$ for all sufficiently large $j$.  Hence
\begin{equation}
\log(1-x_j)
\le-\beta\log j+\beta\varepsilon_j\log j+O(1).
\label{eq:abstract_point_upper_explicit}
\end{equation}
The error sums to $o(n\log n)$.  Indeed, for every $\epsilon>0$ choose $J$ so that $|\varepsilon_j|\le\epsilon$ for $j\ge J$; then
\begin{equation}
\sum_{j=J}^{n}|\varepsilon_j|\log j
\le\epsilon\sum_{j=J}^{n}\log j
=O(\epsilon n\log n),
\end{equation}
and the finitely many $j<J$ contribute $O_J(1)$.  Since
\begin{equation}
\sum_{j=1}^{n}\log j=\log(n!)=n\log n-n+O(\log n),
\end{equation}
we obtain
\begin{equation}
\sum_{j=1}^{n}\log(1-x_j)
\le-\beta n\log n+o(n\log n).
\end{equation}
Thus \eqref{eq:abstract_monic_upper} gives
\begin{equation}
\log h_n\le-2\beta n\log n+o(n\log n).
\label{eq:abstract_h_upper_rigorous}
\end{equation}
Summing \eqref{eq:abstract_h_upper_rigorous} for $0\le n<N$ is legitimate at the same scale: if $r_n=o(n\log n)$, then for every $\epsilon>0$ one has $|r_n|\le\epsilon n\log n$ for all sufficiently large $n$, so
$\sum_{n<N}r_n=o(N^2\log N)$.  Finally,
\begin{equation}
\sum_{n=1}^{N-1}n\log n
=\frac12N^2\log N-\frac14N^2+O(N\log N),
\end{equation}
and therefore
\begin{equation}
\log D_N[\mu]
\le-\beta N^2\log N+o(N^2\log N).
\label{eq:abstract_det_upper}
\end{equation}

\emph{Lower bound.}
Fix $N$ large enough that hypothesis (iii) supplies distinct support parameters
$s_{N,1}<\cdots<s_{N,N}$, and put $y_a:=x(s_{N,a})$.  The finite $N\times N$ moment matrix generated by these selected atoms is
\begin{equation}
H_N^{\mathrm{sel}}=V_NW_NV_N^T,
\end{equation}
where $(V_N)_{r,a}=y_a^{r}$ for $0\le r\le N-1$ and
$W_N=\operatorname{diag}(w(s_{N,1}),\ldots,w(s_{N,N}))$.  Since the full moment matrix is the sum of this positive semidefinite matrix and the contribution of all remaining atoms, the Cauchy--Binet expansion of the full determinant contains the selected term with positive sign; equivalently, applying Cauchy--Binet first to a finite truncation containing the selected atoms and then passing monotonically to the full measure gives
\begin{equation}
D_N[\mu]\ge
\left(\prod_{a=1}^{N}w(s_{N,a})\right)
\prod_{1\le a<b\le N}(y_b-y_a)^2.
\label{eq:abstract_one_term_lower}
\end{equation}

For $a<b$, monotonicity of $x$ and hypothesis (i) give
\begin{align}
y_b-y_a
&=\int_{s_{N,a}}^{s_{N,b}}x'(u)\,\dd u\\
&\ge c\,s_{N,N}^{-\beta-1}(s_{N,b}-s_{N,a}),
\end{align}
because $u\le s_{N,N}$ on the interval and $u\mapsto u^{-\beta-1}$ is decreasing.  By hypothesis (iii),
\begin{align}
y_b-y_a
&\ge c(b-a)
N^{-(\beta+1)(1+\eta_N)}N^{-\eta_N}\\
&=c(b-a)N^{-\beta-1-(\beta+2)\eta_N}.
\label{eq:abstract_node_gap_lower_rigorous}
\end{align}
Taking logarithms and summing over the $\binom N2$ pairs,
\begin{align}
2\sum_{a<b}\log(y_b-y_a)
&\ge 2\sum_{a<b}\log(b-a)
-2[\beta+1+(\beta+2)\eta_N]\binom N2\log N
+O(N^2).
\label{eq:abstract_vdm_lower_expand}
\end{align}
The elementary identity
\begin{equation}
\prod_{1\le a<b\le N}(b-a)=\prod_{k=1}^{N-1}k!
\end{equation}
and Stirling's formula imply
\begin{equation}
\sum_{a<b}\log(b-a)
=\frac12N^2\log N-\frac34N^2+O(N\log N).
\label{eq:abstract_pair_log_sum}
\end{equation}
Since $\eta_N\to0$, \eqref{eq:abstract_vdm_lower_expand} and \eqref{eq:abstract_pair_log_sum} yield
\begin{equation}
2\sum_{a<b}\log(y_b-y_a)
\ge-\beta N^2\log N-o(N^2\log N).
\end{equation}
Hypothesis (iv) states precisely that the selected weights do not decrease this bound at the leading scale.  Inserting the result into \eqref{eq:abstract_one_term_lower} gives
\begin{equation}
\log D_N[\mu]
\ge-\beta N^2\log N-o(N^2\log N).
\label{eq:abstract_det_lower}
\end{equation}
Together with \eqref{eq:abstract_det_upper}, this proves the limit.

\emph{Polynomial bases.}
Let $p_n(x)=c_nx^n+$ terms of degree $<n$, and let $C_N$ be the lower triangular coefficient matrix taking the monomial vector $(1,x,\ldots,x^{N-1})^T$ to $(p_0,\ldots,p_{N-1})^T$.  Its diagonal is $(c_0,\ldots,c_{N-1})$, so
$\det C_N=\prod_{n<N}c_n$.  If $G_N[p]$ denotes the corresponding Gram matrix, then exactly
\begin{equation}
G_N[p]=C_NH_NC_N^T,
\end{equation}
and hence
\begin{equation}
\log\det G_N[p]
=\log D_N[\mu]+2\sum_{n=0}^{N-1}\log|c_n|.
\end{equation}
Hypothesis \eqref{eq:basis_subleading_condition} makes the second term $O(N^2)=o(N^2\log N)$, so the leading coefficient remains $-\beta$.
\end{proof}

\begin{remark}[What is universal]
\label{rem:what_is_universal}
Theorem~\ref{thm:endpoint_exponent_law} shows that the leading constant is the endpoint exponent $-\beta$. The precise counting law is absent from the coefficient. In particular, $t_j\asymp j$, $t_j\asymp j/\log j$, and more generally logarithmic refinements compatible with \eqref{eq:abstract_node_lower_scale} and \eqref{eq:abstract_separated_configuration} give the same $-\beta$. The counting law is used only to verify that enough nodes exist in the required logarithmic scale. Arithmetic information, if present, can therefore enter only below the $N^2\log N$ term.
\end{remark}

\begin{remark}[Relation with earlier discrete Hankel asymptotics]
\label{rem:monien_haynes_position}
The mere occurrence of an $N^2\log N$ scale for a discrete Hankel determinant is not new. Monien \cite{Monien2009} derived a discrete Vandermonde representation for Hankel determinants of Dirichlet series.  For the reciprocal-integer zeta model his Coulomb-gas calculation leads to
\begin{equation}
\log H_n^{(r)}[\zeta]
=-n^2\left(\log(2n)-\frac32\right)+o(n^2)
\label{eq:monien_model_asymptotic}
\end{equation}
at the level of the variational analysis presented there.  Monien explicitly notes that the passage from the discrete optimizer to the continuous density is assumed and suggests Poissonization as a route to a fully rigorous treatment.  Thus \eqref{eq:monien_model_asymptotic} should not be cited here as a proved general asymptotic theorem from that preprint.  Its leading coefficient nevertheless provides an important external consistency check: the nodes $x_m=1/m$ have endpoint exponent $\beta=1$, and Theorem~\ref{thm:endpoint_exponent_law} gives the same rigorous leading term $-n^2\log n$ under its hypotheses.  The present theorem is broader at leading order---arbitrary $\beta$ and non-explicit node sets satisfying coarse separation hypotheses---whereas Monien's special-model variational calculation is finer at order $n^2$.

Haynes and Zudilin \cite{HaynesZudilin2015} consider $\det(\zeta(a(i+j)+b))$.  They distinguish two different statements: the expected full law has coefficient $-a$, while their elementary theorem proves only the weaker upper estimate
\begin{equation}
\log\left|\det(\zeta(a(i+j)+b))_{1\le i,j\le n}\right|
\le-\frac a2 n^2\log n+O(n^2).
\end{equation}
There is therefore no factor-two contradiction with the endpoint law: $m^{-a}$ has $\beta=a$, so Theorem~\ref{thm:endpoint_exponent_law} predicts the conjectural/full leading coefficient $-a$ when its lower-bound hypotheses are available; Haynes--Zudilin prove a one-sided estimate with coefficient $-a/2$ by a different elementary argument.

Johansson's discrete orthogonal-polynomial framework \cite{Johansson2001} is also closely related but does not make Theorem~\ref{thm:endpoint_exponent_law} an immediate corollary, and the reason is structural rather than a matter of coverage.  Every ensemble treated there is supported on a \emph{uniform} lattice: the configurations range over $\mathbb N^N$ or $\{0,\dots,K\}^N$, the Vandermonde is formed in the lattice variable, and the asymptotic analysis proceeds through a constrained equilibrium measure whose saturation constraint $\rho\le1$ is precisely an expression of that uniform spacing.  The detailed results are then obtained for the structured Meixner, Charlier, Krawtchouk, Hahn, and Plancherel weights.

The support relevant here is of a different type.  The nodes $x_\gamma$ of \eqref{eq:universal_li_cayley_map} are nonuniform, do not lie on a lattice, and accumulate at an endpoint.  Thus there is no fixed uniform saturation constraint of the form $\rho\le1$ inherited from unit lattice spacing; any continuum occupancy constraint would instead have to reflect the local nonuniform node density.  The hypotheses used above, \eqref{eq:endpoint_power_law}, \eqref{eq:abstract_node_lower_scale} and \eqref{eq:abstract_separated_configuration}, are coarse and impose no lattice structure at all.  Accordingly, Johansson supplies the rigorous discrete-Coulomb-gas context and is directly relevant to the lower-order constrained-equilibrium questions raised in Remark~\ref{rem:second_order_not_artifact}, but Theorem~\ref{thm:endpoint_exponent_law} is logically independent of the results stated there.  This is a statement about \cite{Johansson2001} only; no general priority is claimed against the wider literature on discrete orthogonal polynomials with an accumulation point, on Stieltjes-type moment problems, or on Hankel determinant asymptotics.

These are different matrices from the Li--Weil reflection pair, with Dirichlet-series or lattice nodes rather than Cayley-transformed zero ordinates. Accordingly, no novelty is claimed here for the generic presence of $N^2\log N$ or for constrained equilibrium methods themselves. The distinction of Theorem~\ref{thm:endpoint_exponent_law} is the elementary two-sided endpoint-exponent principle under coarse geometric hypotheses and its application to the Li--Weil reflection measures below.
\end{remark}

\subsection{\texorpdfstring{The zeta reflection measures and the universal constant $-2$}{The zeta reflection measures and the universal constant -2}}

Let $\gamma_1^*<\gamma_2^*<\cdots$ denote the distinct positive ordinates and set
\begin{equation}
\delta_j:=1-x_{\gamma_j^*}=\frac{2}{4(\gamma_j^*)^2+1}.
\label{eq:delta_distinct_nodes}
\end{equation}
Multiplicity is deliberately omitted from this enumeration.

\begin{lemma}[Coarse location of distinct ordinates]
\label{lem:distinct_ordinate_lower_bound}
There are constants $c>0$ and $j_0$ such that
\begin{equation}
\gamma_j^*\ge c\frac{j}{\log j},
\qquad j\ge j_0.
\label{eq:distinct_gamma_lower_bound}
\end{equation}
This estimate requires no simplicity hypothesis.
\end{lemma}

\begin{proof}
Let $N^*(T)$ count distinct positive ordinates and let $N(T)$ count positive ordinates with multiplicity. Then $N^*(T)\le N(T)$. The Riemann--von Mangoldt formula gives $N(T)\ll T\log T$, hence
$j\le N(\gamma_j^*)\ll \gamma_j^*\log\gamma_j^*$, which implies \eqref{eq:distinct_gamma_lower_bound}. Multiple zeros can only decrease $N^*(T)$ and therefore move the distinct ordinates to larger heights; the estimate is thus robust with respect to multiplicity.
\end{proof}

\begin{lemma}[A deterministically separated zero subsequence]
\label{lem:separated_zero_subsequence}
There are constants $H>0$, $T_0>0$, and distinct positive ordinates $\widetilde\gamma_j$ such that, for all $j\ge1$,
\begin{equation}
T_0+2Hj\le \widetilde\gamma_j\le T_0+2Hj+H.
\label{eq:separated_gamma_boxes}
\end{equation}
Consequently $\widetilde\gamma_j\asymp j$ and
\begin{equation}
\widetilde\gamma_b-\widetilde\gamma_a\ge H(b-a),
\qquad 1\le a<b.
\label{eq:separated_gamma_gap}
\end{equation}
\end{lemma}

\begin{proof}
Write the Riemann--von Mangoldt formula as
\begin{equation}
N(T)=F(T)+R(T),
\qquad
F(T):=\frac{T}{2\pi}\log\frac{T}{2\pi}-\frac{T}{2\pi},
\qquad
R(T)=O(\log T).
\label{eq:RVM_for_boxes}
\end{equation}
Thus there are absolute constants $A,T_1>0$ such that $|R(T)|\le A\log T$ for $T\ge T_1$. For every fixed $H\ge1$,
\begin{equation}
F(T+H)-F(T)=\frac{H}{2\pi}\log T+O_H(1),
\qquad T\to\infty.
\label{eq:RVM_short_fixed_interval}
\end{equation}
Choose once and for all $H>4\pi A$. Then, for sufficiently large $T$,
\[
N(T+H)-N(T)
\ge
\left(\frac{H}{2\pi}-2A+o(1)\right)\log T>0.
\]
Hence every interval $[T,T+H]$ with $T\ge T_0$ contains a zero for some $T_0\ge T_1$. Applying this to the alternating boxes $[T_0+2Hj,T_0+2Hj+H]$ gives \eqref{eq:separated_gamma_boxes}; the boxes are disjoint and \eqref{eq:separated_gamma_gap} follows immediately.
\end{proof}

\begin{theorem}[Li--Weil determinant law]
\label{thm:determinant_first_order}
Assume $\RH$. Then, for both reflection families,
\begin{equation}
\boxed{
\lim_{N\to\infty}
\frac{\log\Delta_N^E}{N^2\log N}
=
\lim_{N\to\infty}
\frac{\log\Delta_N^O}{N^2\log N}
=-2.}
\label{eq:determinant_first_order_limit}
\end{equation}
More quantitatively,
\begin{align}
-2N^2\log N+O(N^2)
&\le \log\Delta_N^E
\le -2N^2\log N+2N^2\log\log N+O(N^2),
\label{eq:E_determinant_two_sided}\\
-2N^2\log N+O(N^2)
&\le \log\Delta_N^O
\le -2N^2\log N+2N^2\log\log N+O(N^2).
\label{eq:O_determinant_two_sided}
\end{align}
\end{theorem}

\begin{proof}
For both zero measures the support map is
\begin{equation}
x(t)=1-\frac{2}{4t^2+1},
\label{eq:x_of_t_decay}
\end{equation}
so
\begin{equation}
1-x(t)\asymp t^{-2},
\qquad
x'(t)\asymp t^{-3}.
\label{eq:zeta_endpoint_beta_two}
\end{equation}
Thus $\beta=2$. Lemma~\ref{lem:distinct_ordinate_lower_bound} verifies the full-sequence condition \eqref{eq:abstract_node_lower_scale}, while Lemma~\ref{lem:separated_zero_subsequence} verifies \eqref{eq:abstract_separated_configuration} with $\eta_N=0$ up to harmless constants. Along that separated subsequence, the multiplicity of every selected ordinate is at least one and \(\widetilde\gamma_j\asymp j\).  Hence each selected one-point weight in \eqref{eq:E_vandermonde_gamma}--\eqref{eq:O_vandermonde_gamma} is bounded below by a constant multiple of \(j^{-2}\).  Consequently its logarithmic product is bounded below by \(-O(N\log N)\), which is exactly the one-sided information required in \eqref{eq:abstract_weight_condition}.  No upper bound on zero multiplicities is used. The Chebyshev-adapted bases have leading coefficient $2^n$, hence their total basis contribution is $O(N^2)$. Theorem~\ref{thm:endpoint_exponent_law} therefore gives the limit \eqref{eq:determinant_first_order_limit}.

For the quantitative upper bound, let $h_n$ be the squared norm of the monic orthogonal polynomial for either $\mu_E$ or $\mu_O$. The competitor
\begin{equation}
q_n(x):=\prod_{j=1}^{n}(x-x_{\gamma_j^*})
\label{eq:extremal_competitor}
\end{equation}
gives
\begin{equation}
h_n\le M\prod_{j=1}^{n}\delta_j^2,
\qquad M=\mu([-1,1]).
\label{eq:monic_norm_upper_bound}
\end{equation}
Lemma~\ref{lem:distinct_ordinate_lower_bound} yields
\begin{equation}
\delta_j\ll\frac{(\log j)^2}{j^2},
\label{eq:delta_upper_RVM}
\end{equation}
and therefore, after restoring the leading coefficient $2^n$ of $U_n$ or $P_n$,
\begin{equation}
\log d_{n+1}^{E,O}
\le -4n\log n+4n\log\log n+O(n).
\label{eq:pivot_upper_asymptotic}
\end{equation}
Summing gives the upper inequalities in \eqref{eq:E_determinant_two_sided}--\eqref{eq:O_determinant_two_sided}.

For the quantitative lower bound, retain the separated term $\widetilde\gamma_1,\ldots,\widetilde\gamma_N$ in the nonnegative Cauchy--Binet expansion. Since $x'(t)\asymp t^{-3}$ and \eqref{eq:separated_gamma_gap} holds,
\begin{equation}
\widetilde x_b-\widetilde x_a
\gg\frac{b-a}{b^3}.
\label{eq:regularized_x_gap}
\end{equation}
Hence
\begin{align}
2\sum_{a<b}\log(\widetilde x_b-\widetilde x_a)
&\ge
2\sum_{a<b}\log(b-a)-6\sum_{a<b}\log b+O(N^2)\\
&=-2N^2\log N+O(N^2),
\label{eq:regularized_vandermonde_lower}
\end{align}
using
\begin{equation}
\sum_{a<b}\log(b-a)=\frac12N^2\log N+O(N^2),
\qquad
\sum_{a<b}\log b=\frac12N^2\log N+O(N^2).
\label{eq:pair_log_sums}
\end{equation}
The selected one-point weights can decrease the logarithm of this retained term by at most \(O(N\log N)\), while the explicit basis prefactor contributes only \(O(N^2)\); this gives the lower inequalities.
\end{proof}

\begin{corollary}[Universal Cayley exponent]
\label{cor:universal_cayley_exponent}
For every positive discrete Gram system satisfying the hypotheses of Theorem~\ref{thm:endpoint_exponent_law} with the Li--Cayley support map
\begin{equation}
x(t)=\frac{4t^2-1}{4t^2+1},
\label{eq:universal_li_cayley_map}
\end{equation}
the leading determinant coefficient is $-2$. In particular, the value $-2$ is unchanged by replacing a $T\log T$ counting law by a $T$ counting law, or by any other logarithmic refinement compatible with the theorem's node hypotheses.
\end{corollary}

\begin{proof}
The map \eqref{eq:universal_li_cayley_map} has endpoint exponent $\beta=2$, so the conclusion is Theorem~\ref{thm:endpoint_exponent_law}.
\end{proof}

\begin{remark}[Automorphic Li coefficients]
\label{rem:automorphic_universality}
Lagarias defines generalized Li coefficients for principal automorphic $L$-functions and relates them to Weil's quadratic functional \cite{Lagarias2007}. Consequently, whenever a self-dual member of that class is placed under GRH and its canonical reflection compression has the same positive Cayley zero measure with the standard coarse zero-counting and separation properties, Corollary~\ref{cor:universal_cayley_exponent} gives the same coefficient $-2$. The degree and conductor can change lower-order terms but not the endpoint exponent. This is a conditional universality statement: the present paper does not identify a real even--odd reflection pair for every non-self-dual automorphic representation.
\end{remark}

\begin{corollary}[Condition-number obstruction]
\label{cor:condition_number_obstruction}
Assume $\RH$. Then
\begin{equation}
\boxed{
\liminf_{N\to\infty}
\frac{\log\kappa_2(\E_N)}{N\log N}\ge2,
\qquad
\liminf_{N\to\infty}
\frac{\log\kappa_2(\Omat_N)}{N\log N}\ge2.}
\label{eq:condition_number_liminf}
\end{equation}
More quantitatively,
\begin{equation}
\log\kappa_2(\E_N),\ \log\kappa_2(\Omat_N)
\ge
2N\log N-2N\log\log N-O(N).
\label{eq:condition_number_quantitative}
\end{equation}
\end{corollary}

\begin{proof}
For a positive definite $N\times N$ matrix $A_N$,
\begin{equation}
\lambda_{\min}(A_N)^N\le\det A_N,
\end{equation}
so
\begin{equation}
\log\kappa_2(A_N)
\ge \log\lambda_{\max}(A_N)-\frac1N\log\det A_N.
\label{eq:condition_from_det}
\end{equation}
For the even family $\lambda_{\max}(\E_N)\ge(\E_N)_{11}=\lambda_2>0$, while for the odd family $\lambda_{\max}(\Omat_N)\ge(\Omat_N)_{11}=\lambda_1>0$. Insert the determinant upper bounds from Theorem~\ref{thm:determinant_first_order}. This gives \eqref{eq:condition_number_quantitative}, and division by $N\log N$ gives \eqref{eq:condition_number_liminf}.
\end{proof}

\begin{corollary}[Divergence of the inverse norm under RH]
\label{cor:inverse_norm_divergence}
Assume \(\RH\).  Then
\begin{equation}
\boxed{
\log\|\E_N^{-1}\|_2
\ge
2N\log N-2N\log\log N-O(N),}
\label{eq:E_inverse_norm_lower}
\end{equation}
and consequently \(\|\E_N^{-1}\|_2\to\infty\) faster than \(C^N\) for every fixed
\(C>1\).  The same statement holds for \(\Omat_N^{-1}\).
\end{corollary}

\begin{proof}
For a positive definite matrix,
\(\|A_N^{-1}\|_2=\lambda_{\min}(A_N)^{-1}\) and
\(\lambda_{\min}(A_N)^N\le\det A_N\).  The determinant upper bound in
Theorem~\ref{thm:determinant_first_order} therefore gives
\[
\log\lambda_{\min}(\E_N)
\le
-2N\log N+2N\log\log N+O(N),
\]
which is equivalent to \eqref{eq:E_inverse_norm_lower}.  The odd case is identical.
\end{proof}

\begin{remark}[Congruence warning for first-difference eigenvalues]
\label{rem:first_difference_inverse_congruence}
Let
\(\widehat E_N:=\tfrac12D_N\E_ND_N^T\).  Then
\begin{equation}
\E_N^{-1}=\frac12D_N^T\widehat E_N^{-1}D_N.
\label{eq:E_inverse_from_Ehat}
\end{equation}
Hence a tiny value of \(\lambda_{\min}(\widehat E_N)\) does force a large inverse after
accounting for the congruence, but one must not write
\(\|\E_N^{-1}\|_2=1/\lambda_{\min}(\widehat E_N)\).  More precisely,
\begin{equation}
\|\E_N^{-1}\|_2
\ge
\frac{\sigma_{\min}(D_N)^2}{2\,\lambda_{\min}(\widehat E_N)}.
\label{eq:E_inverse_Ehat_lower}
\end{equation}
The direct identity is
\(\|\E_N^{-1}\|_2=1/\lambda_{\min}(\E_N)\).
\end{remark}

\begin{remark}[The condition-number bound is far from sharp]
\label{rem:condition_number_not_sharp}
Corollary~\ref{cor:condition_number_obstruction} passes through $\lambda_{\min}(A_N)^N\le\det A_N$, which replaces the smallest eigenvalue by the geometric mean of the spectrum. For these matrices the spectrum is extremely spread, so the substitution is lossy. In a direct computation with the measure carried by the first $1200$ ordinates and eigenvalues evaluated at $1200$ decimal digits, the observed ratios
\begin{equation}
\frac{\log\lambda_{\min}(\E_N)}{N\log N}
\label{eq:observed_lambda_min_ratio}
\end{equation}
are $-11.9,\,-8.85,\,-7.83,\,-7.31,\,-6.94,\,-6.68,\,-6.48,\,-6.32$ for $N=2,4,\dots,16$, decreasing slowly. The scale that \eqref{eq:condition_from_det} induces for the same ratio \eqref{eq:observed_lambda_min_ratio} is only $-2$, since the argument replaces $\log\lambda_{\min}$ by $N^{-1}\log\det$; the sign convention here is that of $\lambda_{\min}$, not of $\kappa_2$. Correspondingly, at $N=16$ the computation gives $\log_{10}\kappa_2(\E_{16})\approx123.9$ against a proved lower bound near $38.5$. The corollary is therefore correct but conservative by a factor of about three at these orders. All numbers refer to the measure truncated to the first $1200$ ordinates and are reported as computation on that truncation, not as values for the exact infinite matrices, and not as an asymptotic claim.
\end{remark}

\begin{remark}[Why the $N^2\log\log N$ term is unresolved]
\label{rem:second_order_not_artifact}
The upper term $2N^2\log\log N$ in Theorem~\ref{thm:determinant_first_order} is not declared to be an artifact. Positivity implies that one favorable Cauchy--Binet subset gives a lower bound, but it does not imply that the logarithm of the full infinite sum equals the logarithm of its largest term; an entropy contribution can survive below the leading scale. Moreover, the Riemann--von Mangoldt clock itself suggests that the displayed logarithmic correction is structurally plausible. If, heuristically, $t_j\asymp j/\log j$ with local spacing $t_{j+1}-t_j\asymp1/\log j$, then for a general endpoint exponent $\beta$ one has at the pairwise level
\begin{equation}
x(t_b)-x(t_a)
\asymp
\frac{(b-a)(\log b)^{\beta}}{b^{\beta+1}},
\label{eq:clock_gap_second_order}
\end{equation}
and the corresponding Vandermonde energy has the scale
\begin{equation}
-\beta N^2\log N
+\beta N^2\log\log N
+O(N^2).
\label{eq:clock_vandermonde_second_order}
\end{equation}
Thus for the Li--Cayley value $\beta=2$ the coefficient $+2$ in the present upper bound is exactly what the smooth clock predicts. Proving a matching second-order law would require substantially finer information on the dense zero configuration and is deliberately separated from the universal leading theorem. Theorem~\ref{thm:conditional_second_order} below isolates exactly what that finer information has to be.
\end{remark}

The next statement isolates a second-order law under explicit strengthened hypotheses on the node scale, the aggregate gap geometry, and the first-$N$ atom weights.  No implication of these hypotheses from $\RH$ alone is asserted. Write
\begin{equation}
\Lambda_N:=
\sum_{1\le a<b\le N}
\left[
\log(t_b-t_a)-\log\frac{b-a}{\log t_b}
\right].
\label{eq:gap_regularity_functional}
\end{equation}

\begin{theorem}[Conditional second-order law]
\label{thm:conditional_second_order}
Assume the hypotheses of Theorem~\ref{thm:endpoint_exponent_law} together with the strengthened node scale
\begin{equation}
\log t_j=\log j-\log\log j+O(1),
\label{eq:RVM_node_scale}
\end{equation}
the aggregate gap regularity hypothesis
\begin{equation}
\Lambda_N=O(N^2),
\label{eq:gap_regularity_hypothesis}
\end{equation}
and the first-$N$ weight condition
\begin{equation}
\sum_{j=1}^{N}\log w_j\ge -O(N^2).
\label{eq:second_order_weight_condition}
\end{equation}
Then
\begin{equation}
\log D_N[\mu]
=-\beta N^2\log N+\beta N^2\log\log N+O(N^2).
\label{eq:conditional_second_order_law}
\end{equation}

For the Li--Weil reflection measures, assume $\RH$ and, in addition, that the distinct positive ordinates satisfy
\begin{equation}
\log\gamma_j^*
=\log j-\log\log j+O(1)
\label{eq:distinct_RVM_scale_hypothesis}
\end{equation}
and that \eqref{eq:gap_regularity_hypothesis} holds with $t_j=\gamma_j^*$.  Then
\begin{equation}
\log\Delta_N^{E},\ \log\Delta_N^{O}
=-2N^2\log N+2N^2\log\log N+O(N^2).
\label{eq:conditional_second_order_li_weil}
\end{equation}
No claim is made that \eqref{eq:distinct_RVM_scale_hypothesis} follows from $\RH$ alone.
\end{theorem}

\begin{proof}
The upper bound is already available. Indeed $1-x(t_j)\asymp t_j^{-\beta}$ and \eqref{eq:RVM_node_scale} give $\log\bigl(1-x(t_j)\bigr)=-\beta\log j+\beta\log\log j+O(1)$, so the extremal-polynomial estimate used in Theorem~\ref{thm:endpoint_exponent_law} yields
\begin{equation}
\log h_n\le-2\beta n\log n+2\beta n\log\log n+O(n),
\end{equation}
and summation over $n=0,\dots,N-1$, together with $\sum_{n<N}n\log n=\tfrac12N^2\log N+O(N^2)$ and $\sum_{n<N}n\log\log n=\tfrac12N^2\log\log N+O(N^2)$, gives the upper bound in \eqref{eq:conditional_second_order_law}.

For the lower bound retain in the nonnegative Cauchy--Binet expansion the single term carried by the first $N$ atoms.  By the lower derivative bound in \eqref{eq:endpoint_power_law},
\begin{align}
x_b-x_a
&=\int_{t_a}^{t_b}x'(u)\,\dd u\\
&\ge c\int_{t_a}^{t_b}u^{-\beta-1}\,\dd u\\
&\ge c\,(t_b-t_a)t_b^{-\beta-1},
\qquad 1\le a<b\le N,
\label{eq:second_order_node_gap_lower}
\end{align}
because $u\le t_b$ on $[t_a,t_b]$.  Hence
\begin{equation}
2\sum_{a<b}\log(x_b-x_a)
\ge
2\sum_{a<b}\log(t_b-t_a)
-2(\beta+1)\sum_{a<b}\log t_b
+O(N^2).
\label{eq:second_order_vandermonde_lower}
\end{equation}
By \eqref{eq:gap_regularity_hypothesis} and \eqref{eq:RVM_node_scale},
\begin{equation}
\sum_{a<b}\log(t_b-t_a)
=\sum_{a<b}\log(b-a)-\sum_{a<b}\log\log t_b+O(N^2)
=\tfrac12N^2\log N-\tfrac12N^2\log\log N+O(N^2),
\end{equation}
while \eqref{eq:RVM_node_scale} gives
\begin{equation}
\sum_{a<b}\log t_b=\tfrac12N^2\log N-\tfrac12N^2\log\log N+O(N^2).
\end{equation}
Combining the preceding estimates gives
\begin{equation}
2\sum_{a<b}\log(x_b-x_a)
\ge
-\beta N^2\log N+\beta N^2\log\log N-O(N^2).
\label{eq:second_order_vandermonde_final_lower}
\end{equation}
Hypothesis \eqref{eq:second_order_weight_condition} shows that the first-$N$ one-point weights can lower the logarithm of the retained Cauchy--Binet term by at most $O(N^2)$. Therefore
\begin{equation}
\log D_N[\mu]
\ge
-\beta N^2\log N+\beta N^2\log\log N-O(N^2),
\end{equation}
which matches the upper bound.

For the Li--Weil measures, \eqref{eq:distinct_RVM_scale_hypothesis} and the explicit weights in \eqref{eq:E_weight_gamma}--\eqref{eq:O_weight_gamma} give
\begin{equation}
\sum_{j\le N}\log w_{E}(\gamma_j^*),
\quad
\sum_{j\le N}\log w_{O}(\gamma_j^*)
\ge -O(N\log N),
\end{equation}
because every zero multiplicity is at least one.  This one-sided lower bound is all that
\eqref{eq:second_order_weight_condition} requires; no upper bound on multiplicities is assumed.  The Chebyshev changes of basis contribute only $O(N^2)$ to the logarithm of the two reflection determinants.  Thus the same second-order formula holds for $\Delta_N^E$ and $\Delta_N^O$ under the additional hypotheses stated above.
\end{proof}

\begin{remark}[Where the gap hypothesis actually bites]
\label{rem:gap_hypothesis_localised}
Hypothesis \eqref{eq:gap_regularity_hypothesis} is an upper and a lower constraint, but only the lower one has content, and only for close pairs. On one side, a counting law of the form $N(T+h)-N(T)\ll h\log T+\log T$ forces $t_b-t_a\gg(b-a)/\log t_b$ whenever $b-a\gg\log t_N$, with no further hypothesis; the error term $\log T$ is what prevents the same deduction for smaller $b-a$. Hence \eqref{eq:gap_regularity_hypothesis} constrains only the $O(N\log N)$ pairs with $b-a\ll\log t_N$, and it permits those separations to be as small as $\exp(-cN/\log N)$ on logarithmic average. It is logically independent of simplicity, because the sequence \(t_j\) contains distinct ordinates and carries no multiplicity information, and it is far weaker than any conjectured pointwise lower bound on normalized gaps. What it is not is a theorem: no unconditional lower bound of this kind for consecutive ordinates is currently available, which is precisely why \eqref{eq:conditional_second_order_law} is stated conditionally.
\end{remark}

\begin{remark}[Size of the Cauchy--Binet defect]
\label{rem:cauchy_binet_defect_numerics}
The gap described in Remark~\ref{rem:second_order_not_artifact} can be measured. With the same finite measure on both sides, namely the first $1200$ ordinates, and $\log\det$ evaluated by Cholesky at $3000$ decimal digits, the difference between $\log\Delta_N^E$ and the logarithm of the single Cauchy--Binet term carried by the first $N$ atoms is as follows.
\begin{center}
\begin{tabular}{rrrrr}
$N$ & $\log\Delta_N^E$ & single term & defect & defect$/|\log\Delta_N^E|$\\
\hline
$2$ & $-17.336$ & $-20.421$ & $3.085$ & $0.178$\\
$4$ & $-88.812$ & $-96.619$ & $7.807$ & $0.088$\\
$8$ & $-414.694$ & $-438.761$ & $24.067$ & $0.058$\\
$12$ & $-1002.853$ & $-1048.617$ & $45.765$ & $0.046$\\
$16$ & $-1869.011$ & $-1943.919$ & $74.908$ & $0.040$
\end{tabular}
\end{center}
Over the computed range the normalized defect decreases, which is consistent with a contribution of lower order relative to $N^2\log N$; a power-law fit over these orders is close to $N^{3/2}$. No asymptotic exponent is claimed, and no growth rate is asserted beyond the tabulated orders: both statements would require an independent proof, and the positivity of the Cauchy--Binet expansion gives only $\text{defect}\ge0$.

What the table does establish, without extrapolation, is that a single Cauchy--Binet term is already quantitatively inadequate at these orders. Consequently no statement below the $N^2\log N$ scale should be extracted from one term without control of $\log\sum_J$. This is the quantitative form of the caveat in Remark~\ref{rem:second_order_not_artifact}.
\end{remark}

\begin{remark}[Zero logarithmic capacity]
\label{rem:zero_capacity_support}
The compact support
\begin{equation}
\{1\}\cup\{x_\gamma:\gamma\in\Gamma_+^*\}
\label{eq:polar_support}
\end{equation}
is countable and therefore has logarithmic capacity zero. This places the two measures outside the usual positive-capacity regular regime of orthogonal-polynomial asymptotics; see Stahl--Totik \cite{StahlTotik1992} and Totik's survey \cite{Totik2005Orthogonal}. Zero carrier capacity alone does not impose a universal regular $n$th-root law for the leading coefficients; the geometry and the weights of a discrete measure remain decisive. Discrete orthogonal polynomials on infinite node sets with an accumulation point form a classical subject in their own right; see \cite{WuLinXuZhao2016} for a modern asymptotic treatment in a specific model. Under the change of variable $u=(1-x)^{-1}$, the endpoint law $1-x(t)\asymp t^{-\beta}$ becomes $u\asymp t^{\beta}$, so Theorem~\ref{thm:endpoint_exponent_law} may also be viewed as a power-counting statement for a Stieltjes-type discrete moment problem. The comparison with Monien and Haynes--Zudilin in Remark~\ref{rem:monien_haynes_position} is therefore essential: the present claim is not that $N^2\log N$ decay is new, but that the Li--Weil coefficient is fixed by the Cayley endpoint exponent and can be proved from the coarse hypotheses above.
\end{remark}

\section{\texorpdfstring{$LDL^T$}{LDL\textasciicircum T}, LU, Cholesky, and Schur complements}
\label{sec:factorizations}

Set
\begin{equation}
\Delta_N^E:=\det\E_N,
\qquad
\Delta_N^O:=\det\Omat_N,
\qquad
\Delta_0^E=\Delta_0^O:=1.
\label{eq:delta_definitions}
\end{equation}

\begin{theorem}[$LDL^T$ decomposition]
\label{thm:LDL}
Under $\RH$, there are unique unit lower-triangular matrices $L_E,L_O$ and positive diagonal matrices $D_E,D_O$ such that
\begin{equation}
\E_N=L_E D_E L_E^T,
\qquad
\Omat_N=L_O D_O L_O^T.
\label{eq:LDL_decompositions}
\end{equation}
Their diagonal pivots are
\begin{align}
(D_E)_{nn}&=\frac{\Delta_n^E}{\Delta_{n-1}^E},
\label{eq:E_LDL_pivots}\\
(D_O)_{nn}&=\frac{\Delta_n^O}{\Delta_{n-1}^O}.
\label{eq:O_LDL_pivots}
\end{align}
\end{theorem}

\begin{proof}
We prove the statement for a generic real symmetric positive definite matrix $A_N$ whose leading principal blocks are denoted $A_n$, $1\le n\le N$; the two reflection families then follow from Theorem~\ref{thm:main_criterion}.  Set $A_0$ equal to the empty matrix and $\det A_0:=1$.

We argue by induction on $N$.  For $N=1$, write $A_1=(a_{11})$.  Positive definiteness gives $a_{11}>0$, and the factorization $A_1=(1)(a_{11})(1)$ is unique.  Assume the result for $A_{N-1}$ and write
\begin{equation}
A_N=\begin{pmatrix}A_{N-1}&r\\ r^T&\alpha\end{pmatrix}.
\end{equation}
Since $A_N\succ0$, its principal block $A_{N-1}$ is positive definite and hence invertible.  By the induction hypothesis,
$A_{N-1}=L_{N-1}D_{N-1}L_{N-1}^T$ with $L_{N-1}$ unit lower triangular and $D_{N-1}$ positive diagonal.  Put
\begin{equation}
y:=L_{N-1}^{-1}r,
\qquad
 d_N:=\alpha-r^TA_{N-1}^{-1}r,
\end{equation}
and define
\begin{equation}
L_N:=\begin{pmatrix}L_{N-1}&0\\ y^TD_{N-1}^{-1}&1\end{pmatrix},
\qquad
D_N:=\begin{pmatrix}D_{N-1}&0\\0&d_N\end{pmatrix}.
\end{equation}
A direct block multiplication gives $L_ND_NL_N^T=A_N$.  Moreover,
\begin{equation}
d_N=
\begin{pmatrix}-A_{N-1}^{-1}r\\1\end{pmatrix}^{T}
A_N
\begin{pmatrix}-A_{N-1}^{-1}r\\1\end{pmatrix}>0,
\end{equation}
so $D_N$ is positive diagonal.  Uniqueness follows recursively: in any factorization $A_N=LDL^T$ with $L$ unit lower triangular and $D$ diagonal, the leading $(N-1)\times(N-1)$ block fixes $L_{N-1},D_{N-1}$ by induction; then the last column below the diagonal is fixed by $r=L_{N-1}D_{N-1}\ell$, and the last diagonal entry is fixed by the $(N,N)$ entry.

Finally, $\det L_n=1$, so
\begin{equation}
\det A_n=\det D_n=\prod_{j=1}^{n}d_j.
\end{equation}
Dividing the identities for $n$ and $n-1$ yields $d_n=\det A_n/\det A_{n-1}$.  Taking $A_n=\E_n$ and $A_n=\Omat_n$ proves \eqref{eq:E_LDL_pivots}--\eqref{eq:O_LDL_pivots}.
\end{proof}

The first pivots are
\begin{align}
d_{E,1}&=\lambda_2,
\label{eq:E_first_pivot}\\
d_{E,2}&=\lambda_4-
\frac{(\lambda_3-\lambda_1)^2}{\lambda_2},
\label{eq:E_second_pivot}\\
d_{O,1}&=\lambda_1,
\label{eq:O_first_pivot}\\
d_{O,2}&=\lambda_3-
\frac{(\lambda_2-\lambda_1)^2}{\lambda_1}.
\label{eq:O_second_pivot}
\end{align}

The ordinary LU factorization with unit diagonal in the lower factor is
\begin{equation}
\E_N=L_EU_E,
\qquad U_E=D_EL_E^T,
\label{eq:E_LU}
\end{equation}
and similarly
\begin{equation}
\Omat_N=L_OU_O,
\qquad U_O=D_OL_O^T.
\label{eq:O_LU}
\end{equation}
Thus the diagonal of $U_E$ and $U_O$ is exactly the pivot sequence in \eqref{eq:E_LDL_pivots}--\eqref{eq:O_LDL_pivots}.

The Cholesky factors are
\begin{equation}
\E_N=R_ER_E^T,
\qquad R_E=L_ED_E^{1/2},
\label{eq:E_cholesky}
\end{equation}
and
\begin{equation}
\Omat_N=R_OR_O^T,
\qquad R_O=L_OD_O^{1/2}.
\label{eq:O_cholesky}
\end{equation}
In particular,
\begin{align}
(R_E)_{nn}&=\sqrt{\frac{\Delta_n^E}{\Delta_{n-1}^E}},
\label{eq:E_cholesky_diagonal}\\
(R_O)_{nn}&=\sqrt{\frac{\Delta_n^O}{\Delta_{n-1}^O}}.
\label{eq:O_cholesky_diagonal}
\end{align}

\begin{proposition}[Schur-complement recursion]
\label{prop:schur_recursion}
Write
\begin{equation}
\E_N=
\begin{pmatrix}
\E_{N-1}&r_N^E\\
(r_N^E)^T&\lambda_{2N}
\end{pmatrix},
\qquad
\Omat_N=
\begin{pmatrix}
\Omat_{N-1}&r_N^O\\
(r_N^O)^T&\lambda_{2N-1}
\end{pmatrix}.
\label{eq:block_decomposition}
\end{equation}
Then under $\RH$,
\begin{align}
\frac{\Delta_N^E}{\Delta_{N-1}^E}
&=\lambda_{2N}-(r_N^E)^T\E_{N-1}^{-1}r_N^E>0,
\label{eq:E_schur_complement}\\
\frac{\Delta_N^O}{\Delta_{N-1}^O}
&=\lambda_{2N-1}-(r_N^O)^T\Omat_{N-1}^{-1}r_N^O>0.
\label{eq:O_schur_complement}
\end{align}
\end{proposition}

\begin{proof}
We again treat a generic block matrix
\begin{equation}
A_N=\begin{pmatrix}A&r\\r^T&\alpha\end{pmatrix},
\end{equation}
where $A=A_{N-1}$ is positive definite.  In particular $A$ is invertible.  Consider the two block triangular matrices
\begin{equation}
L:=\begin{pmatrix}I&0\\-r^TA^{-1}&1\end{pmatrix},
\qquad
R:=\begin{pmatrix}I&-A^{-1}r\\0&1\end{pmatrix}.
\end{equation}
Both have determinant one.  Direct multiplication gives
\begin{equation}
LA_NR=
\begin{pmatrix}
A&0\\
0&\alpha-r^TA^{-1}r
\end{pmatrix}.
\end{equation}
Taking determinants therefore yields
\begin{equation}
\det A_N=\det A\,\bigl(\alpha-r^TA^{-1}r\bigr),
\end{equation}
and, since $\det A>0$,
\begin{equation}
\frac{\det A_N}{\det A}=\alpha-r^TA^{-1}r.
\end{equation}
The scalar on the right is strictly positive because the nonzero vector
\begin{equation}
z:=\begin{pmatrix}-A^{-1}r\\1\end{pmatrix}
\end{equation}
satisfies
\begin{equation}
z^TA_Nz=\alpha-r^TA^{-1}r>0
\end{equation}
by positive definiteness of $A_N$.  Substituting $A=\E_{N-1}$, $r=r_N^E$, $\alpha=\lambda_{2N}$ and then $A=\Omat_{N-1}$, $r=r_N^O$, $\alpha=\lambda_{2N-1}$ proves both formulas and their strict positivity.
\end{proof}

\begin{remark}[Geometric meaning]
\label{rem:schur_geometry}
The quantities in \eqref{eq:E_schur_complement} and \eqref{eq:O_schur_complement} are the squared residual norms after orthogonally projecting the new sine mode onto the span of all preceding modes. They quantify the genuinely new Gram energy contributed by frequency $N$ or $N-1/2$.
\end{remark}

\section{Inverse matrices and reproducing kernels}
\label{sec:inverse}

Whenever the matrices are invertible, the adjugate formula gives
\begin{align}
(\E_N^{-1})_{ij}
&=(-1)^{i+j}\frac{\det \E_N^{(j,i)}}{\Delta_N^E},
\label{eq:E_inverse_adjugate}\\
(\Omat_N^{-1})_{ij}
&=(-1)^{i+j}\frac{\det \Omat_N^{(j,i)}}{\Delta_N^O}.
\label{eq:O_inverse_adjugate}
\end{align}
More structure is visible after the moment-basis transformation.

Let
\begin{align}
u_N(x)&:=(U_0(x),\dots,U_{N-1}(x))^T,
\label{eq:uN_vector}\\
\mathbf p_N(x)&:=(P_0(x),\dots,P_{N-1}(x))^T.
\label{eq:pN_vector}
\end{align}
Define
\begin{align}
\mathcal K_N^E(x,y)
&:=u_N(x)^T\E_N^{-1}u_N(y),
\label{eq:KE_definition}\\
\mathcal K_N^O(x,y)
&:=\mathbf p_N(x)^T\Omat_N^{-1}\mathbf p_N(y).
\label{eq:KO_definition}
\end{align}

\begin{theorem}[Christoffel--Darboux interpretation]
\label{thm:inverse_CD}
Under $\RH$, $\mathcal K_N^E$ and $\mathcal K_N^O$ are the reproducing kernels of the polynomial spaces of degree $<N$ in $L^2(\mu_E)$ and $L^2(\mu_O)$, respectively. If $\{q_j^E\}$ and $\{q_j^O\}$ are the corresponding orthonormal polynomials, then
\begin{align}
\mathcal K_N^E(x,y)&=\sum_{j=0}^{N-1}q_j^E(x)q_j^E(y),
\label{eq:KE_orthonormal_sum}\\
\mathcal K_N^O(x,y)&=\sum_{j=0}^{N-1}q_j^O(x)q_j^O(y).
\label{eq:KO_orthonormal_sum}
\end{align}
\end{theorem}

\begin{proof}
We prove the even statement.  Let
\begin{equation}
\mathcal P_{N-1}:=\operatorname{span}\{U_0,\ldots,U_{N-1}\}
\end{equation}
inside $L^2(\mu_E)$.  Under $\RH$, $\E_N\succ0$, so these basis functions are linearly independent in $L^2(\mu_E)$ and their Gram matrix is invertible.  For fixed $y$, define
\begin{equation}
K_y(x):=u_N(x)^T\E_N^{-1}u_N(y)\in\mathcal P_{N-1}.
\end{equation}
Take an arbitrary $f\in\mathcal P_{N-1}$ and write $f(x)=c^Tu_N(x)$ with $c\in\mathbb R^N$.  Then, using the moment Gram identity,
\begin{align}
\langle f,K_y\rangle_{L^2(\mu_E)}
&=\int c^Tu_N(x)u_N(x)^T\E_N^{-1}u_N(y)\,\dd\mu_E(x)\\
&=c^T\left(\int u_N(x)u_N(x)^T\,\dd\mu_E(x)\right)\E_N^{-1}u_N(y)\\
&=c^T\E_N\E_N^{-1}u_N(y)=f(y).
\end{align}
Thus $K_y$ has the reproducing property and is the unique reproducing kernel section at $y$.  This proves \eqref{eq:KE_definition} as a reproducing kernel; the odd case is identical with $p_N$ and $\Omat_N$.

Now let $q_0^E,\ldots,q_{N-1}^E$ be any orthonormal basis of $\mathcal P_{N-1}$.  The function
\begin{equation}
\widetilde K_N^E(x,y):=\sum_{j=0}^{N-1}q_j^E(x)q_j^E(y)
\end{equation}
has the same reproducing property, because for
$f=\sum_j\widehat f_jq_j^E$ one has
$\langle f,\widetilde K_N^E(\cdot,y)\rangle=\sum_j\widehat f_jq_j^E(y)=f(y)$.  Uniqueness of the Riesz representer in the finite-dimensional Hilbert space therefore gives
$\widetilde K_N^E=\mathcal K_N^E$.  This is \eqref{eq:KE_orthonormal_sum}; the odd formula follows in the same way.
\end{proof}

In the monomial basis, Theorem \ref{thm:hankel_congruence} gives the exact inverse identities
\begin{align}
\E_N^{-1}
&=(B_E^{(N)})^{-T}(H_N^E)^{-1}(B_E^{(N)})^{-1},
\label{eq:E_inverse_hankel}\\
\Omat_N^{-1}
&=(B_O^{(N)})^{-T}(H_N^O)^{-1}(B_O^{(N)})^{-1}.
\label{eq:O_inverse_hankel}
\end{align}
Together with Corollary \ref{cor:inverse_displacement}, these formulas show that the inverses simultaneously possess moment-kernel and low-displacement-rank descriptions.

\begin{remark}[No universal sign pattern]
\label{rem:no_M_matrix}
The off-diagonal entries $\lambda_{m+n}-\lambda_{|m-n|}$ and $\lambda_{m+n-1}-\lambda_{|m-n|}$ do not have an algebraically forced sign. Consequently the matrices are not naturally Stieltjes $M$-matrices, and no universal checkerboard or nonpositive off-diagonal sign pattern for their inverses follows from the present theory.
\end{remark}


\section{Column recursion and the closed two-generator bordering scheme}
\label{sec:two_generator}

The central matrix-first problem is whether nesting and rank-two displacement can turn the
Schur complements \eqref{eq:E_schur_complement} into a propagation law whose sign can be
controlled.  This section solves the recurrence part in closed form.  What remains is the
sign mechanism isolated later in Section~\ref{sec:route_A_even};
Remark~\ref{rem:recursion_is_not_a_sign_mechanism} records the distinction.

Throughout, $\lambda_0=0$ and $\lambda_{-j}=\lambda_j$.  Every numerical statement in this section and in Section~\ref{sec:christoffel_bridge} uses the Li coefficients produced by the arithmetic construction of Subsection~\ref{subsec:arithmetic_reconstruction}, which is stated later only because it belongs with the archimedean--arithmetic split; it takes Stieltjes constants as input and uses no zero of $\zeta$ at any stage.

\begin{definition}[Bordering column]
\label{def:bordering_column}
For $N\ge1$ let $\mathsf b^{(N)}$ be the semi-infinite vector
\begin{equation}
 \mathsf b^{(N)}_m:=\lambda_{m+N}-\lambda_{|m-N|},
 \qquad m\ge1 .
 \label{eq:bordering_column}
\end{equation}
Its first $N-1$ entries form the off-diagonal block of $\E_N$ in
\eqref{eq:block_decomposition}, and $\mathsf b^{(N)}_N=\lambda_{2N}$.
\end{definition}

\begin{proposition}[Column wave equation]
\label{prop:column_wave_equation}
Let $T$ denote the formal nearest-neighbour shift on semi-infinite sequences with Dirichlet
condition at $0$, that is $(Tv)_m=v_{m+1}+v_{m-1}$ with $v_0:=0$.  No $\ell^2$ assumption is
made on the bordering columns.  Then for every $N\ge2$,
\begin{equation}
 \boxed{\;
 \mathsf b^{(N+1)}+\mathsf b^{(N-1)}=T\,\mathsf b^{(N)} \;}
 \label{eq:column_wave_equation}
\end{equation}
as an identity between semi-infinite vectors, valid unconditionally.  Equivalently,
\begin{equation}
 \mathsf b^{(N)}=U_{N-1}\!\left(\tfrac{T}{2}\right)\mathsf b^{(1)},
 \label{eq:column_chebyshev}
\end{equation}
where $U_n$ is the Chebyshev polynomial of the second kind.
\end{proposition}

\begin{proof}
Fix $m\ge1$.  The left-hand side of \eqref{eq:column_wave_equation} at index $m$ is
\begin{equation}
 \bigl(\lambda_{m+N+1}-\lambda_{|m-N-1|}\bigr)
 +\bigl(\lambda_{m+N-1}-\lambda_{|m-N+1|}\bigr),
\end{equation}
and the right-hand side is
\begin{equation}
 \bigl(\lambda_{m+1+N}-\lambda_{|m+1-N|}\bigr)
 +\bigl(\lambda_{m-1+N}-\lambda_{|m-1-N|}\bigr),
\end{equation}
where the convention $v_0=0$ is automatic because $\mathsf b^{(N)}_0=\lambda_N-\lambda_N=0$.
The two positive parts agree termwise.  The two negative parts agree as well, since
$|m-N-1|=|m-1-N|$ and $|m-N+1|=|m+1-N|$.  This proves
\eqref{eq:column_wave_equation}.  Since $\mathsf b^{(0)}=0$, the recursion
$\mathsf b^{(N+1)}=T\mathsf b^{(N)}-\mathsf b^{(N-1)}$ with initial data
$\mathsf b^{(0)}=0$, $\mathsf b^{(1)}$ is exactly the Chebyshev recursion, giving
\eqref{eq:column_chebyshev}.
\end{proof}

\begin{remark}[Interpretation]
\label{rem:column_wave_meaning}
Under $\RH$ the identity \eqref{eq:column_wave_equation} is the statement
$\sin((N+1)\theta)+\sin((N-1)\theta)=2\cos\theta\,\sin(N\theta)$ transported through the
Gram representation \eqref{eq:intro_gram_E}.  The content of
Proposition~\ref{prop:column_wave_equation} is that it holds \emph{unconditionally}, as an
identity in the Li coefficients: the same discrete wave equation that
Section~\ref{sec:displacement} exhibits in the row index also governs the sequence of
bordering columns.  The pair $\bigl(\mathsf b^{(1)},\mathsf b^{(2)}\bigr)$ therefore
generates the entire semi-infinite array, which is the sense in which the scheme below is
a \emph{two-generator} scheme.
\end{remark}

\begin{proposition}[Closed bordering recursion]
\label{prop:closed_bordering}
Assume $\E_{N-1}$ is invertible and set $G:=\E_{N-1}^{-1}$.  Write $e:=e_{N-1}\in\mathbb R^{N-1}$
for the last standard basis vector, let
\begin{equation}
 \omega_N:=\lambda_{2N-1}-\lambda_1,
 \qquad
 r^{(N)}:=\bigl(\mathsf b^{(N)}_1,\dots,\mathsf b^{(N)}_{N-2},0\bigr)^T,
 \label{eq:omegaN_and_r}
\end{equation}
so that the off-diagonal block of $\E_N$ is
\begin{equation}
 b_N=r^{(N)}+\omega_N\,e .
 \label{eq:b_split}
\end{equation}
Put $g:=Ge$, $h:=Gr^{(N)}$ and $q:=h+\omega_Ng$.  Then the next Schur pivot is
\begin{equation}
 p_N:=\frac{\Delta_N^E}{\Delta_{N-1}^E}
 =\lambda_{2N}-\bigl(r^{(N)}\bigr)^Th-2\omega_N h_{N-1}-\omega_N^2 g_{N-1}.
 \label{eq:pN_two_generator}
\end{equation}
If, in addition, \(p_N\ne0\)---equivalently, \(\E_N\) is invertible---then
\begin{equation}
 \E_N^{-1}=
 \begin{pmatrix}
 G+p_N^{-1}qq^T & -p_N^{-1}q\\[2pt]
 -p_N^{-1}q^T & p_N^{-1}
 \end{pmatrix},
 \qquad
 g^{(N)}:=\E_N^{-1}e_N=\begin{pmatrix}-q/p_N\\ 1/p_N\end{pmatrix}.
 \label{eq:gN_two_generator}
\end{equation}
In particular \(g^{(N)}_N\,p_N=1\) whenever the inverse exists.
\end{proposition}

\begin{proof}
Equation \eqref{eq:b_split} is the definition of $r^{(N)}$ together with
$\mathsf b^{(N)}_{N-1}=\lambda_{2N-1}-\lambda_1=\omega_N$.  Since $G$ is symmetric,
\begin{equation}
 b_N^TGb_N
 =\bigl(r^{(N)}+\omega_Ne\bigr)^TG\bigl(r^{(N)}+\omega_Ne\bigr)
 =\bigl(r^{(N)}\bigr)^Th+2\omega_N\,e^Th+\omega_N^2\,e^Tg,
\end{equation}
and $e^Th=h_{N-1}$, $e^Tg=g_{N-1}$.  Substituting into the Schur complement
\eqref{eq:E_schur_complement} gives \eqref{eq:pN_two_generator}.  Formula
\eqref{eq:gN_two_generator} is the standard block inverse of
$\begin{psmallmatrix}\E_{N-1}&b_N\\b_N^T&\lambda_{2N}\end{psmallmatrix}$ with
$q=Gb_N=h+\omega_Ng$, which is valid whenever $p_N\ne0$.
\end{proof}


\begin{proposition}[Closed two-vector state update]
\label{prop:two_vector_state_update}
Under the hypotheses and notation of Proposition~\ref{prop:closed_bordering}, assume in addition
that \(p_N\ne0\), so that \(\E_N^{-1}\) exists, and let
\begin{equation}
 \rho_N:=\bigl(\mathsf b^{(N+1)}_1,\dots,\mathsf b^{(N+1)}_{N-1}\bigr)^T\in\mathbb R^{N-1}.
 \label{eq:rhoN_two_vector}
\end{equation}
Then the column wave equation gives
\begin{equation}
 \boxed{\;
 \rho_N=2\J_E^{(N-1)}b_N-\E_{N-1}e+\lambda_{2N}e .\;}
 \label{eq:rhoN_wave}
\end{equation}
If $q=h+\omega_Ng=\E_{N-1}^{-1}b_N$ and
\begin{equation}
 \eta_N:=-\frac{q^T\rho_N}{p_N},
 \label{eq:etaN_two_vector}
\end{equation}
then the second inverse generator
\begin{equation}
 h^{(N)}:=\E_N^{-1}r_E^{(N)},
 \qquad
 r_E^{(N)}=(\rho_N^T,0)^T,
\end{equation}
is obtained without a new inverse column:
\begin{equation}
 \boxed{
 h^{(N)}=
 \begin{pmatrix}
 2\J_E^{(N-1)}q+(q_{N-1}-\eta_N)q+p_Ng-e\\[2pt]
 \eta_N
 \end{pmatrix}.}
 \label{eq:hN_two_vector_update}
\end{equation}
Together with \eqref{eq:gN_two_generator}, equations
\eqref{eq:pN_two_generator} and \eqref{eq:hN_two_vector_update} close the inverse information
needed at the next step on the pair $(g^{(N)},h^{(N)})$ and the bordering columns.
\end{proposition}

\begin{proof}
The first identity is the finite restriction of
$\mathsf b^{(N+1)}+\mathsf b^{(N-1)}=T\mathsf b^{(N)}$ at the new boundary; the missing
$N$-th coupling contributes the last two terms in \eqref{eq:rhoN_wave}.  For completeness,
the inverse displacement identity gives
\begin{equation}
 \J_E^{(N-1)}G-G\J_E^{(N-1)}
 =\frac12(gh^T-hg^T).
 \label{eq:inverse_displacement_gh_local}
\end{equation}
Apply $G$ to \eqref{eq:rhoN_wave}, substitute $q=Gb_N=h+\omega_Ng$, and use
$p_N=\lambda_{2N}-b_N^Tq$.  The scalar terms combine to give
\begin{equation}
 G\rho_N=2\J_E^{(N-1)}q+q_{N-1}q+p_Ng-e.
 \label{eq:Grho_two_vector}
\end{equation}
Finally multiply $r_E^{(N)}=(\rho_N^T,0)^T$ by the block inverse
\eqref{eq:gN_two_generator}.  Its last component is $-q^T\rho_N/p_N=\eta_N$, and its top
block is $G\rho_N-\eta_Nq$.  Equation \eqref{eq:hN_two_vector_update} follows.
\end{proof}

\begin{theorem}[Local $LDL^T$ pivot dynamics]
\label{thm:local_ldl_pivot_ratio}
Let $A$ denote either the even or the odd semi-infinite reflection array, and suppose
$p_1,\dots,p_n$ are nonzero.  Gaussian $LDL^T$ elimination is then well defined through
column $n$; in particular the leading block $A_{n+1}$ has the factorization
\begin{equation}
 A_{n+1}=L_{n+1}\,\diag(p_1,\dots,p_{n+1})\,L_{n+1}^T,
 \qquad (L_{n+1})_{jj}=1,
 \label{eq:local_LDL_factor}
\end{equation}
where $p_{n+1}$ is allowed to vanish.  The $n$th elimination column in the next
leading block is also defined because its construction divides only by
$p_1,\dots,p_n$.  Define the first two subdiagonals
\begin{equation}
 \ell_n:=(L_{n+1})_{n+1,n},
 \qquad
 \tau_n:=(L_{n+2})_{n+2,n},
 \qquad n\ge1.
 \label{eq:ell_tau_def}
\end{equation}
Set
\begin{equation}
 \ell_0=\sigma,
 \qquad \tau_0=0,
 \qquad
 \sigma=\begin{cases}0,&A=\E,\\1,&A=\Omat.\end{cases}
 \label{eq:ell_tau_boundary}
\end{equation}
Then, for every $n\ge1$,
\begin{equation}
 \boxed{\;
 \frac{p_{n+1}}{p_n}
 =1+(\tau_n-\tau_{n-1})
 -\ell_n(\ell_n-\ell_{n-1}).\;}
 \label{eq:local_pivot_ratio}
\end{equation}
Equivalently, if
\begin{equation}
 \mathcal H_n:=\tau_n+n-\frac12\ell_n^2
 -\frac12\sum_{j=1}^{n}(\ell_j-\ell_{j-1})^2,
 \qquad \mathcal H_0=-\frac12\sigma^2,
 \label{eq:Hlocal_def}
\end{equation}
then
\begin{equation}
 \boxed{\;\mathcal H_n-\mathcal H_{n-1}=\frac{p_{n+1}}{p_n}.\;}
 \label{eq:Hlocal_increment}
\end{equation}
Thus, once $p_n>0$ is known, positivity of the next pivot is equivalent to the strict
increase $\mathcal H_n>\mathcal H_{n-1}$.
\end{theorem}

\begin{proof}
First suppose $p_{n+1}\ne0$.  We give the even case.  On the semi-infinite reflection array the discrete wave
identity is the commutation relation $\J_EA=A\J_E$.  Insert
$A=LDL^T$ and put $B=L^{-1}\J_EL$.  Since $\J_E$ is tridiagonal and $L$ is unit lower
triangular, $B$ is upper Hessenberg.  The relation $BD=DB^T$ forces it to be tridiagonal.
Moreover
\begin{equation}
 2B_{n-1,n}=1,
 \qquad
 2B_{n,n}=\ell_n-\ell_{n-1},
 \qquad
 2B_{n+1,n}=\frac{p_{n+1}}{p_n}.
 \label{eq:B_local_entries}
\end{equation}
The identity $\J_EL=LB$, read at entry $(m,n)$, therefore gives
\begin{equation}
 L_{m+1,n}+L_{m-1,n}
 =L_{m,n-1}+(\ell_n-\ell_{n-1})L_{m,n}
 +\frac{p_{n+1}}{p_n}L_{m,n+1}.
 \label{eq:L_discrete_wave}
\end{equation}
Set $m=n+1$.  Since $L_{n,n}=L_{n+1,n+1}=1$,
$L_{n+2,n}=\tau_n$ and $L_{n+1,n-1}=\tau_{n-1}$, this is exactly
\eqref{eq:local_pivot_ratio}.  In the odd family the only change is the $-1/2$ entry at the
upper-left corner of $\J_O$; the same calculation is recovered by the boundary convention
$\ell_0=1$.  Finally,
\begin{equation}
 \ell_n(\ell_n-\ell_{n-1})
 =\frac12(\ell_n^2-\ell_{n-1}^2)
 +\frac12(\ell_n-\ell_{n-1})^2,
\end{equation}
which telescopes \eqref{eq:local_pivot_ratio} into \eqref{eq:Hlocal_increment}.

It remains to allow $p_{n+1}=0$.  For fixed $n$, the elimination coefficients
$\ell_j,\tau_j$ appearing above are rational functions of the reflection entries whose
denominators involve only $p_1,\dots,p_n$.  Perturb the scalar datum \(\lambda_{2n+2}\) in the ambient reflection sequence, without asserting
that the perturbed sequence is itself a Li sequence.  This leaves the previously formed pivots
\(p_1,\ldots,p_n\) unchanged, preserves the algebraic reflection identities, and makes
\(p_{n+1}\) nonzero except at one parameter value.  Both sides of
\eqref{eq:local_pivot_ratio} are continuous as the perturbation is removed.  Hence the
identity extends to $p_{n+1}=0$ as well.
\end{proof}

\begin{corollary}[Defect-energy form]
\label{cor:defect_energy_pivot}
With $\sigma$ as in Theorem~\ref{thm:local_ldl_pivot_ratio}, define
\begin{align}
 d_n&:=2n+\sigma-\ell_n,\label{eq:dn_local_def}\\
 c_n&:=2n^2+(1+2\sigma)n-\tau_n,\label{eq:cn_local_def}\\
 V_n&:=(2n+\sigma+2)d_n-\frac12d_n^2-c_n,
 \qquad d_0=c_0=V_0=0.
 \label{eq:Vn_local_def}
\end{align}
Then
\begin{equation}
 \boxed{\;
 \frac{p_{n+1}}{p_n}
 =V_n-V_{n-1}-\frac12(d_n-d_{n-1})^2.\;}
 \label{eq:defect_energy_ratio}
\end{equation}
In particular, under the inductive hypothesis $p_1,\dots,p_n>0$,
\begin{equation}
 p_{n+1}>0
 \quad\Longleftrightarrow\quad
 V_n-V_{n-1}>\frac12(d_n-d_{n-1})^2.
 \label{eq:defect_energy_sign}
\end{equation}
Equivalently, the symmetric cell
\begin{equation}
 \mathcal C_n=
 \begin{pmatrix}
 V_n-V_{n-1} & (d_n-d_{n-1})/\sqrt2\\[2pt]
 (d_n-d_{n-1})/\sqrt2 & 1
 \end{pmatrix}
 \label{eq:local_cell_def}
\end{equation}
satisfies
\begin{equation}
 \det\mathcal C_n=\frac{p_{n+1}}{p_n}.
 \label{eq:local_cell_det}
\end{equation}
Thus the next reflection pivot is controlled by a $2\times2$ local cell built only from
the first two subdiagonals of the $LDL^T$ elimination factor.
\end{corollary}

\begin{proof}
Substitute the definitions \eqref{eq:dn_local_def} and
\eqref{eq:cn_local_def} into \eqref{eq:local_pivot_ratio}, and simplify.  Equation \eqref{eq:local_cell_det} is the
Schur-complement identity for the lower-right entry $1$.
\end{proof}

\begin{corollary}[Cost]
\label{cor:two_generator_cost}
Suppose the inverse-producing steps under consideration have nonzero Schur pivots.
Given $\E_{N-1}^{-1}$ and $\mathsf b^{(N-1)},\mathsf b^{(N-2)}$, one step of
Propositions~\ref{prop:column_wave_equation} and \ref{prop:closed_bordering} produces
$\mathsf b^{(N)}$, $p_N$, $g^{(N)}$ and $\E_N^{-1}$ using one banded matrix--vector product
and two dense matrix--vector products, that is $O(N^2)$ operations and no determinant
evaluation.  In particular the whole pivot sequence $p_1,\dots,p_N$ is obtained in
$O(N^3)$ operations without ever forming $\Delta_N^E$.
\end{corollary}

\begin{proof}
The column-wave update is banded and costs \(O(N)\).  Forming the required inverse
generator vectors uses dense matrix--vector products with an \((N-1)\times(N-1)\)
matrix and therefore costs \(O(N^2)\); the rank-one block-inverse update also costs
\(O(N^2)\).  Thus one nonsingular bordering step is \(O(N^2)\), and summing these
costs through order \(N\) gives \(O(N^3)\).  No determinant evaluation is required,
because \(p_n\) is obtained directly from the Schur formula.
\end{proof}

\begin{experiment}[Numerical verification of the two-generator scheme]
\label{exp:two_generator}
With the arithmetic $\lambda$ of Experiment~\ref{exp:arithmetic_lambda} and $400$ decimal
working digits, the identity \eqref{eq:column_wave_equation} holds to $6.8\times10^{-400}$
for all $N\le23$ and $m\le20$, that is to working precision; it is an exact algebraic
identity.  The recursion \eqref{eq:pN_two_generator}--\eqref{eq:gN_two_generator} was run
from $N=2$ to $N=22$, updating $\E_N^{-1}$ by \eqref{eq:gN_two_generator} alone and never
recomputing an inverse.  Against direct multiprecision evaluation of
$\Delta_N^E/\Delta_{N-1}^E$ and of $\E_N^{-1}e_N$:
\begin{center}
\small
\begin{tabular}{rlll}
\toprule
$N$ & $p_N$ (recursion) & rel.\ error in $p_N$ & rel.\ error in $g^{(N)}$\\
\midrule
$2$  & $3.360255490\times10^{-7}$   & $1.6\times10^{-395}$ & $1.6\times10^{-395}$\\
$5$  & $6.620331547\times10^{-26}$  & $7.6\times10^{-356}$ & $7.6\times10^{-356}$\\
$8$  & $1.511647984\times10^{-46}$  & $4.8\times10^{-311}$ & $4.8\times10^{-311}$\\
$12$ & $1.214721178\times10^{-75}$  & $8.3\times10^{-248}$ & $8.3\times10^{-248}$\\
$16$ & $2.446315652\times10^{-106}$ & $3.8\times10^{-182}$ & $3.8\times10^{-182}$\\
$20$ & $4.565211501\times10^{-138}$ & $1.1\times10^{-113}$ & $1.1\times10^{-113}$\\
$22$ & $2.854130423\times10^{-154}$ & $7.1\times10^{-79}\phantom{0}$ & $7.1\times10^{-79}$\\
\bottomrule
\end{tabular}
\end{center}
The relative error degrades by roughly fifteen digits per step, which is the expected
conditioning cost of the bordering and matches the decay rate of $p_N$; it is not an
instability of the formulas.  The identity $g^{(N)}_Np_N=1$ holds exactly at every step.

The two generators have strictly alternating signs of opposite phase for all $N\le22$:
\begin{equation}
 \operatorname{sign}\bigl(g_i\bigr)=(-1)^{N-1-i},
 \qquad
 \operatorname{sign}\bigl(h_i\bigr)=(-1)^{N-i},
 \qquad 1\le i\le N-1,
 \label{eq:checkerboard_generators}
\end{equation}
and $q=h+\omega_Ng$ inherits the phase of $g$.  This holds with the arithmetic $\lambda$, hence
independently of any zero data.
\end{experiment}

\begin{remark}[The remaining sign problem is local]
\label{rem:recursion_is_not_a_sign_mechanism}
Propositions~\ref{prop:closed_bordering} and \ref{prop:two_vector_state_update} close the
bordering dynamics algebraically, while Theorem~\ref{thm:local_ldl_pivot_ratio} removes the
large inverse from the sign test altogether: once $p_n>0$ is known, the next step is
exactly the local inequality \eqref{eq:defect_energy_sign}.  The checkerboard pattern
\eqref{eq:checkerboard_generators} is therefore evidence about the elimination geometry,
not a substitute for the missing proof.  The unresolved task is to derive
\eqref{eq:defect_energy_sign} from structural properties of the reflection matrices under
the inductive hypothesis, without replacing the local identity by a termwise bound on the
ill-conditioned inverse.
\end{remark}

\section{Orthogonal polynomials and infinite Jacobi operators}
\label{sec:jacobi}

Assume $\RH$. Let $\pi_n^E$ and $\pi_n^O$ denote the monic orthogonal polynomials associated with $\mu_E$ and $\mu_O$.  The following elementary argument records the hypotheses behind their three-term recurrences.  For either positive measure $\mu=\mu_E$ or $\mu_O$, let $h_n=\|\pi_n\|_{L^2(\mu)}^2>0$.  Since $x\pi_n$ is monic of degree $n+1$, it has an expansion
\begin{equation}
x\pi_n=\pi_{n+1}+\sum_{j=0}^{n}c_{n,j}\pi_j.
\end{equation}
For $j\le n-2$, symmetry of multiplication by the real variable $x$ gives
\begin{equation}
\langle x\pi_n,\pi_j\rangle
=\langle\pi_n,x\pi_j\rangle=0,
\end{equation}
because $x\pi_j$ has degree at most $n-1$.  Thus only the $j=n$ and $j=n-1$ coefficients can remain.  Moreover
\begin{equation}
c_{n,n-1}
=\frac{\langle x\pi_n,\pi_{n-1}\rangle}{h_{n-1}}
=\frac{\langle\pi_n,x\pi_{n-1}\rangle}{h_{n-1}}
=\frac{h_n}{h_{n-1}}>0,
\end{equation}
because $x\pi_{n-1}=\pi_n+$ lower-degree terms.  Writing
$b_n:=c_{n,n}$ and $(a_n)^2:=h_n/h_{n-1}$ gives
\begin{align}
x\pi_n^E(x)
&=\pi_{n+1}^E(x)+b_n^E\pi_n^E(x)+(a_n^E)^2\pi_{n-1}^E(x),
\label{eq:E_monic_recurrence}\\
x\pi_n^O(x)
&=\pi_{n+1}^O(x)+b_n^O\pi_n^O(x)+(a_n^O)^2\pi_{n-1}^O(x),
\label{eq:O_monic_recurrence}
\end{align}
with $a_n^E,a_n^O>0$ for $n\ge1$ and the convention $\pi_{-1}=0$.  After normalizing $q_n=\pi_n/\sqrt{h_n}$, multiplication by $x$ has the symmetric tridiagonal matrix with diagonal $b_n$ and positive off-diagonal $a_{n+1}$.  Because the measures are supported in $[-1,1]$, multiplication by $x$ is a bounded self-adjoint operator of norm at most one on each $L^2$ space; once polynomial completeness is established below, the corresponding Jacobi matrices define bounded self-adjoint operators $\mathcal J_E$ and $\mathcal J_O$ unitarily equivalent to multiplication by $x$.

The first recurrence coefficients are explicit in the Li coefficients:
\begin{align}
b_0^E
&=\frac{\lambda_3-\lambda_1}{2\lambda_2},
\label{eq:b0E}\\
(a_1^E)^2
&=\frac{\lambda_2\lambda_4-(\lambda_3-\lambda_1)^2}{4\lambda_2^2},
\label{eq:a1E}\\
b_0^O
&=\frac{\lambda_2-2\lambda_1}{2\lambda_1},
\label{eq:b0O}\\
(a_1^O)^2
&=\frac{\lambda_1\lambda_3-(\lambda_2-\lambda_1)^2}{4\lambda_1^2}.
\label{eq:a1O}
\end{align}

\begin{proposition}[Jacobi off-diagonals from the principal determinants]
\label{prop:jacobi_from_determinants}
Let
\begin{equation}
d_n^E:=\frac{\Delta_n^E}{\Delta_{n-1}^E},
\qquad
d_n^O:=\frac{\Delta_n^O}{\Delta_{n-1}^O}.
\label{eq:pivot_short_notation}
\end{equation}
Under $\RH$, for every $n\ge1$,
\begin{align}
(a_n^E)^2
&=\frac14\frac{d_{n+1}^E}{d_n^E}
=\frac14\frac{\Delta_{n+1}^E\Delta_{n-1}^E}{(\Delta_n^E)^2},
\label{eq:anE_from_determinants}\\
(a_n^O)^2
&=\frac14\frac{d_{n+1}^O}{d_n^O}
=\frac14\frac{\Delta_{n+1}^O\Delta_{n-1}^O}{(\Delta_n^O)^2}.
\label{eq:anO_from_determinants}
\end{align}
\end{proposition}

\begin{proof}
The basis polynomial $U_{n-1}$, and also $P_{n-1}$, has leading coefficient $2^{n-1}$.  Therefore the squared norm $h_{n-1}$ of the corresponding monic orthogonal polynomial is
\begin{equation}
h_{n-1}^E=\frac{d_n^E}{4^{n-1}},
\qquad
h_{n-1}^O=\frac{d_n^O}{4^{n-1}}.
\label{eq:monic_norm_from_pivot}
\end{equation}
For monic orthogonal polynomials $(a_n)^2=h_n/h_{n-1}$, which yields \eqref{eq:anE_from_determinants}--\eqref{eq:anO_from_determinants}.
\end{proof}

Let $\Gamma_+^*$ denote the set of \emph{distinct} positive ordinates of nontrivial zeros and let $m_\gamma$ denote multiplicity.

The measures $\mu_E$ and $\mu_O$ are convenient unnormalized representing measures. The spectral measure of the first Jacobi basis vector is instead the corresponding probability measure
\begin{equation}
\nu_E:=\frac{\mu_E}{\lambda_2},
\qquad
\nu_O:=\frac{\mu_O}{\lambda_1},
\label{eq:normalized_spectral_measures}
\end{equation}
because $\mu_E([-1,1])=\lambda_2$ and $\mu_O([-1,1])=\lambda_1$. This normalization does not alter the orthogonal polynomials except for their orthonormal scaling, and it does not alter the associated Jacobi operator.

\begin{theorem}[Exact spectrum of the infinite Jacobi operators]
\label{thm:jacobi_zero_support}
Assume $\RH$. Then
\begin{equation}
\spec(\mathcal J_E)=\spec(\mathcal J_O)
=\{1\}\cup\{x_\gamma:\gamma\in\Gamma_+^*\},
\label{eq:jacobi_support}
\end{equation}
where
\begin{equation}
x_\gamma=\frac{4\gamma^2-1}{4\gamma^2+1}.
\label{eq:jacobi_xgamma_repeat}
\end{equation}
Each $x_\gamma$ is an isolated eigenvalue of the multiplication operator with spectral mass determined by $m_\gamma$, while $1$ is the unique accumulation point of the spectrum. The inverse transformation is
\begin{equation}
\gamma=\frac12\sqrt{\frac{1+x}{1-x}}
\qquad(-1<x<1).
\label{eq:gamma_from_x}
\end{equation}
Moreover,
\begin{equation}
I-\mathcal J_E\in\mathcal S_1,
\qquad
I-\mathcal J_O\in\mathcal S_1,
\label{eq:Jacobi_trace_class}
\end{equation}
so both are trace-class perturbations of the identity and have essential spectrum $\{1\}$.
\end{theorem}

\begin{proof}
We treat either measure $\mu=\mu_E$ or $\mu_O$; the support is the same in both cases.  Every atom has positive mass, and
$x_\gamma\to1$ as $\gamma\to\infty$.  Hence the topological support is exactly
\begin{equation}
K:=\operatorname{supp}\mu
=\{x_\gamma:\gamma\in\Gamma_+^*\}\cup\{1\}.
\label{eq:measure_support_exact}
\end{equation}
There is no atom at $1$; it is present because it is the unique accumulation point of the atoms.

Because $K$ is a compact subset of $\mathbb R$, the restrictions of algebraic polynomials form a unital subalgebra of $C(K)$ that separates points.  The Stone--Weierstrass theorem therefore makes polynomials uniformly dense in $C(K)$.  Continuous functions are dense in $L^2(\mu)$ for every finite Borel measure on a compact metric space, hence polynomials are dense in $L^2(\mu)$.  Gram--Schmidt applied to $1,x,x^2,\ldots$ therefore produces a complete orthonormal polynomial basis $\{p_n\}_{n\ge0}$.

Let $\mathcal F:\ell^2(\mathbb N_0)\to L^2(\mu)$ send the standard vector $e_n$ to $p_n$.  Completeness makes $\mathcal F$ unitary.  Multiplication by $x$ is a bounded self-adjoint operator $M_x$ with $\|M_x\|\le1$.  Its matrix in the orthonormal polynomial basis is the Jacobi matrix $\mathcal J$ coming from the three-term recurrence, so
\begin{equation}
\mathcal J=\mathcal F^{-1}M_x\mathcal F.
\label{eq:Jacobi_multiplication_equivalence}
\end{equation}

Now identify $L^2(\mu)$ explicitly with a weighted sequence space.  If $w_\gamma=\mu(\{x_\gamma\})$, then
\begin{equation}
Tf:=\bigl(\sqrt{w_\gamma}\,f(x_\gamma)\bigr)_{\gamma\in\Gamma_+^*}
\end{equation}
is unitary from $L^2(\mu)$ onto $\ell^2(\Gamma_+^*)$, and
\begin{equation}
TM_xT^{-1}=\operatorname{diag}(x_\gamma).
\end{equation}
The spectrum of a bounded diagonal operator is the closure of its diagonal entries.  Since the only accumulation point is $1$,
\begin{equation}
\operatorname{spec}(\mathcal J)=\{x_\gamma:\gamma\in\Gamma_+^*\}\cup\{1\}.
\end{equation}
Each $x_\gamma$ is an isolated eigenvalue with one-dimensional eigenspace generated by the normalized atom indicator.  The point $1$ belongs to the spectrum as a limit point but is not an eigenvalue because $\mu(\{1\})=0$.  Solving
$x=(4\gamma^2-1)/(4\gamma^2+1)$ for $\gamma>0$ gives \eqref{eq:gamma_from_x}.

Under the same unitary identification,
\begin{equation}
I-\mathcal J\sim\operatorname{diag}(1-x_\gamma),
\qquad
1-x_\gamma=\frac{2}{4\gamma^2+1}.
\label{eq:one_minus_xgamma}
\end{equation}
To prove trace-class membership without appealing to an unnamed estimate, enumerate distinct ordinates increasingly as $\gamma_j^*$.  Lemma~\ref{lem:distinct_ordinate_lower_bound} gives
$\gamma_j^*\gg j/\log j$, hence
\begin{equation}
\sum_j(1-x_{\gamma_j^*})
\ll\sum_{j\ge2}\frac{(\log j)^2}{j^2}<\infty.
\end{equation}
A positive diagonal operator is trace class exactly when the sum of its diagonal eigenvalues is finite.  Therefore
\begin{equation}
\operatorname{Tr}(I-\mathcal J)
=\sum_{\gamma\in\Gamma_+^*}\frac{2}{4\gamma^2+1},
\label{eq:Jacobi_trace_formula}
\end{equation}
and the same expression holds for both $E$ and $O$ because their supports coincide.

For every integer $k\ge1$, the positive operator $(I-\mathcal J)^k$ is trace class as well, since
$0\le1-x_\gamma<1$ and $\sum(1-x_\gamma)<\infty$.  Its eigenvalues are $(1-x_\gamma)^k$, so for both reflection families
\begin{equation}
\operatorname{Tr}\!\left((I-\mathcal J_E)^k\right)
=\operatorname{Tr}\!\left((I-\mathcal J_O)^k\right)
=\sum_{\gamma\in\Gamma_+^*}\left(\frac{2}{4\gamma^2+1}\right)^k.
\label{eq:Jacobi_trace_powers}
\end{equation}
Finally, trace invariance under change of orthonormal basis gives
\begin{equation}
\sum_{n\ge0}(1-b_n^E)
=\sum_{n\ge0}(1-b_n^O)
=\sum_{\gamma\in\Gamma_+^*}\frac{2}{4\gamma^2+1}.
\label{eq:Jacobi_diagonal_trace}
\end{equation}
Indeed, for either family the left side is $\sum_n\langle(I-\mathcal J)e_n,e_n\rangle=\operatorname{Tr}(I-\mathcal J)$; monotone convergence is automatic because $I-\mathcal J\ge0$.
\end{proof}

\begin{remark}[Finite truncations are not formulas for the zeros]
\label{rem:finite_jacobi_not_zeros}
The $N\times N$ leading Jacobi truncations have eigenvalues equal to the zeros of $\pi_N^E$ or $\pi_N^O$, equivalently Gaussian-quadrature nodes for the corresponding measure. These nodes are generally not points $x_\gamma$. Exact recovery of the support in \eqref{eq:jacobi_support} is an infinite inverse-moment statement requiring the full coefficient sequence. Thus \eqref{eq:gamma_from_x} converts an already recovered spectral support point into an ordinate; it is not a finite closed formula for a zeta zero.
\end{remark}

\begin{proposition}[Spectral masses and multiplicities]
\label{prop:spectral_multiplicity_weights}
For $\gamma\in\Gamma_+^*$, the unnormalized representing measures satisfy
\begin{align}
\mu_O(\{x_\gamma\})
&=\frac{4m_\gamma}{4\gamma^2+1},
\label{eq:muO_atom_mass}\\
\mu_E(\{x_\gamma\})
&=\frac{64m_\gamma\gamma^2}{(4\gamma^2+1)^2}.
\label{eq:muE_atom_mass}
\end{align}
Accordingly, the spectral measures of the first Jacobi basis vector are
\begin{align}
\nu_O(\{x_\gamma\})
&=\frac{4m_\gamma}{\lambda_1(4\gamma^2+1)},
\label{eq:nuO_atom_mass}\\
\nu_E(\{x_\gamma\})
&=\frac{64m_\gamma\gamma^2}{\lambda_2(4\gamma^2+1)^2}.
\label{eq:nuE_atom_mass}
\end{align}
Thus zero multiplicity is encoded by the distinguished cyclic spectral measure, even though the abstract multiplication operator has only one one-dimensional eigenspace for each distinct support point $x_\gamma$.
\end{proposition}

\begin{proof}
The first pair follows from \eqref{eq:mu_E}, \eqref{eq:mu_O}, and \eqref{eq:E_weight_gamma}--\eqref{eq:O_weight_gamma}. The second pair follows by the normalization \eqref{eq:normalized_spectral_measures}.
\end{proof}

\begin{corollary}[Trace and simplicity defect]
\label{cor:trace_simplicity_defect}
Under $\RH$,
\begin{equation}
\frac{\lambda_1}{2}
-
\operatorname{Tr}(I-\mathcal J_O)
=
2\sum_{\gamma\in\Gamma_+^*}
\frac{m_\gamma-1}{4\gamma^2+1}
\ge0.
\label{eq:trace_simplicity_defect}
\end{equation}
The same trace occurs for $\mathcal J_E$. Equality in \eqref{eq:trace_simplicity_defect} holds if and only if every positive ordinate is simple.
\end{corollary}

\begin{proof}
Equation \eqref{eq:li_cosine_representation} with $n=1$ and \eqref{eq:theta_cos} give
\begin{equation}
\frac{\lambda_1}{2}
=2\sum_{\gamma\in\Gamma_+^*}\frac{m_\gamma}{4\gamma^2+1}.
\label{eq:lambda1_multiplicity_sum}
\end{equation}
Subtract \eqref{eq:Jacobi_trace_formula}. Every summand in the resulting series is nonnegative, and it vanishes exactly when $m_\gamma=1$.
\end{proof}

\begin{proposition}[Operator bounds and recovery of the first distinct ordinate]
\label{prop:operator_bounds_first_zero}
Assume $\RH$ and let $\gamma_1$ be the smallest positive zero ordinate.  Then
\begin{equation}
x_{\gamma_1}I\le \mathcal J_E,\mathcal J_O\le I,
\label{eq:Jacobi_operator_bounds}
\end{equation}
and
\begin{equation}
\|I-\mathcal J_E\|
=\|I-\mathcal J_O\|
=\frac{2}{4\gamma_1^2+1}.
\label{eq:Jacobi_defect_norm}
\end{equation}
Consequently
\begin{equation}
\gamma_1
=\frac12\sqrt{\frac{2}{\|I-\mathcal J_E\|}-1}
=\frac12\sqrt{\frac{2}{\|I-\mathcal J_O\|}-1}.
\label{eq:first_zero_from_infinite_operator_norm}
\end{equation}
This is an identity for the complete infinite operator, not a finite formula for a zeta zero.
\end{proposition}

\begin{proof}
The spectrum in \eqref{eq:jacobi_support} is contained in $[x_{\gamma_1},1]$, which gives \eqref{eq:Jacobi_operator_bounds} by the spectral theorem.  The positive compact operator $I-\mathcal J$ has largest eigenvalue $1-x_{\gamma_1}=2/(4\gamma_1^2+1)$, proving \eqref{eq:Jacobi_defect_norm} and its inversion.
\end{proof}

\begin{theorem}[Singular values and the Schatten threshold]
\label{thm:Schatten_threshold}
Assume $\RH$, and enumerate the \emph{distinct} positive ordinates increasingly as
\begin{equation}
0<\gamma_1^*<\gamma_2^*<\cdots.
\label{eq:distinct_ordinate_enumeration}
\end{equation}
Then the singular values of either defect operator $K_E:=I-\mathcal J_E$ or $K_O:=I-\mathcal J_O$ are
\begin{equation}
s_k(K_E)=s_k(K_O)
=\frac{2}{4(\gamma_k^*)^2+1}.
\label{eq:singular_values_distinct_zeros}
\end{equation}
There exist positive constants $c,C$ such that, for all sufficiently large $k$,
\begin{equation}
c\frac{(\log\log k)^2}{k^2}
\le s_k(K_E)=s_k(K_O)
\le C\frac{(\log k)^2}{k^2}.
\label{eq:singular_value_two_sided_bounds}
\end{equation}
Consequently, for every $p>0$,
\begin{equation}
I-\mathcal J_E,\ I-\mathcal J_O\in\mathcal S_p
\quad\Longleftrightarrow\quad
p>\frac12.
\label{eq:Schatten_threshold}
\end{equation}
No asymptotic equivalent for \eqref{eq:singular_values_distinct_zeros} is asserted here, because the operator counts distinct ordinates rather than zeros with multiplicity.
\end{theorem}

\begin{proof}
Equation \eqref{eq:singular_values_distinct_zeros} follows from the diagonal multiplication representation and positivity of $K_E,K_O$.  Let $N(T)$ count zeros with multiplicity and $N^*(T)$ count distinct positive ordinates.  The Riemann--von Mangoldt formula gives $N(T)\asymp T\log T$.  Under $\RH$, Littlewood's classical bound
\begin{equation}
S(T)=O\!\left(\frac{\log T}{\log\log T}\right)
\label{eq:Littlewood_S_bound}
\end{equation}
implies the corresponding multiplicity bound
\begin{equation}
m(\tfrac12+i\gamma)
=O\!\left(\frac{\log\gamma}{\log\log\gamma}\right);
\label{eq:RH_multiplicity_bound}
\end{equation}
see, for example, \cite{Titchmarsh1986}.  Hence
\begin{equation}
T\log\log T\ll N^*(T)\ll T\log T.
\label{eq:distinct_zero_count_bounds}
\end{equation}
Inverting these bounds gives
\begin{equation}
\frac{k}{\log k}\ll\gamma_k^*\ll\frac{k}{\log\log k},
\label{eq:distinct_ordinate_bounds}
\end{equation}
which yields \eqref{eq:singular_value_two_sided_bounds}.  The upper bound is $p$-summable when $2p>1$; the lower bound forces divergence when $2p\le1$.  This proves \eqref{eq:Schatten_threshold}.
\end{proof}

\begin{proposition}[Endpoint exponent and critical Schatten exponent]
\label{prop:schatten_endpoint_exponent}
Let a positive discrete Gram system satisfy \eqref{eq:endpoint_power_law} with exponent $\beta>0$, and assume the two-sided node scale
\begin{equation}
\log t_j=(1+o(1))\log j.
\label{eq:two_sided_node_scale}
\end{equation}
Let $\mathcal J$ denote the Jacobi operator unitarily equivalent to multiplication by $x$ on $L^2(\mu)$, and put $K:=I-\mathcal J$. Since the support consists of the distinct atoms $x(t_j)$ with unique accumulation point $1$, the nonzero eigenvalues, hence the singular values, of the positive compact operator $K$ are the endpoint distances $1-x(t_j)$. Then the critical Schatten exponent of $K$ is
\begin{equation}
p_c(K):=\inf\{p>0:\ K\in\mathcal S_p\}=\frac1\beta.
\label{eq:critical_schatten_exponent}
\end{equation}
If in addition
\begin{equation}
\sum_j\frac1{t_j}=\infty,
\label{eq:endpoint_divergence_hypothesis}
\end{equation}
then membership at the critical exponent fails, and
\begin{equation}
K\in\mathcal S_p
\quad\Longleftrightarrow\quad
p>\frac1\beta.
\label{eq:schatten_threshold_general_beta}
\end{equation}
\end{proposition}

\begin{proof}
By \eqref{eq:endpoint_power_law},
\begin{equation}
c_1 t_j^{-\beta}\le 1-x(t_j)\le c_2 t_j^{-\beta}
\end{equation}
for all sufficiently large $j$. From \eqref{eq:two_sided_node_scale}, for every $\varepsilon>0$ there is $J_\varepsilon$ such that
\begin{equation}
j^{1-\varepsilon}\le t_j\le j^{1+\varepsilon},
\qquad j\ge J_\varepsilon.
\label{eq:two_sided_power_bounds}
\end{equation}
Therefore, for fixed $p>0$ and large $j$,
\begin{equation}
c_1^p j^{-\beta p(1+\varepsilon)}
\le
\bigl(1-x(t_j)\bigr)^p
\le
c_2^p j^{-\beta p(1-\varepsilon)}.
\label{eq:schatten_term_scale}
\end{equation}
If $p>1/\beta$, choose $\varepsilon>0$ so small that $\beta p(1-\varepsilon)>1$; the upper comparison in \eqref{eq:schatten_term_scale} gives convergence. If $p<1/\beta$, choose $\varepsilon>0$ so small that $\beta p(1+\varepsilon)<1$; the lower comparison gives divergence. This proves \eqref{eq:critical_schatten_exponent}.

Hypothesis \eqref{eq:two_sided_node_scale} alone does not decide the endpoint $p=1/\beta$. The admissible sequence $t_j=j(\log j)^2$ satisfies \eqref{eq:two_sided_node_scale} while $\sum_j t_j^{-1}<\infty$, so in that case $K\in\mathcal S_{1/\beta}$. Under the additional hypothesis \eqref{eq:endpoint_divergence_hypothesis} one has
\begin{equation}
\sum_j\bigl(1-x(t_j)\bigr)^{1/\beta}\asymp\sum_j\frac1{t_j}=\infty,
\end{equation}
so $K\notin\mathcal S_{1/\beta}$, and \eqref{eq:schatten_threshold_general_beta} follows.
\end{proof}

\begin{remark}[One geometric constant, two invariants]
\label{rem:beta_two_invariants}
Theorem~\ref{thm:endpoint_exponent_law} and Proposition~\ref{prop:schatten_endpoint_exponent} record the same quantity:
\begin{equation}
-\lim_{N\to\infty}\frac{\log D_N}{N^2\log N}
=\beta
=\frac{1}{p_c(I-\mathcal J)}.
\label{eq:determinant_schatten_unification}
\end{equation}
Once their respective coarse node hypotheses are satisfied, both invariants are fixed by the same endpoint exponent.  The determinant additionally uses the separated-subconfiguration and weight hypotheses of Theorem~\ref{thm:endpoint_exponent_law}, whereas the critical Schatten exponent uses the two-sided logarithmic node scale \eqref{eq:two_sided_node_scale}.  Neither invariant records the logarithmic Riemann--von Mangoldt refinement.  The determinant is a macroscopic quantity attached to the finite compressions, the critical exponent is an ideal-theoretic invariant of the infinite defect operator, and \eqref{eq:determinant_schatten_unification} states that, at this coarse level, they encode the same endpoint exponent.

For the Li--Cayley map \eqref{eq:universal_li_cayley_map} one has $\beta=2$, so the coefficient $-2$ of Theorem~\ref{thm:determinant_first_order} and the value $1/2$ of Theorem~\ref{thm:Schatten_threshold} are one geometric constant seen twice. The divergence hypothesis \eqref{eq:endpoint_divergence_hypothesis} holds in that case: by \eqref{eq:distinct_ordinate_bounds} one has $\gamma_k^*\ll k/\log\log k$, so $\sum_k(\gamma_k^*)^{-1}\gg\sum_k(\log\log k)/k=\infty$, and the strict form \eqref{eq:schatten_threshold_general_beta} is recovered.

In particular the value $1/2$ here is not the critical line. Replacing the endpoint map by one with exponent $\beta=3$ moves the critical Schatten exponent to $1/3$ while leaving the zeros untouched, and the value is likewise unrelated to abscissas of convergence of prime series, where a coincidental $1/2$ also occurs. This reinforces Remark~\ref{rem:what_is_universal}: the leading-order data of the pair is geometric, and arithmetic information can enter only below it.
\end{remark}

\begin{corollary}[Quadratic Jacobi sum rule]
\label{cor:quadratic_jacobi_sum_rule}
Under $\RH$,
\begin{equation}
\sum_{n\ge0}(1-b_n^E)^2+2\sum_{n\ge0}(a_{n+1}^E)^2
=
\sum_{\gamma\in\Gamma_+^*}\left(\frac{2}{4\gamma^2+1}\right)^2,
\label{eq:E_quadratic_sum_rule}
\end{equation}
and the identical formula holds with $E$ replaced by $O$.
\end{corollary}

\begin{proof}
For $K=I-\mathcal J$, the Hilbert--Schmidt norm squared is the sum of the squares of all matrix entries in the Jacobi basis and also the sum of the squares of its eigenvalues.  Since the off-diagonal entries occur twice, the stated identity follows.
\end{proof}

\begin{proposition}[Local Jacobi inequalities and summability]
\label{prop:local_jacobi_inequalities}
Under $\RH$, for either Jacobi family,
\begin{equation}
a_{n+1}^2\le(1-b_n)(1-b_{n+1})
\qquad(n\ge0).
\label{eq:local_jacobi_inequality}
\end{equation}
Moreover
\begin{equation}
\sum_{n\ge0}(1-b_n)<\infty,
\qquad
\sum_{n\ge0}a_{n+1}<\infty,
\label{eq:jacobi_l1_summability}
\end{equation}
and therefore
\begin{equation}
b_n\longrightarrow1,
\qquad
a_n\longrightarrow0.
\label{eq:jacobi_coeff_limits}
\end{equation}
\end{proposition}

\begin{proof}
The operator $I-\mathcal J$ is positive.  Positivity of each consecutive $2\times2$ principal minor gives \eqref{eq:local_jacobi_inequality}.  Equation \eqref{eq:Jacobi_diagonal_trace} gives the first series in \eqref{eq:jacobi_l1_summability}.  Since $1-b_n\ge0$,
\begin{equation}
a_{n+1}\le\sqrt{(1-b_n)(1-b_{n+1})}
\le\frac{(1-b_n)+(1-b_{n+1})}{2},
\label{eq:an_amgm_bound}
\end{equation}
which gives the second series and the limits.
\end{proof}

\begin{proposition}[Fredholm determinant of the Jacobi defect]
\label{prop:Fredholm_determinant}
Under $\RH$, the trace-class operators $K_E=I-\mathcal J_E$ and $K_O=I-\mathcal J_O$ have the same Fredholm determinant
\begin{equation}
D(z):=\det(I-zK_E)=\det(I-zK_O)
=\prod_{\gamma\in\Gamma_+^*}
\left(1-\frac{2z}{4\gamma^2+1}\right).
\label{eq:Fredholm_product}
\end{equation}
The product converges absolutely and locally uniformly, and the zeros of $D$ are
\begin{equation}
z_\gamma=\frac{4\gamma^2+1}{2},
\qquad \gamma\in\Gamma_+^*,
\label{eq:Fredholm_zeros}
\end{equation}
one for each distinct positive ordinate.  Zeta multiplicity is not encoded as the zero multiplicity of $D$.
\end{proposition}

\begin{proof}
By Theorem~\ref{thm:jacobi_zero_support}, each defect $K=I-\mathcal J$ is positive and trace class, and under the multiplication representation it is diagonal on the normalized atom indicators with eigenvalues
\begin{equation}
\delta_\gamma:=1-x_\gamma=\frac{2}{4\gamma^2+1}>0,
\qquad \gamma\in\Gamma_+^*.
\end{equation}
The only additional spectral point is $0$, coming from the accumulation $x_\gamma\to1$; it contributes no factor to the Fredholm determinant.  Since
\begin{equation}
\sum_{\gamma\in\Gamma_+^*}\delta_\gamma<\infty,
\end{equation}
the trace-class Fredholm determinant satisfies
\begin{equation}
\det(I-zK)=\prod_{\gamma\in\Gamma_+^*}(1-z\delta_\gamma).
\end{equation}
For completeness, local uniform convergence follows directly: on $|z|\le R$, choose a tail for which $R\delta_\gamma\le1/2$; then
$|\log(1-z\delta_\gamma)|\le2R\delta_\gamma$, and the Weierstrass test applies.  Thus the product defines an entire function and can vanish only when one factor vanishes, namely at
\begin{equation}
z=\delta_\gamma^{-1}=\frac{4\gamma^2+1}{2}.
\end{equation}
Each distinct support point supplies one one-dimensional eigenspace of the multiplication operator, so these zeros are indexed by distinct ordinates rather than by zeta multiplicity.  The two reflection families have the same support and hence the same eigenvalue list $\{\delta_\gamma\}$, proving equality of their Fredholm determinants.
\end{proof}

\begin{remark}[What the Jacobi operator does and does not retain]
\label{rem:operator_support_vs_multiplicity}
The abstract operator $\mathcal J$ determines the set of distinct transformed ordinates through its spectrum, but not the zeta multiplicities.  The normalized cyclic pair $(\mathcal J,e_0)$ determines the probability spectral measure and therefore the relative masses.  To recover the integer multiplicities from the odd realization one must also retain its scale $\lambda_1$, equivalently the unnormalized measure $\mu_O$.  Thus
\begin{equation}
(\mathcal J_O,e_0,\lambda_1)
\quad\text{determines}\quad
\{(\gamma,m_\gamma):\gamma\in\Gamma_+^*\}.
\label{eq:operator_cyclic_scale_data}
\end{equation}
The original semi-infinite matrix $\Omat$ already retains this scale because $o_{11}=\lambda_1$.
\end{remark}

\begin{proposition}[Unitary equivalence of the two Jacobi operators]
\label{prop:Jacobi_unitary_equivalence}
Under $\RH$, $\mathcal J_E$ and $\mathcal J_O$ are unitarily equivalent.
\end{proposition}

\begin{proof}
From \eqref{eq:christoffel_measure_relation},
\begin{equation}
\dd\mu_E(x)=2(1+x)\,\dd\mu_O(x).
\label{eq:measure_RN_repeat}
\end{equation}
The actual support is contained in $[x_{\gamma_1},1)$ with $x_{\gamma_1}>-1$, so $2(1+x)$ is bounded above and bounded away from zero on the support.  Define
\begin{equation}
(Uf)(x):=\sqrt{2(1+x)}\,f(x).
\label{eq:measure_unitary_map}
\end{equation}
Then for every $f\in L^2(\mu_E)$,
\begin{equation}
\|Uf\|_{L^2(\mu_O)}^2
=\int 2(1+x)|f(x)|^2\,\dd\mu_O(x)
=\int |f(x)|^2\,\dd\mu_E(x),
\end{equation}
so $U$ is an isometry.  Because the multiplier $[2(1+x)]^{-1/2}$ is bounded on the support, every $g\in L^2(\mu_O)$ has the preimage
$f=g/[2(1+x)]^{1/2}\in L^2(\mu_E)$; hence $U$ is onto and therefore unitary.

Let $M_x$ denote multiplication by $x$.  Since both $U$ and $M_x$ are multiplication operators,
$UM_x=M_xU$ on all of $L^2(\mu_E)$.  Let $\mathcal F_E:\ell^2(\mathbb N_0)\to L^2(\mu_E)$ and
$\mathcal F_O:\ell^2(\mathbb N_0)\to L^2(\mu_O)$ be the unitary maps sending the standard basis to the corresponding orthonormal polynomial bases.  By definition
\begin{equation}
\mathcal J_E=\mathcal F_E^{-1}M_x\mathcal F_E,
\qquad
\mathcal J_O=\mathcal F_O^{-1}M_x\mathcal F_O.
\end{equation}
Thus $V:=\mathcal F_O^{-1}U\mathcal F_E$ is unitary and
\begin{equation}
V\mathcal J_E
=\mathcal F_O^{-1}UM_x\mathcal F_E
=\mathcal F_O^{-1}M_xU\mathcal F_E
=\mathcal J_OV,
\end{equation}
which proves unitary equivalence.
\end{proof}

\begin{theorem}[Darboux factorization of the Jacobi pair]
\label{thm:Darboux_factorization}
Assume $\RH$.  Since $-1$ lies strictly below the common spectral support, $\mathcal J_O+I$ is strictly positive.  Let
\begin{equation}
\mathcal J_O+I=LL^*
\label{eq:JO_Cholesky_Darboux}
\end{equation}
be its positive lower-bidiagonal Cholesky factorization.  Then
\begin{equation}
\mathcal J_E+I=L^*L.
\label{eq:JE_Darboux_partner}
\end{equation}
Writing
\begin{equation}
L=
\begin{pmatrix}
r_0&&&\\
s_0&r_1&&\\
&s_1&r_2&\\
&&\ddots&\ddots
\end{pmatrix},
\qquad r_n>0,
\label{eq:Darboux_bidiagonal_L}
\end{equation}
one obtains
\begin{align}
1+b_n^O&=r_n^2+s_{n-1}^2,
& a_{n+1}^O&=r_ns_n,
\label{eq:Darboux_O_coefficients}\\
1+b_n^E&=r_n^2+s_n^2,
& a_{n+1}^E&=s_nr_{n+1},
\label{eq:Darboux_E_coefficients}
\end{align}
with $s_{-1}:=0$.
\end{theorem}

\begin{proof}
Let $\{p_n^E\}_{n\ge0}$ and $\{p_n^O\}_{n\ge0}$ be the orthonormal polynomial bases in $L^2(\mu_E)$ and $L^2(\mu_O)$, chosen with positive leading coefficients.  Consider the bounded operator
\begin{equation}
T:L^2(\mu_E)\longrightarrow L^2(\mu_O),
\qquad
(Tf)(x):=(x+1)f(x).
\label{eq:Darboux_T_operator}
\end{equation}
It is bounded because, using $\dd\mu_E=2(1+x)\dd\mu_O$ and $0<1+x\le2$ on the support,
\begin{align}
\|Tf\|_{L^2(\mu_O)}^2
&=\int (1+x)^2|f(x)|^2\,\dd\mu_O(x)\\
&=\frac12\int (1+x)|f(x)|^2\,\dd\mu_E(x)
\le\|f\|_{L^2(\mu_E)}^2.
\end{align}
For $f\in L^2(\mu_E)$ and $g\in L^2(\mu_O)$,
\begin{equation}
\langle Tf,g\rangle_{O}
=\int (x+1)f\overline g\,\dd\mu_O
=\frac12\int f\overline g\,\dd\mu_E.
\end{equation}
The identity embedding $g\mapsto g$ from $L^2(\mu_O)$ to $L^2(\mu_E)$ is bounded, since
\begin{equation}
\|g\|_{L^2(\mu_E)}^2
=2\int(1+x)|g|^2\,\dd\mu_O
\le4\|g\|_{L^2(\mu_O)}^2.
\end{equation}
Hence the adjoint is the everywhere-defined bounded operator
\begin{equation}
T^*g=\frac12g.
\label{eq:Darboux_T_adjoint}
\end{equation}
Consequently
\begin{equation}
T^*Tf=\frac12(x+1)f,
\qquad
TT^*g=\frac12(x+1)g.
\label{eq:Darboux_operator_products}
\end{equation}

Let $R$ be the matrix of $T$ from the $E$ orthonormal basis to the $O$ orthonormal basis,
\begin{equation}
R_{mn}:=\langle p_m^O,Tp_n^E\rangle_O.
\end{equation}
The polynomial $(x+1)p_n^E$ has degree $n+1$ and, by the same orthogonality calculation used in Proposition~\ref{prop:christoffel_polynomials}, is orthogonal in $L^2(\mu_O)$ to every polynomial of degree $<n$.  Therefore
\begin{equation}
Tp_n^E=\alpha_n p_n^O+\beta_n p_{n+1}^O,
\end{equation}
with no other orthonormal-polynomial components.  Comparison of positive leading coefficients gives \(\beta_n>0\).  The diagonal coefficient is also strictly positive:
\[
\alpha_n
=\langle p_n^O,(x+1)p_n^E\rangle_O
=\frac12\langle p_n^O,p_n^E\rangle_E.
\]
Writing \(p_n^O=(\kappa_n^O/\kappa_n^E)p_n^E+\) lower-degree terms, where
\(\kappa_n^{E,O}>0\) are the leading coefficients, and using the \(E\)-orthogonality
of \(p_n^E\) to lower-degree polynomials gives
\[
\alpha_n=\frac12\,\frac{\kappa_n^O}{\kappa_n^E}>0.
\]
Thus \(R\) is lower bidiagonal with positive diagonal and positive subdiagonal.  In coefficient space, \eqref{eq:Darboux_operator_products} becomes
\begin{equation}
R^*R=\frac12(\mathcal J_E+I),
\qquad
RR^*=\frac12(\mathcal J_O+I).
\end{equation}
Set
\begin{equation}
L:=\sqrt2\,R.
\end{equation}
Then $L$ is lower bidiagonal and
\begin{equation}
\mathcal J_O+I=LL^*,
\qquad
\mathcal J_E+I=L^*L.
\end{equation}
Since the common spectrum is contained in $[x_{\gamma_1},1]$ with $x_{\gamma_1}>-1$, both operators $\mathcal J_{E,O}+I$ are bounded below by $(1+x_{\gamma_1})I>0$.  Hence the bidiagonal factor has no zero diagonal entry; with the positive-leading-coefficient convention we take $r_n>0$.  This establishes the asserted Cholesky/Darboux factorization without an appeal to a formal infinite Gaussian elimination.

Finally, write $L_{nn}=r_n$ and $L_{n+1,n}=s_n$, with $s_{-1}=0$.  Direct multiplication gives
\begin{align}
(LL^*)_{nn}&=r_n^2+s_{n-1}^2,
&(LL^*)_{n,n+1}&=r_ns_n,\\
(L^*L)_{nn}&=r_n^2+s_n^2,
&(L^*L)_{n,n+1}&=s_nr_{n+1}.
\end{align}
Comparison with the Jacobi entries proves \eqref{eq:Darboux_O_coefficients}--\eqref{eq:Darboux_E_coefficients}.
\end{proof}

\begin{proposition}[Exact relation between the Weyl functions]
\label{prop:Weyl_function_relation}
Define, for $z$ off the common spectrum,
\begin{equation}
m_E(z):=\int\frac{\dd\nu_E(x)}{x-z},
\qquad
m_O(z):=\int\frac{\dd\nu_O(x)}{x-z}.
\label{eq:Weyl_functions_definition}
\end{equation}
Then
\begin{equation}
m_E(z)
=\frac{2\lambda_1}{\lambda_2}
\left[1+(z+1)m_O(z)\right],
\label{eq:Weyl_E_from_O}
\end{equation}
and, for $z\ne-1$, equivalently
\begin{equation}
m_O(z)
=\frac{\dfrac{\lambda_2}{2\lambda_1}m_E(z)-1}{z+1}.
\label{eq:Weyl_O_from_E}
\end{equation}
\end{proposition}

\begin{proof}
From \eqref{eq:christoffel_measure_relation} and \eqref{eq:normalized_spectral_measures},
\begin{equation}
\dd\nu_E(x)
=\frac{2\lambda_1}{\lambda_2}(1+x)\dd\nu_O(x).
\label{eq:normalized_christoffel_relation}
\end{equation}
Use
\begin{equation}
\frac{1+x}{x-z}=1+\frac{z+1}{x-z}
\label{eq:Weyl_fraction_identity}
\end{equation}
and $\nu_O([-1,1])=1$.
\end{proof}

\begin{proposition}[Christoffel relation between orthogonal systems]
\label{prop:christoffel_polynomials}
Under $\RH$,
\begin{equation}
\pi_n^E(x)=
\frac{
\pi_{n+1}^O(x)-
\dfrac{\pi_{n+1}^O(-1)}{\pi_n^O(-1)}\pi_n^O(x)
}{x+1}.
\label{eq:christoffel_polynomial_formula}
\end{equation}
The denominator $\pi_n^O(-1)$ never vanishes.
\end{proposition}

\begin{proof}
The scalar factor $2$ in $\dd\mu_E=2(1+x)\dd\mu_O$ does not affect orthogonality.  Let $q$ be any polynomial with $\deg q<n$.  Since $\pi_n^E$ is orthogonal to every such $q$ in $L^2(\mu_E)$,
\begin{equation}
0=\int \pi_n^E(x)q(x)\,\dd\mu_E(x)
=2\int (x+1)\pi_n^E(x)q(x)\,\dd\mu_O(x).
\end{equation}
Thus the monic polynomial $(x+1)\pi_n^E(x)$ has degree $n+1$ and is orthogonal, with respect to $\mu_O$, to all polynomials of degree at most $n-1$.

Because the monic orthogonal polynomials $\pi_0^O,\ldots,\pi_{n+1}^O$ form a basis of the polynomials of degree at most $n+1$, write
\begin{equation}
(x+1)\pi_n^E(x)=\sum_{j=0}^{n+1}c_j\pi_j^O(x).
\end{equation}
Taking the $L^2(\mu_O)$ inner product with $\pi_j^O$ for $j<n$ shows $c_j=0$.  Comparing leading coefficients gives $c_{n+1}=1$, so
\begin{equation}
(x+1)\pi_n^E(x)=\pi_{n+1}^O(x)+c_n\pi_n^O(x).
\end{equation}
Every zero of $\pi_n^O$ lies in the open convex hull of the support of the positive measure $\mu_O$; the support is contained in $( -1,1)$ and has left endpoint $x_{\gamma_1}>-1$.  Hence $\pi_n^O(-1)\ne0$.  Evaluating the preceding identity at $x=-1$ gives
$c_n=-\pi_{n+1}^O(-1)/\pi_n^O(-1)$, which is exactly \eqref{eq:christoffel_polynomial_formula} after division by $x+1$.  The numerator vanishes at $x=-1$, so the quotient is indeed a polynomial of degree $n$.
\end{proof}

\begin{corollary}[Determinant bridge]
\label{cor:determinant_bridge}
Under $\RH$,
\begin{equation}
\det\E_N
=2^N(-1)^N\pi_N^O(-1)\det\Omat_N.
\label{eq:determinant_bridge}
\end{equation}
\end{corollary}

\begin{proof}
We first derive the required Christoffel determinant identity from the moment determinants.  Let $\mu_k:=\int x^k\,\dd\mu_O(x)$ and
\begin{equation}
D_N[\mu_O]:=\det(\mu_{i+j})_{i,j=0}^{N-1}>0.
\end{equation}
The monic orthogonal polynomial of degree $N$ has the determinantal representation
\begin{equation}
\pi_N^O(c)
=\frac{1}{D_N[\mu_O]}
\det\begin{pmatrix}
\mu_0&\mu_1&\cdots&\mu_N\\
\mu_1&\mu_2&\cdots&\mu_{N+1}\\
\vdots&\vdots&&\vdots\\
\mu_{N-1}&\mu_N&\cdots&\mu_{2N-1}\\
1&c&\cdots&c^N
\end{pmatrix}.
\label{eq:monic_OP_determinant_formula}
\end{equation}
Indeed, expansion along the last row shows that the right side is monic of degree $N$; integrating it against $c^k\dd\mu_O(c)$ for $0\le k<N$ replaces the last row by the already present moment row $(\mu_k,\ldots,\mu_{k+N})$, so the determinant vanishes.  Thus it is exactly the monic orthogonal polynomial.

Starting from the determinant in \eqref{eq:monic_OP_determinant_formula}, perform for $j=N-1,N-2,\ldots,0$ the column operation
\begin{equation}
C_{j+1}\leftarrow C_{j+1}-cC_j.
\end{equation}
The determinant is unchanged.  The last row becomes $(1,0,\ldots,0)$, while after deleting its first column the remaining $N\times N$ minor has entries
$\mu_{i+j+1}-c\mu_{i+j}$, $0\le i,j<N$.  Expansion along the last row therefore gives
\begin{equation}
D_N[\mu_O]\,\pi_N^O(c)
=(-1)^N\det(\mu_{i+j+1}-c\mu_{i+j})_{i,j=0}^{N-1}.
\end{equation}
But the matrix on the right is precisely the moment matrix of the measure $(x-c)\dd\mu_O(x)$.  Hence
\begin{equation}
D_N[(x-c)\mu_O]
=(-1)^N\pi_N^O(c)D_N[\mu_O].
\label{eq:heine_christoffel_general}
\end{equation}
Taking $c=-1$ proves
\begin{equation}
D_N[(x+1)\mu_O]
=(-1)^N\pi_N^O(-1)D_N[\mu_O].
\label{eq:heine_christoffel_det}
\end{equation}

Now $\mu_E=2(x+1)\mu_O$.  Multiplying a measure by the scalar $2$ multiplies every entry of its $N\times N$ moment matrix by $2$, and therefore its determinant by $2^N$.  Finally, Theorem~\ref{thm:hankel_congruence} gives
\begin{equation}
\det\E_N=2^{N(N-1)}D_N[\mu_E],
\qquad
\det\Omat_N=2^{N(N-1)}D_N[\mu_O].
\end{equation}
The identical basis factors cancel in the ratio.  Combining these identities with \eqref{eq:heine_christoffel_det} yields
\begin{equation}
\det\E_N
=2^N(-1)^N\pi_N^O(-1)\det\Omat_N,
\end{equation}
as claimed.
\end{proof}

\begin{corollary}[Bridge between the elimination pivots]
\label{cor:pivot_bridge}
Under $\RH$,
\begin{equation}
d_n^E
=-2\frac{\pi_n^O(-1)}{\pi_{n-1}^O(-1)}d_n^O
\qquad(n\ge1),
\label{eq:pivot_bridge}
\end{equation}
where $d_n^E,d_n^O$ are defined in \eqref{eq:pivot_short_notation}.
\end{corollary}

\begin{proof}
Divide \eqref{eq:determinant_bridge} for $N=n$ by the same identity for $N=n-1$.
\end{proof}


\section{The Christoffel margin: exact factorisation of the pivot ratio and a two-sided band}
\label{sec:christoffel_bridge}

Corollary~\ref{cor:determinant_bridge} links the two determinant sequences through the
single number $\pi_N^O(-1)$.  This section identifies that number as a truncated Fredholm
determinant of the Jacobi defect, deduces an exact factorisation of the ratio of the two
pivot sequences into a universal constant times an arithmetic factor, and turns the
resulting monotonicity into a two-sided band which is a necessary condition for $\RH$
requiring no arithmetic input at all.

\begin{definition}[Christoffel factor]
\label{def:christoffel_factor}
For $N\ge1$ set
\begin{equation}
 F_N:=\frac{(-1)^N\pi_N^O(-1)}{2^N},
 \qquad F_0:=1 .
 \label{eq:FN_definition}
\end{equation}
Corollary~\ref{cor:determinant_bridge} then reads
\begin{equation}
 \boxed{\;\Delta_N^E=4^N F_N\,\Delta_N^O\;}
 \label{eq:determinant_bridge_normalized}
\end{equation}
\end{definition}

\begin{proposition}[$F_N$ is a truncated Fredholm determinant of the Jacobi defect]
\label{prop:FN_spectral}
Under $\RH$, let $\mathcal J_O^{(N)}$ be the leading $N\times N$ truncation of the odd Jacobi
operator and let $\mathcal D_N:=I_N-\mathcal J_O^{(N)}$ be its defect (the symbol $\mathcal K$ is reserved for the Christoffel--Darboux kernels of Section~\ref{sec:inverse}).  Then
\begin{equation}
 \boxed{\;
 F_N=\det\!\left(\frac{I_N+\mathcal J_O^{(N)}}{2}\right)
 =\det\!\left(I_N-\tfrac12\mathcal D_N\right)
 =\sum_{k=0}^{N}\frac{(-1)^k}{2^k}\,e_k\bigl(\mathcal D_N\bigr),\;}
 \label{eq:FN_spectral}
\end{equation}
where $e_k$ denotes the $k$-th elementary symmetric function of the eigenvalues.
\end{proposition}

\begin{proof}
The monic orthogonal polynomial of degree $N$ is the characteristic polynomial of the
leading Jacobi truncation, $\pi_N^O(x)=\det\bigl(xI_N-\mathcal J_O^{(N)}\bigr)$; this is the
statement recalled in Remark~\ref{rem:finite_jacobi_not_zeros}.  Evaluating at $x=-1$,
\begin{equation}
 \pi_N^O(-1)=\det\bigl(-I_N-\mathcal J_O^{(N)}\bigr)=(-1)^N\det\bigl(I_N+\mathcal J_O^{(N)}\bigr),
\end{equation}
so $(-1)^N\pi_N^O(-1)=\det\bigl(I_N+\mathcal J_O^{(N)}\bigr)$ and \eqref{eq:FN_definition} gives the
first equality in \eqref{eq:FN_spectral}.  The second is $I+\J=2I-(I-\J)$ divided by $2$,
and the third is the expansion of $\det(I-zK)$ in elementary symmetric functions at
$z=1/2$.
\end{proof}

\begin{theorem}[Margin factorisation]
\label{thm:margin_factorisation}
Let $p_N^E:=\Delta_N^E/\Delta_{N-1}^E$ and $p_N^O:=\Delta_N^O/\Delta_{N-1}^O$ be the
Cholesky pivots \eqref{eq:E_LDL_pivots}--\eqref{eq:O_LDL_pivots}, and define the
\emph{margin}
\begin{equation}
 \kappa_N:=\frac{p_N^E}{p_N^O}.
 \label{eq:kappa_definition}
\end{equation}
Then, under $\RH$,
\begin{equation}
 \boxed{\;\kappa_N=4\,\frac{F_N}{F_{N-1}}\;}
 \qquad(N\ge1).
 \label{eq:margin_factorisation}
\end{equation}
Within any positive even--odd pair with the same Christoffel relation
\(\dd\mu_E=2(1+x)\dd\mu_O\) and the same polynomial normalizations used here, the
constant \(4\) is independent of the particular measure.  The measure-dependent
information in this factorization sits in \(F_N/F_{N-1}\).
\end{theorem}

\begin{proof}
Divide \eqref{eq:determinant_bridge_normalized} at $N$ by the same identity at $N-1$.
\end{proof}

\begin{theorem}[Monotonicity, band, and Fredholm limit]
\label{thm:margin_band}
Let $\mathcal D=I-\mathcal J_O$ be the Jacobi defect of Proposition~\ref{prop:local_jacobi_inequalities},
a positive trace-class operator satisfying $\Tr\mathcal D\le\lambda_1/2$, with equality if and only if every positive zero ordinate is simple.  Under $\RH$:
\begin{enumerate}[label=(\roman*)]
\item $F_N$ is nonincreasing in $N$ and $0<F_N\le1$;
\item $\det\bigl(I-\tfrac12\mathcal D\bigr)\le F_N\le1$ for every $N$, and
      $F_N\to\det\bigl(I-\tfrac12\mathcal D\bigr)$ as $N\to\infty$, where
      \begin{equation}
       \det\Bigl(I-\tfrac12\mathcal D\Bigr)
       =\prod_{\gamma\in\Gamma_+^*}\frac{4\gamma^2}{4\gamma^2+1};
       \label{eq:Fredholm_at_half}
      \end{equation}
\item consequently
      \begin{equation}
       \boxed{\;
       4\det\Bigl(I-\tfrac12\mathcal D\Bigr)\;\le\;\kappa_N\;\le\;4
       \qquad(N\ge1),\;}
       \label{eq:margin_band}
      \end{equation}
      and the weaker but entirely explicit form
      \begin{equation}
       4-\lambda_1\;\le\;\kappa_N\;\le\;4 .
       \label{eq:margin_band_explicit}
      \end{equation}
\end{enumerate}
\end{theorem}

\begin{proof}
Write $A:=\tfrac12(I+\mathcal J_O)=I-\tfrac12\mathcal D$ and $A_N$ for its leading $N\times N$
truncation, so that $F_N=\det A_N$ by Proposition~\ref{prop:FN_spectral}.  By
Theorem~\ref{thm:jacobi_zero_support} the spectrum of $\mathcal J_O$ lies in
$[x_{\gamma_1},1]$, hence $0\prec A\preceq I$ and every diagonal entry of $A$ in the
Jacobi basis satisfies $0<A_{nn}\le1$.

(i) Positivity of $F_N$ is positivity of the principal minors of the positive definite
operator $A$.  For monotonicity, bordering gives
\begin{equation}
 \det A_N=\det A_{N-1}\bigl(A_{NN}-v^TA_{N-1}^{-1}v\bigr)\le \det A_{N-1}\,A_{NN}
 \le\det A_{N-1},
\end{equation}
where $v$ is the last off-diagonal block and $A_{N-1}^{-1}\succ0$.  Also $F_1=A_{11}\le1$.

(ii) Let $\mu_1\ge\mu_2\ge\cdots$ be the eigenvalues of $\mathcal D$ and
$\mu_1^{(N)}\ge\cdots\ge\mu_N^{(N)}$ those of the compression
$\mathcal D_N=P_N\mathcal D P_N$.  Cauchy interlacing for compressions gives
$\mu_j^{(N)}\le\mu_j$ for $1\le j\le N$, and $0\le\mu_j<1$ for all $j$ because
$\|\mathcal D\|=1-x_{\gamma_1}<1$ by \eqref{eq:Jacobi_defect_norm}.  Hence
\begin{equation}
 F_N=\prod_{j=1}^{N}\Bigl(1-\tfrac12\mu_j^{(N)}\Bigr)
 \ge\prod_{j=1}^{N}\Bigl(1-\tfrac12\mu_j\Bigr)
 \ge\prod_{j\ge1}\Bigl(1-\tfrac12\mu_j\Bigr)
 =\det\Bigl(I-\tfrac12\mathcal D\Bigr),
\end{equation}
each factor being in $(0,1]$.  Since $\mathcal D$ is trace class, $\mathcal D_N\to\mathcal D$ in trace
norm and the Fredholm determinant is trace-norm continuous, so $F_N\to\det(I-\mathcal D/2)$.
Formula \eqref{eq:Fredholm_at_half} is $D(1/2)$ in the notation of
Proposition~\ref{prop:Fredholm_determinant}, since
$1-\tfrac{1}{4\gamma^2+1}=\tfrac{4\gamma^2}{4\gamma^2+1}$.

(iii) Combining (i), (ii) with \eqref{eq:margin_factorisation} gives
$\kappa_N=4F_N/F_{N-1}\le4$ and $\kappa_N\ge4\det(I-\mathcal D/2)$, because
$F_N\ge\det(I-\mathcal D/2)$ and $F_{N-1}\le1$.  Finally, for a positive operator with
eigenvalues in $[0,1)$ one has $\prod(1-\mu_j/2)\ge1-\tfrac12\sum\mu_j$, and
$\Tr\mathcal D\le\lambda_1/2$ by Corollary~\ref{cor:trace_simplicity_defect}, whence
$\det(I-\mathcal D/2)\ge1-\lambda_1/4$ and $4\det(I-\mathcal D/2)\ge4-\lambda_1$.
\end{proof}

\begin{corollary}[An arithmetic-free necessary condition]
\label{cor:kappa_le_4}
If $\RH$ holds, then for every $N\ge1$,
\[
\Delta_N^E\,\Delta_{N-1}^O
\le4\,\Delta_{N-1}^E\,\Delta_N^O.
\]
Equivalently $F_N\le F_{N-1}$.  No zero, prime, or Li coefficient beyond those defining the two
determinants enters this inequality.
\end{corollary}

\begin{proof}
This is the upper half of \eqref{eq:margin_band}, whose only input is $\mathcal J_O\preceq I$,
that is, that the odd reflection measure is supported in $(-\infty,1]$.  Under $\RH$ the
support is $\{x_\gamma\}\cup\{1\}\subset(0,1]$.
\end{proof}

\begin{remark}[What is universal and what is arithmetic]
\label{rem:kappa_universal_vs_arithmetic}
Identity \eqref{eq:margin_factorisation} is an exact algebraic statement about \emph{any}
positive sine Gram pair, not a small-angle approximation: it was verified to working
precision for angles drawn uniformly from $(0,0.08)$, where the first tested values are
\[
\kappa_N=3.9959\ldots,\ 3.9967\ldots,\ 3.9970\ldots .
\]
For angles drawn uniformly from $(0,\pi)$, the same identity holds exactly; for $N=1,\ldots,6$ the
tested values are
\[
1.3484,\ 1.3844,\ 1.1409,\ 1.4586,\ 1.2812,\ 1.4200,
\]
and $F_N$ falls to $1.4\times10^{-3}$.  What requires the zeta geometry is not the factorisation but the \emph{value}.  In the finite
zeta computations reported below, $\kappa_N$ is close to $4$, consistently with the small Cayley
angles and with $F_N/F_{N-1}$ being close to $1$.  The universal object is the constant $4$
together with the exact factorisation, while the arithmetic object is $F_N$.
\end{remark}



\begin{remark}[Trace normalization]
\label{rem:trace_normalisation}
Writing $q_n:=2\Tr\mathcal J_O^{(n)}$, the identity $2n-q_n=2\Tr\mathcal D_n$ is immediate from
$\mathcal D_n=I_n-\mathcal J_O^{(n)}$, and by \eqref{eq:Jacobi_diagonal_trace} the increasing sequence
$\Tr\mathcal D_n=\sum_{j<n}(1-b_j^O)$ converges to $\Tr\mathcal D\le\lambda_1/2$, with equality exactly in the simple-zero case.  This is the
first-order shadow of Theorem~\ref{thm:margin_band}(ii): $\log F_n$ decreases to
$\log\det(I-\mathcal D/2)$, whose linearization is $-\tfrac12\Tr\mathcal D=-\lambda_1/4$.
\end{remark}

\begin{experiment}[The Christoffel margin for the arithmetic $\lambda$]
\label{exp:christoffel_margin}
All entries below use the arithmetic $\lambda$ of Experiment~\ref{exp:arithmetic_lambda}
at $250$ working digits.  The columns $F_N$ and $\det\bigl((I+\mathcal J_O^{(N)})/2\bigr)$ were
computed by two independent routes: the first from
$\Delta_N^E/(4^N\Delta_N^O)$, the second from the Jacobi matrix reconstructed from the
moments of $\mu_O$ by Cholesky.
\begin{center}
\small
\begin{tabular}{rlllr}
\toprule
$N$ & $F_N$ & $\kappa_N=p_N^E/p_N^O$ & $|\kappa_N-4F_N/F_{N-1}|$ & $\Tr\mathcal D_N$\\
\midrule
$1$  & $0.99959840336042631$ & $3.9983936134417052$ & $0$                  & $0.000803193$\\
$2$  & $0.99872447452420145$ & $3.9965028802235502$ & $7.0\times10^{-251}$ & $0.002551298$\\
$3$  & $0.99827559875164635$ & $3.9982022037748942$ & $0$                  & $0.003450031$\\
$4$  & $0.99797192293554315$ & $3.9987831984815297$ & $3.5\times10^{-251}$ & $0.004058371$\\
$5$  & $0.99770423468019879$ & $3.9989270709959176$ & $3.5\times10^{-251}$ & $0.004594797$\\
$6$  & $0.99748822571541767$ & $3.9991339759529021$ & $0$                  & $0.005027778$\\
$7$  & $0.99733892798897730$ & $3.9994013053083076$ & $7.0\times10^{-251}$ & $0.005327110$\\
$8$  & $0.99720226204496688$ & $3.9994518776308632$ & $0$                  & $0.005601162$\\
$9$  & $0.99706026967019839$ & $3.9994304370129394$ & $3.5\times10^{-251}$ & $0.005885933$\\
$10$ & $0.99694804248538865$ & $3.9995497676992109$ & $3.5\times10^{-251}$ & $0.006111041$\\
\bottomrule
\end{tabular}
\end{center}
The two routes to $F_N$ agree to $1.1\times10^{-182}$ at $N=10$, and
\eqref{eq:margin_factorisation} holds to the last working digit; the entries in the fourth
column are pure rounding.  The band \eqref{eq:margin_band} is satisfied with
$4\det(I-\mathcal D/2)=3.976966$ from the first $1500$ ordinates with the
Riemann--von Mangoldt tail, against the explicit bound $4-\lambda_1=3.976904$.  The
Fredholm value itself is
\begin{equation}
 \det\Bigl(I-\tfrac12\mathcal D\Bigr)
 =\prod_{\gamma>0}\frac{4\gamma^2}{4\gamma^2+1}
 =0.9942415\ldots,
 \label{eq:Fredholm_numeric}
\end{equation}
whose truncation to the first $1200$ ordinates is $0.9943992$ and to the first $1500$
ordinates $0.9943764$; the difference is the tail
$\sum_{\gamma>T}(4\gamma^2+1)^{-1}=1.36\times10^{-4}$.  Independently, a first-$1500$-ordinate evaluation supplemented by the same
Riemann--von Mangoldt tail approximation gives
$\sum 2/(4\gamma^2+1)=0.0115479184$, within $6.4\times10^{-8}$ of the rigorous arithmetic
upper bound $\lambda_1/2=0.0115478545$.  Since Corollary~\ref{cor:trace_simplicity_defect}
proves $\Tr\mathcal D\le\lambda_1/2$, the small numerical overshoot measures the tail
approximation error and is not interpreted as an excess of the exact trace.  Convergence
of $F_N$ to \eqref{eq:Fredholm_numeric} is slow: at $N=10$ only about a fifth of the
estimated total defect has been accumulated.
\end{experiment}

\begin{experiment}[The binding component of the local Jacobi inequalities]
\label{exp:binding_component}
Proposition~\ref{prop:local_jacobi_inequalities} gives
$a_{n+1}^2\le(1-b_n)(1-b_{n+1})$ for every $n\ge0$.  Define the normalized margin
\begin{equation}
 \varrho_n:=\frac{a_{n+1}^2}{(1-b_n)(1-b_{n+1})}\in[0,1],
 \label{eq:local_jacobi_margin}
\end{equation}
so that $\varrho_n$ close to $1$ means the corresponding $2\times2$ minor is close to
degenerate.  With the arithmetic $\lambda$ one finds, for both families,
\begin{center}
\small
\begin{tabular}{lcccccccc}
\toprule
$n$ & $0$ & $1$ & $2$ & $3$ & $4$ & $5$ & $6$ & $7$\\
\midrule
$\varrho_n$ (even) & $0.6482$ & $0.2100$ & $0.2222$ & $0.2353$ & $0.2665$ & $0.2401$ & $0.2258$ & $0.2633$\\
$\varrho_n$ (odd)  & $0.6482$ & $0.2099$ & $0.2222$ & $0.2353$ & $0.2665$ & $0.2401$ & $0.2258$ & $0.2633$\\
\bottomrule
\end{tabular}
\end{center}
The binding index is $n=0$ in every positive definite case tested, with
$\varrho_0=0.648160$ (even) and $0.648207$ (odd), a factor of about three above all later
values.  This is the precise sense in which ``the first component binds''.  The two
families agree to four digits because they differ by a Christoffel transform at a point
outside the support, which perturbs the low-index recurrence coefficients only at order
$1-x_{\gamma_1}$.
\end{experiment}

\section{Hermitian congruence reduction and inertia}
\label{sec:hermitian_reduction}

For the real symmetric matrices considered here, the intrinsic analogue of a reduction of the quadratic form is congruence diagonalization. From \eqref{eq:LDL_decompositions},
\begin{align}
L_E^{-1}\E_NL_E^{-T}
&=D_E
=\diag\left(
\frac{\Delta_1^E}{\Delta_0^E},
\frac{\Delta_2^E}{\Delta_1^E},\ldots,
\frac{\Delta_N^E}{\Delta_{N-1}^E}
\right),
\label{eq:E_hermitian_reduction}\\
L_O^{-1}\Omat_NL_O^{-T}
&=D_O
=\diag\left(
\frac{\Delta_1^O}{\Delta_0^O},
\frac{\Delta_2^O}{\Delta_1^O},\ldots,
\frac{\Delta_N^O}{\Delta_{N-1}^O}
\right).
\label{eq:O_hermitian_reduction}
\end{align}
Thus the signature is read from the signs of the pivot quotients whenever no principal determinant vanishes.

\begin{remark}[Hermite normal form]
\label{rem:HNF}
The integer Hermite normal form is not intrinsic to $\E_N$ or $\Omat_N$: the entries are real numbers built from Li coefficients, and no natural embedding of these matrices into a specified principal ideal domain is part of the construction. No transcendence property of the Li coefficients is assumed or known here. Hermite normal form is an arithmetic reduction after a coefficient ring such as $\mathbb Z$ has been fixed, as in \cite{KannanBachem1979}. For the present real symmetric forms, the intrinsic reduction is the congruence diagonalization in \eqref{eq:E_hermitian_reduction}--\eqref{eq:O_hermitian_reduction}.
\end{remark}

\section{Geometric transformations}
\label{sec:geometry}

Under $\RH$, each matrix defines an inner product on $\mathbb R^N$:
\begin{equation}
\langle c,d\rangle_{E}:=c^T\E_Nd,
\qquad
\langle c,d\rangle_{O}:=c^T\Omat_Nd.
\label{eq:matrix_inner_products}
\end{equation}
The associated synthesis maps are
\begin{align}
(\mathcal S_Ec)_\gamma
&:=2\sum_{n=1}^{N}c_n\sin(n\theta_\gamma),
\label{eq:SE_synthesis}\\
(\mathcal S_Oc)_\gamma
&:=2\sum_{n=1}^{N}c_n\sin((n-\tfrac12)\theta_\gamma),
\label{eq:SO_synthesis}
\end{align}
so that
\begin{equation}
\E_N=\mathcal S_E^*\mathcal S_E,
\qquad
\Omat_N=\mathcal S_O^*\mathcal S_O.
\label{eq:synthesis_gram}
\end{equation}
Hence the matrices are pullback metrics from the zero-side $\ell^2$ geometry.

If $\sigma_1\ge\cdots\ge\sigma_N>0$ are the eigenvalues of $\E_N$, then
\begin{align}
\sigma_1
&=4\max_{\|c\|_2=1}
\sum_{\gamma>0}
\left|\sum_{n=1}^{N}c_n\sin(n\theta_\gamma)\right|^2,
\label{eq:E_largest_eigenvalue_variational}\\
\sigma_N
&=4\min_{\|c\|_2=1}
\sum_{\gamma>0}
\left|\sum_{n=1}^{N}c_n\sin(n\theta_\gamma)\right|^2.
\label{eq:E_smallest_eigenvalue_variational}
\end{align}
Thus $\sigma_N$ measures the least zero-side energy of a normalized sine polynomial of degree $N$. Analogous formulas hold for $\Omat_N$ with half-integer modes.

The positive square root $\E_N^{1/2}$ maps the Euclidean unit ball to the ellipsoid whose principal semiaxes have lengths $\sqrt{\sigma_j}$. Cholesky gives a triangular coordinate realization of the same geometry, while $L_E^{-1}$ gives the Gram--Schmidt coordinates in which the form becomes diagonal. The Chebyshev changes $B_E^{(N)}$ and $B_O^{(N)}$ convert from reflection-adapted harmonic coordinates to monomial moment coordinates. Passing to the infinite orthonormal-polynomial bases yields the Jacobi operators $\mathcal J_E$ and $\mathcal J_O$, in which multiplication by $x$ is tridiagonal.

\begin{corollary}[Eigenvalue interlacing]
\label{cor:eigenvalue_interlacing}
Under $\RH$, the eigenvalues of $\E_N$ interlace those of $\E_{N+1}$, and the same holds for $\Omat_N$ and $\Omat_{N+1}$.
\end{corollary}

\begin{proof}
Each $N\times N$ matrix is a leading principal submatrix of the next truncation, so Cauchy's interlacing theorem applies.
\end{proof}

\section{Generating functions}
\label{sec:generating_functions}

The complete infinite arrays have compact bivariate generating functions. Let $\Lambda(z)$ be given by \eqref{eq:li_generating_function}.
\begin{theorem}[Bivariate generating functions]
\label{thm:bivariate_generating_functions}
As identities of formal power series, one has
\begin{align}
\sum_{m,n\ge1}e_{mn}x^my^n
&=
\frac{y\Lambda(x)-x\Lambda(y)}{x-y}
-
\frac{xy}{1-xy}\bigl(\Lambda(x)+\Lambda(y)\bigr),
\label{eq:E_bivariate_generating_function}\\
\sum_{m,n\ge1}o_{mn}x^my^n
&=
\frac{xy\bigl(\Lambda(x)-\Lambda(y)\bigr)}{x-y}
-
\frac{xy}{1-xy}\bigl(\Lambda(x)+\Lambda(y)\bigr).
\label{eq:O_bivariate_generating_function}
\end{align}
Analytically, these formulas hold whenever $|x|,|y|<R_\xi$; under $\RH$ this is the full bidisk $|x|,|y|<1$.
\end{theorem}

\begin{proof}
For the Hankel part,
\begin{equation}
\sum_{m,n\ge1}\lambda_{m+n}x^my^n
=
\frac{y\Lambda(x)-x\Lambda(y)}{x-y},
\label{eq:Hankel0_generating}
\end{equation}
and
\begin{equation}
\sum_{m,n\ge1}\lambda_{m+n-1}x^my^n
=
\frac{xy(\Lambda(x)-\Lambda(y))}{x-y}.
\label{eq:Hankelminus1_generating}
\end{equation}
For the Toeplitz term,
\begin{equation}
\sum_{m,n\ge1}\lambda_{|m-n|}x^my^n
=
\frac{xy}{1-xy}\bigl(\Lambda(x)+\Lambda(y)\bigr),
\label{eq:Toeplitz_generating}
\end{equation}
because $\lambda_0=0$. Subtraction yields the result.
\end{proof}

Substituting \eqref{eq:li_generating_function} gives bivariate expressions directly in $\xi'/\xi$ evaluated at the Cayley variable $1/(1-z)$.

\section{Finite-dimensional inequalities}
\label{sec:inequalities}

Positive semidefiniteness generates a hierarchy of inequalities among Li coefficients.

\begin{corollary}[Two-by-two inequalities]
\label{cor:two_by_two_inequalities}
Under $\RH$, for all $m,n\ge1$,
\begin{align}
\bigl(\lambda_{m+n}-\lambda_{|m-n|}\bigr)^2
&\le \lambda_{2m}\lambda_{2n},
\label{eq:E_2x2_inequality}\\
\bigl(\lambda_{m+n-1}-\lambda_{|m-n|}\bigr)^2
&\le \lambda_{2m-1}\lambda_{2n-1}.
\label{eq:O_2x2_inequality}
\end{align}
\end{corollary}

\begin{proof}
Apply nonnegativity of every $2\times2$ principal determinant.
\end{proof}

Special cases include
\begin{align}
(\lambda_3-\lambda_1)^2&\le\lambda_2\lambda_4,
\label{eq:first_even_inequality}\\
(\lambda_2-\lambda_1)^2&\le\lambda_1\lambda_3,
\label{eq:first_odd_inequality}\\
(\lambda_{n+1}-\lambda_{n-1})^2&\le\lambda_2\lambda_{2n},
\label{eq:even_difference_bound}\\
(\lambda_n-\lambda_{n-1})^2&\le\lambda_1\lambda_{2n-1}.
\label{eq:odd_difference_bound}
\end{align}

For three indices $a,b,c$, set
\begin{equation}
x_{ab}^E:=\lambda_{a+b}-\lambda_{|a-b|}.
\label{eq:xabE}
\end{equation}
Then the $3\times3$ principal determinant gives
\begin{equation}
\begin{aligned}
&\lambda_{2a}\lambda_{2b}\lambda_{2c}
+2x_{ab}^Ex_{ac}^Ex_{bc}^E
-\lambda_{2a}(x_{bc}^E)^2\\
&\qquad
-\lambda_{2b}(x_{ac}^E)^2
-\lambda_{2c}(x_{ab}^E)^2
\ge0.
\end{aligned}
\label{eq:E_3x3_inequality}
\end{equation}
The odd analogue follows after replacing $x_{ab}^E$ by
\begin{equation}
x_{ab}^O:=\lambda_{a+b-1}-\lambda_{|a-b|}
\label{eq:xabO}
\end{equation}
and the diagonal factors $\lambda_{2a}$ by $\lambda_{2a-1}$.

\section{Trace, norm, and conditioning data}
\label{sec:trace_conditioning}

The elementary spectral invariants have direct Li expressions:
\begin{align}
\operatorname{tr}\E_N
&=\sum_{n=1}^{N}\lambda_{2n},
\label{eq:E_trace}\\
\operatorname{tr}\Omat_N
&=\sum_{n=1}^{N}\lambda_{2n-1},
\label{eq:O_trace}\\
\operatorname{tr}\E_N+\operatorname{tr}\Omat_N
&=\sum_{k=1}^{2N}\lambda_k.
\label{eq:combined_trace}
\end{align}
The Frobenius norms are
\begin{align}
\|\E_N\|_F^2
&=\sum_{m,n=1}^{N}
\bigl(\lambda_{m+n}-\lambda_{|m-n|}\bigr)^2,
\label{eq:E_frobenius}\\
\|\Omat_N\|_F^2
&=\sum_{m,n=1}^{N}
\bigl(\lambda_{m+n-1}-\lambda_{|m-n|}\bigr)^2.
\label{eq:O_frobenius}
\end{align}
Under $\RH$, the condition numbers
\begin{equation}
\kappa_2(\E_N)=\frac{\sigma_{\max}(\E_N)}{\sigma_{\min}(\E_N)},
\qquad
\kappa_2(\Omat_N)=\frac{\sigma_{\max}(\Omat_N)}{\sigma_{\min}(\Omat_N)}
\label{eq:condition_numbers}
\end{equation}
measure the anisotropy of the integer and half-integer zero-side sampling geometries. Corollary~\ref{cor:condition_number_obstruction} gives the nonperturbative lower bound
\begin{equation}
\log\kappa_2(\E_N),\ \log\kappa_2(\Omat_N)
\ge 2N\log N-2N\log\log N-O(N),
\label{eq:conditioning_practical_bound}
\end{equation}
so ill-conditioning is forced asymptotically by the determinant law and is not merely an implementation issue. The smallest eigenvalue is particularly informative because it is the minimum energy in \eqref{eq:E_smallest_eigenvalue_variational} and its odd analogue.

\section{Literature position and comparison with other matrix criteria}
\label{sec:literature_position}

For avoidance of any ambiguity about origin, the input used to manufacture the reflection matrices is copied algebraically from the published Li--Weil framework: Li supplies the coefficients and criterion \cite{Li1997}; Bombieri--Lagarias supply the general multiset criterion, explicit-formula interpretation and Li test functions in inverse-Mellin coordinates \cite{BombieriLagarias1999}; and Lagarias supplies the exact bilateral Gram identity used here \cite[Theorem~3.1]{Lagarias2007}:
\begin{equation}
\langle G_n,G_m\rangle_W
=\lambda_n+\lambda_{-m}-\lambda_{n-m}.
\label{eq:literature_Lagarias_exact}
\end{equation}
Accordingly, neither the Li positivity mechanism nor the full Li--Weil Gram kernel is new here.
The specific origin of \(\E_N\) is the reflection map
\(G_n\mapsto (G_n-G_{-n})/\sqrt2\) applied to Lagarias's known bilateral kernel; the odd family is
obtained from \(G_n\mapsto(G_n-G_{1-n})/\sqrt2\).  This is the precise sense in which the reflection
matrices are transformations of known Li--Weil data rather than independent replacements for it.
Bombieri studied other finite-dimensional truncations of Weil's quadratic functional
\cite{Bombieri2000}.  Matiyasevich gave matrix/eigenvalue reformulations based on Taylor
coefficients and Toeplitz/Hankel rearrangements \cite{Matiyasevich2007a,Matiyasevich2007b}.
Suzuki developed norm, screw-function, moment, and Hilbert-space realizations of the same general
positivity framework \cite{Suzuki2023Norms,Suzuki2023Screw,Suzuki2025Hilbert}.


The localisation of Li-type criteria in height, and the extraction of zero-free regions
from finitely many coefficients, are also established subjects.  Freitas constructed a
family of real functions interpolating the Li coefficients and characterising zero-free
strips inside the critical strip \cite{Freitas2006}.  Droll developed variations of Li's
criterion for an extension of the Selberg class \cite{Droll2012}, and Bucur,
Ernvall-Hyt\"onen, Od\v{z}ak and Smajlovi\'c treated Li-type criteria for zero-free
regions of Dirichlet series with real coefficients \cite{BucurErnvallOdzakSmajlovic2016}.
Palojärvi proved $\tau$-Li-type criteria with explicit indices $N_1,N_2$ such that
nonnegativity of the coefficients between them forces a zero-free region, and a converse
producing a zero from a single negative coefficient \cite{Palojarvi2018}.  Sekatskii
generalized the Bombieri--Lagarias theorem and Li's criterion and gave an arithmetic
interpretation \cite{Sekatskii2014}.  For Hecke $L$-functions, Omar, Ouni and Mazhouda
computed the associated Li coefficients and conjectured the relevant monotonicity in 2014
\cite{OmarOuniMazhouda2014}; the equivalence with the Riemann hypothesis was proved
subsequently by Omar and Ouni \cite{OmarOuni2017}.  Accordingly, the detection
experiments in Subsection~\ref{subsec:detection} are presented as calibration of the
sensitivity of the reflection pivots against the scalar criterion, and no height-localised
criterion is claimed as new here.

The scalar Laguerre direction also has substantial precedent.  Coffey reformulated the Li/Keiper constants using theta--Laguerre calculus and integral representations, and decomposed the associated scalar quantities into superdominant and subdominant components whose sufficiently strong estimation would imply the Riemann hypothesis \cite{Coffey2007ThetaLaguerre}.  Related work develops the polygamma--Stieltjes-constant decomposition and further sums over the nontrivial zeros \cite{Coffey2010Polygamma,Coffey2010ZeroSums}.  Accordingly, no novelty is claimed here for the scalar appearance of associated Laguerre polynomials, for scalar archimedean--arithmetic decomposition as such, or for the general strategy of estimating a subdominant Keiper--Li component.  The distinction retained in Section~\ref{sec:arch_prime_laguerre} is the reflection-matrix cancellation: the entries are differences of scalar coefficients, and the corresponding pole counterterms disappear identically before estimation.

The orthogonal-polynomial machinery used after the reflection reduction is also standard. Hankel moments and the Cauchy--Binet, Vandermonde, Christoffel, Christoffel--Darboux, Jacobi, and Darboux constructions are classical tools of moment and orthogonal-polynomial theory; see \cite{Chihara1978,StahlTotik1992}. Infinite discrete node sets with an accumulation point have an established asymptotic theory in model situations \cite{WuLinXuZhao2016}. Section \ref{sec:determinant_decay} therefore applies these standard tools rather than presenting them as new techniques.

It is especially important to separate the determinant result here from earlier arithmetic Hankel asymptotics. Monien's determinants are
\begin{equation}
H_n^{(r)}[\zeta]=\det\bigl(\zeta(i+j+r)\bigr)_{1\le i,j\le n},
\label{eq:literature_monien_matrix}
\end{equation}
and their discrete Vandermonde representation is supported on reciprocal integers; Monien obtains the leading $-n^2\log n$ scale for the zeta model \cite{Monien2009}. Haynes--Zudilin study the different arithmetic-progression matrices $\det(\zeta(a(i+j)+b))$ and obtain general $N^2\log N$ upper decay estimates, while recording the expected coefficient $-a$ for the full leading law \cite{HaynesZudilin2015}. Thus the generic presence of $N^2\log N$ in discrete Hankel determinants is explicitly \emph{not} a novelty claim of this paper. The Li--Weil matrices here have entries formed from reflection differences of Li coefficients, and their moment nodes are the Cayley transforms of zero ordinates rather than Dirichlet-series nodes. Theorem \ref{thm:endpoint_exponent_law} identifies the common geometric mechanism in a general sufficient-condition form: the coefficient is the endpoint exponent $-\beta$. The zeta reflection matrices then have $\beta=2$. In particular, the leading $-2$ does not encode the $T\log T$ refinement of Riemann--von Mangoldt; that refinement can first affect lower orders.

Within that background, the paper singles out the pair
\begin{equation}
\E_N=\bigl(\lambda_{m+n}-\lambda_{|m-n|}\bigr),
\qquad
\Omat_N=\bigl(\lambda_{m+n-1}-\lambda_{|m-n|}\bigr),
\label{eq:literature_pair_repeat}
\end{equation}
as the two reflection-antisymmetric compressions associated with the involutions $k\mapsto-k$ and $k\mapsto1-k$.  The pair-specific results are the exact algebraic reconstruction and information kernel of Section \ref{sec:algebraic_reconstruction}, the universal rank-at-most-two displacement identities of Section \ref{sec:displacement}, the two explicit zero measures and their relation
\begin{equation}
\dd\mu_E=2(1+x)\,\dd\mu_O,
\label{eq:literature_measure_bridge}
\end{equation}
the resulting determinant bridge, Weyl-function relation and Darboux partner formulas, and the infinite Jacobi realization whose spectral support is the Cayley image of the distinct positive zero ordinates.  Theorem \ref{thm:determinant_first_order} is likewise a statement about this pair, while Theorem \ref{thm:endpoint_exponent_law} shows that its leading constant is geometric: $-2$ is the endpoint exponent of the Li--Cayley map. The precise Riemann--von Mangoldt density, including its logarithmic factor, is not recorded at this order.

The matrices are different from Suzuki's Hankel moment matrices in their native coordinates.  Their conversion to Hankel form by an explicit Chebyshev change of basis is a derived representation of the reflection coordinates, not an independent RH criterion.  Likewise, Theorem \ref{thm:main_criterion} is an exact reformulation of Li/Weil positivity for this pair and is not claimed as new arithmetic information.

Separate matrix approaches to the Riemann hypothesis include divisibility, Mertens-type, Redheffer, and multiplicative Toeplitz constructions; examples are \cite{Roesler1986,Roesler1987,Vaughan1996,Hilberdink2006,Hilberdink2015}.  Those encode different arithmetic data.  The scope of the present paper is narrower: a complete structural characterization of the two reflection compressions and the consequences that are specific to them.

\section{A de Bruijn--Newman falsification filter for matrix proof mechanisms}
\label{sec:newman_filter}

The de Bruijn--Newman deformation provides a way to separate identities that belong to the
reflection/functional-equation geometry from input that is genuinely specific to the Euler product.
In the normalization of de Bruijn, Newman, and Rodgers--Tao, set
\begin{equation}
H_t(z):=\int_0^\infty \e^{tu^2}\Phi(u)\cos(zu)\,\dd u,
\qquad t\in\mathbb R,
\label{eq:newman_heat_flow}
\end{equation}
where, in the normalization of Rodgers--Tao,
\begin{equation}
\Phi(u)=\sum_{n=1}^{\infty}
\left(2\pi^2n^4\e^{9u}-3\pi n^2\e^{5u}\right)
\exp\!\left(-\pi n^2\e^{4u}\right).
\label{eq:rodgers_tao_Phi}
\end{equation}
The kernel is even by the theta transformation and is strictly positive.  For \(u\ge0\),
strict positivity is already termwise in \eqref{eq:rodgers_tao_Phi}, since
\[
2\pi^2n^4\e^{9u}-3\pi n^2\e^{5u}
=\pi n^2\e^{5u}\bigl(2\pi n^2\e^{4u}-3\bigr)>0.
\]
At \(t=0\), \(H_0\) is a fixed rescaling of the Riemann \(\Xi\)-function.  De Bruijn proved forward preservation of
real zeros under the heat flow, Newman proved the existence of a finite constant \(\Lambda\) such
that all zeros of \(H_t\) are real exactly for \(t\ge\Lambda\), and Rodgers--Tao proved
\begin{equation}
\boxed{\Lambda\ge0.}
\label{eq:newman_nonnegative}
\end{equation}
See \cite{deBruijn1950,Newman1976,RodgersTao2020}.

For each real \(t\), introduce the centered entire function
\begin{equation}
\Xi_t(s):=
\frac{H_t\!\left(-2i(s-\tfrac12)\right)}{H_t(-i)}.
\label{eq:newman_Xit}
\end{equation}
The denominator is strictly positive.  Indeed, because \(\Phi(u)>0\),
\begin{equation}
H_t(iy)=\int_0^\infty \e^{tu^2}\Phi(u)\cosh(yu)\,\dd u>0
\qquad(y\in\mathbb R).
\label{eq:newman_imaginary_axis_positive}
\end{equation}
Consequently \(\Xi_t(x)>0\) for every real \(x\).  Evenness and reality of \(H_t\) give
\begin{equation}
\Xi_t(s)=\Xi_t(1-s)=\overline{\Xi_t(\overline s)}.
\label{eq:newman_Xit_symmetry}
\end{equation}
At \(t=0\), \(\Xi_0\) is the Riemann \(\xi\)-function divided by its value at \(1\), so the
following coefficients agree with the usual Li coefficients.  Define \(\lambda_n^{(t)}\) by the
local generating identity
\begin{equation}
\log \Xi_t\!\left(\frac1{1-z}\right)
=
\sum_{n\ge1}\frac{\lambda_n^{(t)}}{n}z^n
\qquad(|z|\text{ sufficiently small}),
\label{eq:newman_Li_generating}
\end{equation}
and form
\begin{equation}
\E_N^{(t)}
:=
\bigl(\lambda_{m+n}^{(t)}-\lambda_{|m-n|}^{(t)}\bigr)_{m,n=1}^N.
\label{eq:newman_Et}
\end{equation}
This definition avoids imposing a zero-summation convention at negative heat time.

\begin{lemma}[Even Li positivity detects the heat-flow critical line]
\label{lem:newman_pringsheim_even}
Fix \(t\in\mathbb R\).  If
\begin{equation}
\lambda_{2n}^{(t)}\ge0\qquad(n\ge1),
\label{eq:newman_even_nonnegative}
\end{equation}
then every zero of \(H_t\) is real.
\end{lemma}

\begin{proof}
Put
\[
L_t(z):=\log \Xi_t\!\left(\frac1{1-z}\right)
\]
for the branch determined near \(z=0\), and define the even coefficient series
\begin{equation}
F_t(w):=\sum_{n\ge1}\frac{\lambda_{2n}^{(t)}}{2n}w^n.
\label{eq:newman_even_series}
\end{equation}
Near the origin one has the exact identity
\begin{equation}
F_t(z^2)=\frac12\bigl(L_t(z)+L_t(-z)\bigr).
\label{eq:newman_even_part_identity}
\end{equation}

Suppose that \(H_t\) has a nonreal zero.  Under the map
\(\rho=\tfrac12+\tfrac{i}{2}z\), this gives a zero \(\rho\) of \(\Xi_t\) with
\(\Re\rho\ne\tfrac12\).  Set \(q=1-1/\rho\).  The symmetry
\(\rho\mapsto1-\rho\) sends \(q\) to \(q^{-1}\); hence one member of the reflected pair has
\(|q|<1\).  Therefore \(L_t(z)\) has a logarithmic singularity at some point of modulus strictly
less than one.  Passing to the even part cannot cancel such a singularity: \(\Xi_t\) is entire, so
all logarithmic singularities here come from zeros and carry nonnegative integer multiplicity; if
both terms in \eqref{eq:newman_even_part_identity} are singular at the same point, their
multiplicities add.  Thus the radius of convergence \(R_t\) of \(F_t\) satisfies \(R_t<1\).

Under \eqref{eq:newman_even_nonnegative}, Pringsheim's theorem forces the positive real point
\(w=R_t\) to be a singularity of \(F_t\).  But for every \(0<w<1\), with \(z=\sqrt w\), the two
arguments
\[
\frac1{1-z}>1,
\qquad
\frac1{1+z}\in\left(\frac12,1\right)
\]
are real.  Equation \eqref{eq:newman_imaginary_axis_positive} therefore shows that both terms on
the right of \eqref{eq:newman_even_part_identity} are analytic there.  Hence \(F_t\) is analytic
on \((0,1)\), a contradiction.  All zeros of \(H_t\) are therefore real.
\end{proof}

\begin{proposition}[Heat-flow sign test]
\label{prop:newman_matrix_filter}
For every \(t<\Lambda\) there exists at least one \(N\) for which
\(\E_N^{(t)}\not\succeq0\).  Consequently, since \(\Lambda\ge0\), for every
\(\varepsilon>0\) there exists \(N=N(\varepsilon)\) such that
\begin{equation}
\boxed{\E_N^{(-\varepsilon)}\not\succeq0.}
\label{eq:newman_negative_time_failure}
\end{equation}
For \(t\ge\Lambda\), the zero set is real by definition of the de Bruijn--Newman constant.
The negative-time implication above, which is the part used as a falsification test, requires only
Lemma~\ref{lem:newman_pringsheim_even} and no Bombieri--Lagarias summability assumption for the
deformed zero set.
\end{proposition}

\begin{proof}
Let \(t<\Lambda\).  By the definition of \(\Lambda\), \(H_t\) has a nonreal zero.  If
\(\E_N^{(t)}\succeq0\) for every \(N\), then its diagonal would give
\[
\lambda_{2N}^{(t)}=(\E_N^{(t)})_{NN}\ge0
\qquad(N\ge1),
\]
contradicting Lemma~\ref{lem:newman_pringsheim_even}.  Hence some finite even reflection block
fails to be positive semidefinite.  Taking \(t=-\varepsilon\) and using \(\Lambda\ge0\) proves
\eqref{eq:newman_negative_time_failure}.
\end{proof}

\begin{corollary}[Falsification principle]
\label{cor:newman_falsification}
No proof of \(\E_N\succeq0\) for all \(N\) can be based solely on identities and inequalities that
remain valid, with the same sign content, for every negative-time family \(H_{-\varepsilon}\).
Any successful proof must use an input that distinguishes the zeta function from those heat
deformations.
\end{corollary}

\begin{proof}
If such a proof used only statements that remained valid with the same sign content for every
\(H_{-\varepsilon}\), it would prove \(\E_N^{(-\varepsilon)}\succeq0\) for every \(N\).
This contradicts Proposition~\ref{prop:newman_matrix_filter}.  Hence at least one
zeta-specific input that fails under negative heat deformation is necessary.
\end{proof}

\begin{remark}[What survives and what does not]
\label{rem:newman_survival}
The algebraic reflection form, the discrete wave equation, rank-two displacement, the Schur
bordering identities, the inverse displacement identity, and the finite-rank update formulas depend
only on the coefficient/reflection algebra and therefore survive the heat deformation.  The
positive-measure consequences do \emph{not} survive once the zeros leave the line, because their
positivity is exactly the property being tested.

There is also a useful distinction on the archimedean--arithmetic split.  The sequence \(A_n\) is
\emph{fixed by decree as a reference sequence}; it is not inherited from \(H_t\).  After making that
choice one may write
\(\lambda_n^{(t)}=A_n+Z_n^{(t)}\), and then the explicit matrix \(\mathsf S_N\) and the scalar
\(c_\infty\) are identical by construction.  What fails for \(H_t\) is precisely the zeta-specific
von Mangoldt--Laguerre representation of \(Z_n\), equivalently the representation of
\(\widehat{\mathsf Z}_N^E\) through \(\Lambda(n)\), \(\psi(x)-x\), and the Euler product.  Thus,
within the machinery actually developed in this paper, the prime--continuum representation is the
only identified zeta-specific door through which a sign proof can distinguish \(\xi\) from
\(H_{-\varepsilon}\).  This statement is internal to the present framework; it does not claim that
every conceivable proof of RH must literally use formula \eqref{eq:prime_continuum_matrix_discrepancy}.
\end{remark}

\begin{remark}[Criticality]
\label{rem:newman_criticality}
Rodgers--Tao gives \(\Lambda\ge0\), while \(\RH\) is equivalent to \(\Lambda\le0\).
Hence, \emph{conditional on RH}, one has \(\Lambda=0\): the Riemann \(\Xi\)-function then lies
exactly on the heat-flow boundary of the all-real-zero region.  The very small reflection pivots,
near-unit low-zero leverage, and other finite margins reported below are consistent with such
criticality, but they do not establish \(\Lambda=0\), nor does the heat-flow theorem by itself
identify a quantitative rate for any one of those finite margins.
\end{remark}

\section{Computational realization}
\label{sec:computational_realization}

For a supplied list $\lambda_1,\dots,\lambda_{2N}$, both matrices are formed in $O(N^2)$ scalar operations. With the zero final component convention in \eqref{eq:rE_definition}--\eqref{eq:rO_definition}, the displacement checks also require no coefficient beyond $\lambda_{2N}$. The basic numerical diagnostics are:
\begin{enumerate}[label=(\roman*)]
\item the smallest eigenvalues of $\E_N$ and $\Omat_N$;
\item the determinant ratios $\Delta_n^E/\Delta_{n-1}^E$ and $\Delta_n^O/\Delta_{n-1}^O$;
\item the Cholesky factors and Schur residuals in \eqref{eq:E_schur_complement}--\eqref{eq:O_schur_complement};
\item the initial Jacobi recurrence coefficients determined by the available moments, together with eigenvalues of the corresponding finite Jacobi truncations as quadrature nodes;
\item the displacement residuals in \eqref{eq:E_displacement}--\eqref{eq:O_displacement}, which provide exact algebraic consistency checks;
\item the Christoffel identity \eqref{eq:christoffel_measure_relation} and determinant bridge \eqref{eq:determinant_bridge}.
\end{enumerate}


\begin{experiment}[Single-zero leverage and transport at $N=10$]
\label{exp:single_zero_transport}
The exact formulas above were evaluated for $\E_{10}$.  The Li coefficients $\lambda_1,\ldots,\lambda_{20}$ were computed at $160$ decimal digits from the Taylor expansion of $\log\xi(1+z)$, and the required ordinates were evaluated at high precision.  The determinant responses were computed from the $2\times2$ formula \eqref{eq:exact_moved_zero_ratio}, not by subtracting two ill-conditioned determinants.  As an independent audit, for $j=1,10,100$ with a $10\%$ local-gap motion, direct multiprecision recomputation of the perturbed $10\times10$ determinant agrees with \eqref{eq:exact_moved_zero_ratio} to more than $80$ decimal digits.

For a zero $\gamma_j$, let
\begin{equation}
h_j:=\gamma_{j+1}-\gamma_j,
\qquad
\gamma_j(c):=\gamma_j+c h_j.
\label{eq:local_gap_transport}
\end{equation}
The table gives the even leverage $\ell_{10,\gamma_j}^E$, the exact determinant ratio after deletion through $1-\ell_{10,\gamma_j}^E$, and the logarithmic determinant response when the zero is moved by $10\%$ and $50\%$ of its next local gap.

\begin{center}
\small
\begin{tabular}{rrrrrr}
\toprule
$j$ & $\gamma_j$ & $\ell_{10,\gamma_j}^E$ & $1-\ell_{10,\gamma_j}^E$
& $\log(\Delta'/\Delta)_{c=0.1}$ & $\log(\Delta'/\Delta)_{c=0.5}$\\
\midrule
1   & 14.1347  & $1-6.90\times10^{-19}$ & $6.90\times10^{-19}$ & $-2.2354$ & $-10.9973$\\
2   & 21.0220  & $0.9999999998$ & $1.58\times10^{-10}$ & $-0.9565$ & $-5.1352$\\
3   & 25.0109  & $0.9999995136$ & $4.86\times10^{-7}$ & $-0.8360$ & $-4.8068$\\
4   & 30.4249  & $0.9975186583$ & $2.48\times10^{-3}$ & $-0.3917$ & $-2.2879$\\
5   & 32.9351  & $0.9644640587$ & $3.55\times10^{-2}$ & $-0.1175$ & $-1.1180$\\
10  & 49.7738  & $0.2310567145$ & $0.76894329$ & $-6.20\times10^{-3}$ & $-3.24\times10^{-2}$\\
20  & 77.1448  & $0.0656564393$ & $0.93434356$ & $-7.04\times10^{-4}$ & $-3.44\times10^{-3}$\\
50  & 143.1118 & $0.0160389932$ & $0.98396101$ & $-9.89\times10^{-5}$ & $-4.89\times10^{-4}$\\
100 & 236.5242 & $0.0034507536$ & $0.99654925$ & $-3.80\times10^{-6}$ & $-1.89\times10^{-5}$\\
200 & 396.3819 & $0.0016652308$ & $0.99833477$ & $-8.26\times10^{-7}$ & $-4.12\times10^{-6}$\\
\bottomrule
\end{tabular}
\end{center}

The total leverage is exactly $10$.  Numerically, the cumulative even leverage carried by the first $M$ positive zeros is
\begin{center}
\small
\begin{tabular}{rcccccccccc}
\toprule
$M$&1&2&3&4&5&10&20&50&100&200\\
\midrule
$\sum_{j\le M}\ell_{10,\gamma_j}^E$
&1.0000&2.0000&3.0000&3.9975&4.9620&6.7290&7.7976&8.6834&9.0382&9.2658\\
\bottomrule
\end{tabular}
\end{center}
Thus the first five zeros already carry about $49.6\%$ of the total leverage, the first ten about $67.3\%$, and the first two hundred about $92.7\%$.  This is a property of the \emph{full matrix determinant}: the leverage is obtained from $\E_{10}^{-1}$ and is the exact marginal of the complete Cauchy--Binet sum.

For illustration of the Li-side representation, moving the tenth zero by $0.1h_{10}=0.3196489000\ldots$ changes the contribution of that conjugate pair by
\begin{equation}
\delta\lambda_1=-5.13\times10^{-6},\quad
\delta\lambda_5=-1.28\times10^{-4},\quad
\delta\lambda_{10}=-5.10\times10^{-4},\quad
\delta\lambda_{20}=-1.999\times10^{-3},
\end{equation}
while the exact determinant response is
\begin{equation}
\log\frac{\det\E_{10}'}{\det\E_{10}}
=-0.00620213676254\ldots,
\qquad
\frac{\det\E_{10}'}{\det\E_{10}}
=0.9938170567868\ldots .
\label{eq:N10_tenth_zero_response}
\end{equation}
For comparison, the same $10\%$ local-gap motion gives
\begin{equation}
\frac{\det\E_{10}' }{\det\E_{10}}=0.1069489018\ldots\quad (j=1),
\qquad
\frac{\det\E_{10}' }{\det\E_{10}}=0.9999962023\ldots\quad (j=100).
\end{equation}
This strong localization is consistent with Theorem~\ref{thm:high_zero_leverage_decay}.
\end{experiment}

\begin{experiment}[Conditioning of the reflection matrices]
\label{exp:conditioning}
Direct determinant evaluation is unstable at very small orders.  In the computation used here, Li coefficients carried to $40$ significant decimal digits and passed to a direct $LU$ determinant produced a negative value for $\Delta_7^E$.  This sign is numerical: under $\RH$, Theorem \ref{thm:main_criterion} gives $\Delta_7^E>0$, and the zero-side Cauchy--Binet representation is a sum of nonnegative terms.

As a diagnostic, retaining only one Cauchy--Binet term, formed from tabulated low ordinates, gives positive one-term estimates whose base-$10$ logarithms for $N=1,2,4,7,10,14,20$ are approximately
\begin{equation}
-1.7,\ -8.9,\ -42.0,\ -142.5,\ -308.0,\ -634.6,\ -1361.2.
\label{eq:experiment_determinant_lower_estimates}
\end{equation}
These decimal values are numerical estimates rather than interval-certified bounds.  Their role is only to illustrate the scale proved analytically in Theorem \ref{thm:determinant_first_order}.
\end{experiment}

\begin{experiment}[Critical sine-polynomial variational identity and finite profile data]
\label{exp:critical_vector_scaling}
Consider the even prime--archimedean normal form introduced in Section~\ref{sec:arch_prime_laguerre},
and let $v^{(N)}$ denote a unit critical direction obtained numerically from the
smallest-eigenvalue problem, with a consistent global sign convention.  A recheck using a
$1500$-ordinate model, inverse iteration and $250$ decimal working digits gives peak indices
$m=3,3,3,5$ for $N=12,20,32,48$.  Representative normalized amplitudes are
\begin{center}
\small
\begin{tabular}{rccc}
\toprule
$N$ & $m/N=0.05$ & $m/N=0.30$ & $m/N=0.50$\\
\midrule
$12$ & $0.636$ & $0.750$ & $0.206$\\
$20$ & $0.493$ & $0.523$ & $0.032$\\
$32$ & $0.893$ & $0.308$ & $0.011$\\
$48$ & $0.570$ & $0.188$ & $0.056$\\
\bottomrule
\end{tabular}
\end{center}
Over this finite range the mass remains concentrated at low indices and the data do not support a
stable nonzero profile in the variable $m/N$.  They do not establish the absence of such a profile
as $N\to\infty$.

The correct exact interpretation is instead variational.  Under $\RH$, for every real vector $v=(v_1,\ldots,v_N)^T$ define
\begin{equation}
 P_v(\theta):=\sum_{m=1}^N v_m\sin(m\theta).
 \label{eq:critical_sine_polynomial}
\end{equation}
Then the sine Gram representation gives
\begin{equation}
 \boxed{
 v^T\E_N v
 =4\sum_{\gamma>0}\bigl|P_v(\theta_\gamma)\bigr|^2.}
 \label{eq:critical_sine_energy}
\end{equation}
Hence the smallest-eigenvalue direction is exactly the unit coefficient vector whose sine polynomial minimizes the discrete squared energy on the Cayley nodes $\{\theta_\gamma\}$.  Since $\theta_\gamma\downarrow0$, suppression of this energy naturally favors high-order flatness at the origin.  Taylor expansion gives
\begin{equation}
 P_v(\theta)
 =\sum_{j\ge0}\frac{(-1)^j\theta^{2j+1}}{(2j+1)!}
 \left(\sum_{m=1}^N m^{2j+1}v_m\right),
 \label{eq:critical_sine_moment_expansion}
\end{equation}
so the exact conditions
\begin{equation}
 \sum_{m=1}^N m^{2j+1}v_m=0,
 \qquad 0\le j<r,
 \label{eq:critical_odd_moment_flatness}
\end{equation}
are equivalent to $P_v(\theta)=O(\theta^{2r+1})$ at the origin.  The observed strictly alternating signs are consistent with these successive odd-moment cancellations, although alternation alone does not prove maximal flatness.

Numerically, the visible core at the largest tested orders is concentrated on low indices of size comparable with
\begin{equation}
 \theta_1^{-1}
 =\left(2\arctan\frac{1}{2\gamma_1}\right)^{-1}
 \approx 14.1.
 \label{eq:first_zero_visible_core_scale}
\end{equation}
This suggests that the first Cayley node sets the scale of the dominant visible core.  It does not imply finite support: if the minimum energy tends to zero, an exactly fixed finite-support limiting polynomial would vanish at infinitely many nodes accumulating at $0$ and hence would be identically zero.  The numerically tiny high-index tail may therefore be essential for purchasing additional orders of flatness as $N$ grows.
\end{experiment}

\begin{remark}[Matrix-specific flatness problem]
\label{rem:critical_vector_continuum}
The finite computations do not exhibit a stable $m/N$ continuum profile over the tested orders.
The matrix-specific object is instead the sequence of odd moments
\begin{equation}
 M_{2j+1}^{(N)}:=\sum_{m=1}^N m^{2j+1}v_m^{(N)}.
 \label{eq:critical_odd_moments}
\end{equation}
A direct high-precision diagnostic is to determine how many consecutive moments $M_1^{(N)},M_3^{(N)},\ldots$ are suppressed to the numerical scale of the smallest-eigenvalue computation.  Equivalently, if
\begin{equation}
 V_{r,N}:=\bigl(m^{2j+1}\bigr)_{\substack{0\le j<r\\1\le m\le N}},
 \label{eq:odd_moment_matrix}
\end{equation}
then one should measure the distance of $v^{(N)}$ to $\ker V_{r,N}$ as $r$ grows.  A growing flatness order would explain how additional matrix dimensions can reduce the sine energy without producing a macroscopic $m/N$ profile.  This phenomenon has no scalar Keiper--Li analogue because it is a statement about the geometry of a smallest-eigenvalue direction in the reflection Gram matrix.
\end{remark}

\begin{remark}[Numerical evaluation]
\label{rem:precision_growth}
Theorem \ref{thm:determinant_first_order} shows that
\begin{equation}
\Delta_N^{E,O}=\exp\!\bigl(-(2+o(1))N^2\log N\bigr),
\label{eq:determinant_numerical_scale}
\end{equation}
so forming a determinant directly in fixed precision rapidly loses meaning.  The theorem by itself does not prescribe a sharp number of working digits for a stable factorization, because that also depends on the smallest pivots and on entrywise cancellation.  Numerically one should work with multiprecision $LDL^T$ or Cholesky pivots and accumulate $\log\Delta_N^{E,O}=\sum_n\log d_n^{E,O}$; when certified positivity is required, interval arithmetic or certified zero-side estimates are preferable.  The Cauchy--Binet formulas are especially useful for sign certification under $\RH$ because their individual terms are nonnegative.
\end{remark}

The rank-at-most-two displacement structure is particularly useful computationally: structured Toeplitz--Hankel algorithms can exploit low displacement rank, and Corollary \ref{cor:inverse_displacement} shows that this feature survives inversion. The matrices therefore offer a natural setting for numerical experiments on Li coefficients without destroying the underlying reflection geometry.


\subsection{Off-line detection: how early a reflection pivot sees a zero off the line}
\label{subsec:detection}

The experiments of this subsection use a deliberately constructed \emph{model} measure and
therefore do not bear on $\zeta$ itself; their purpose is to calibrate the sensitivity of
the reflection pivots against the scalar Li criterion.

\begin{definition}[Count- and symmetry-preserving deformation]
\label{def:offline_model}
Fix an ordinate $\gamma$ and $\beta\in(0,1)$.  The model $M(\beta,\gamma)$ is obtained
from the on-line configuration by replacing the conjugate pair
$\{\tfrac12\pm i\gamma\}$ by the four points
$\{\beta\pm i\gamma,\;1-\beta\pm i\gamma\}$, each carrying weight $\tfrac12$; all other
ordinates remain on the critical line.  The resulting coefficients
\begin{equation}
 \lambda_k^{M}
 =\sum_{\rho}w_\rho\Bigl[1-\bigl(1-\rho^{-1}\bigr)^k\Bigr]
 \label{eq:offline_lambda}
\end{equation}
are real, the total zero mass is unchanged, the set is invariant under both
$\rho\mapsto\bar\rho$ and $\rho\mapsto1-\rho$, and $M(\tfrac12,\gamma)$ coincides
identically with the on-line configuration.  The last property is what makes
$\beta\mapsto\lambda_k^M$ a genuine deformation with vanishing signal at $\beta=1/2$;
models that replace a pair by an unweighted quadruple fail it and produce a spurious
$\beta$-independent signal.
\end{definition}

\begin{experiment}[Detection index of the reflection pivot]
\label{exp:detection_index}
Let $N^*(\beta,\gamma)$ be the smallest $N$ for which the Cholesky pivot $p_N^E$ of the
model $M(\beta,\gamma)$ fails to be positive.  Computations use $K=300$ background
ordinates, $N\le34$ and $900$ decimal working digits.  As a control, the unperturbed
on-line configuration has all pivots strictly positive up to $N=34$, with
$\log_{10}p_{34}=-255.2$, so the working precision exceeds the pivot scale by a wide
margin and every reported sign change is genuine.
\begin{center}
\small
\begin{tabular}{rrrrrr}
\toprule
$\gamma$ & $\beta=0.45$ & $\beta=0.49$ & $\beta=0.499$ & $\beta=0.4999$ & $\gamma/2$\\
\midrule
$14.1347$ & $5$  & $6$  & $7$  & $8$  & $7.07$\\
$21.0220$ & $8$  & $9$  & $10$ & $11$ & $10.51$\\
$25.0109$ & $10$ & $11$ & $12$ & $14$ & $12.51$\\
$30.4249$ & $13$ & $14$ & $16$ & $17$ & $15.21$\\
$32.9351$ & $14$ & $16$ & $17$ & $18$ & $16.47$\\
$40.9187$ & $19$ & $20$ & $22$ & $23$ & $20.46$\\
$49.7738$ & $24$ & $26$ & $28$ & $29$ & $24.89$\\
$65.1125$ & $34$ & $>34$ & $>34$ & $>34$ & $32.56$\\
$77.1448$ & $>34$ & $>34$ & $>34$ & $>34$ & $38.57$\\
\bottomrule
\end{tabular}
\end{center}
The detection index remained unchanged across the tested background truncations: $K=300$, $600$ and $1200$ give
identical values, for instance $13,14,16$ at $\gamma=30.4249$ for the three smallest
$\beta$.  This is a numerical stability statement for the tested truncations only; no
background-independence theorem is inferred from it.

The natural variable is not $\gamma$ but the Cayley angle.  With
$\theta_\gamma=2\arctan(1/2\gamma)$ and $\beta=0.499$,
\begin{center}
\small
\begin{tabular}{rcccccc}
\toprule
$\gamma$ & $14.1347$ & $21.0220$ & $25.0109$ & $30.4249$ & $40.9187$ & $49.7738$\\
$N^*\theta_\gamma$ & $0.495$ & $0.476$ & $0.480$ & $0.526$ & $0.538$ & $0.563$\\
\bottomrule
\end{tabular}
\end{center}
so that
\begin{equation}
 N^*\,\theta_\gamma\approx\frac12 .
 \label{eq:detection_nyquist}
\end{equation}
The observed values suggest a resolution scale in the Cayley angle.  The top available
sine mode is $\sin(N\theta)$, so departure from the small-angle linear regime is a
natural heuristic for the numerical scale, but the data do not prove that this mechanism
determines the threshold.  Since $\theta_\gamma\sim1/\gamma$, the measured values in
\eqref{eq:detection_nyquist} are consistent with $N^*\approx\gamma/2$ over the tested
range.  The dependence on $\beta$ is also empirical: in the table $N^*$ increases as
$\beta\uparrow1/2$ toward the observed near-critical scale.
\end{experiment}

\begin{experiment}[Comparison with the scalar Li criterion]
\label{exp:vs_scalar_li}
At fixed low height the numerical separation from the scalar Li crossing estimate is large.  For an off-line pair the
Cayley image $z=(\rho-1)/\rho$ of the zero with $\Re\rho<1/2$ has $|z|>1$, and the scalar
coefficient $\lambda_n$ first becomes negative near the crossing of $|z|^n$ with the
on-line growth $\tfrac{n}{2}\log n$.  Solving
$n\log|z|=\log\bigl(\tfrac{n}{2}\log n\bigr)$ at $\gamma=14.1347$:
\begin{center}
\small
\begin{tabular}{lrrr}
\toprule
$\beta$ & $\log|z|$ & scalar $n$ & matrix $N^*$\\
\midrule
$0.45$   & $2.50\times10^{-4}$ & $5.0\times10^{4}$ & $5$\\
$0.49$   & $5.00\times10^{-5}$ & $2.9\times10^{5}$ & $6$\\
$0.499$  & $5.00\times10^{-6}$ & $3.4\times10^{6}$ & $7$\\
$0.4999$ & $5.00\times10^{-7}$ & $3.9\times10^{7}$ & $8$\\
\bottomrule
\end{tabular}
\end{center}
The reflection pivot at $N^*=8$ uses $\lambda_1,\dots,\lambda_{16}$.  These figures show
a large finite-height separation between the model pivot detector and the scalar crossing
estimate.  The same data also show that the required matrix dimension grows with height,
so the comparison is not a uniform complexity theorem and should not be read as a general
improvement of Li's criterion.
\end{experiment}

\begin{experiment}[Relative signal across raw and processed observables]
\label{exp:signal_law}
Define the relative signal of an observable $X$ as
$|X(M(\beta,\gamma))-X(M(\tfrac12,\gamma))|/|X(M(\tfrac12,\gamma))|$, maximized over
$n\le30$ for $X=\lambda_n$ and over $N\le10$ for $X=\kappa_N$ and $X=F_N$.  With
$K=300$ and $400$ working digits:
\begin{center}
\small
\begin{tabular}{rlccc}
\toprule
$\gamma$ & $\beta$ & $\lambda_n$ & $\kappa_N$ & $F_N$\\
\midrule
$30.425$  & $0.45$ & $4.07\times10^{-7}$  & $2.98\times10^{-7}$  & $5.05\times10^{-7}$\\
$49.774$  & $0.45$ & $5.79\times10^{-8}$  & $7.56\times10^{-10}$ & $4.76\times10^{-10}$\\
$77.145$  & $0.45$ & $1.07\times10^{-8}$  & $3.00\times10^{-11}$ & $1.51\times10^{-11}$\\
$143.112$ & $0.45$ & $9.29\times10^{-10}$ & $4.23\times10^{-13}$ & $3.65\times10^{-13}$\\
$248.102$ & $0.45$ & $1.04\times10^{-10}$ & $3.85\times10^{-14}$ & $3.85\times10^{-14}$\\
$248.102$ & $0.49$ & $4.15\times10^{-12}$ & $1.54\times10^{-15}$ & $1.54\times10^{-15}$\\
\bottomrule
\end{tabular}
\end{center}
A least-squares fit in the logarithms gives exponent $2.000$ in $(\tfrac12-\beta)$ and
$-3.938$ in $\gamma$.  Over the range tested, the data are therefore consistent with the
empirical scaling
\begin{equation}
 \text{signal}\;\text{ of order }\;\Bigl(\tfrac12-\beta\Bigr)^2\gamma^{-4}.
 \label{eq:signal_law}
\end{equation}
The symmetry $\beta\mapsto1-\beta$ forces an even expansion in $\beta-\tfrac12$; the
observed quadratic leading term and the fourth-power height law are empirical, not proved
asymptotics.  The numerical attenuation is strongly height dependent.  For the five
$\beta=0.45$ rows, the base-$10$ logarithms of the ratio between the largest raw
$\lambda_n$ signal and the $\kappa_N$ signal are, respectively,
\begin{equation}
 0.14,\quad 1.88,\quad 2.55,\quad 3.34,\quad 3.43.
 \label{eq:signal_ratio_lambda_kappa}
\end{equation}
At the lowest tested height $F_N$ actually has a slightly larger relative signal than the
raw coefficients: $5.05\times10^{-7}$ versus $4.07\times10^{-7}$.  Only at the higher
heights does a multi-order attenuation appear.  These numbers quantify numerical
conditioning; they do not imply that forming determinants destroys mathematical
information, nor that an exact matrix identity, recurrence, convexity argument or sign
mechanism is weaker than the scalar sequence.
\end{experiment}

\begin{experiment}[What the band detects, and what it does not]
\label{exp:band_detection}
Corollary~\ref{cor:kappa_le_4} is a necessary condition; the question is whether it is
sharp enough to be a diagnostic.  At $\gamma=14.1347$, $\beta=0.45$, where the pivot first
fails at $N^*=5$:
\begin{center}
\small
\begin{tabular}{rll}
\toprule
$N$ & $\kappa_N$ (model) & $\kappa_N$ (on-line)\\
\midrule
$3$ & $3.9982260960$ & $3.9982215228$\\
$4$ & $3.9989005054$ & $3.9987896058$\\
$5$ & $3.9846420023$ & $3.9989285277$\\
$6$ & $\mathbf{4.0092903964}$ & $3.9991406580$\\
$7$ & $3.9991204664$ & $3.9994048590$\\
$8$ & $3.9993769054$ & $3.9994491550$\\
\bottomrule
\end{tabular}
\end{center}
The upper bound $\kappa_N\le4$ is violated at $N=6$, one step after the pivot sign change,
and the violation is by $2.3\times10^{-3}$, four orders above the working precision.  The
lower bound is never approached.  In this deformation the band is therefore a nontrivial numerical detector rather than merely a consistency check, but it detects strictly later than the pivot itself, so it adds robustness rather than sensitivity.  Second, the proof of the bound $\kappa_N\le4$ under RH uses $\mathcal J_O\preceq I$, equivalently the support condition for the positive RH reflection measure.  The synthetic off-line deformation no longer belongs to that positive-measure setting, but it does violate the resulting determinant inequality one step after the pivot changes sign.  A wider band of the shape
$(1-\Tr\Kmat)\,\kappa_N^{\rm ref}$, which is what a first-order estimate of $F_N/F_{N-1}$
suggests, is satisfied by the model at every $N$ including $N=6$ and therefore does not detect this deformation.  In this example the sharp constant is essential.
\end{experiment}



\section{The archimedean--prime split in reflection and Laguerre coordinates}
\label{sec:arch_prime_laguerre}

The reflection matrices admit a second exact realization which is independent of the zero-side moment representation.  The scalar theta--Laguerre formulation of the Li/Keiper coefficients, and scalar decompositions into dominant and subdominant components, have classical precedents in Coffey's work \cite{Coffey2007ThetaLaguerre,Coffey2010Polygamma}.  The purpose of this section is therefore not to claim the scalar Laguerre representation as new.  What is specific to the present reflection pair is that the entries are coefficient differences: the pole counterterms required by the individual scalar coefficients cancel identically inside $\lambda_{m+n}-\lambda_{|m-n|}$ (and analogously in the odd family).  This leaves a matrix pairing directly against the prime-number-theorem error before any estimate is made.  The section records that cancellation, the resulting matrix cosine/positive-definite structure, and the scalar coefficient criterion only as a diagnostic comparison with the classical Keiper--Li mechanism.

\subsection{An analytic split of the Li coefficients}

Put
\begin{equation}
 s(z):=\frac{1}{1-z},
 \qquad |z|\ll1,
 \label{eq:laguerre_cayley_s}
\end{equation}
and factor
\begin{equation}
 \xi(s)=
 \underbrace{\frac{s}{2}\pi^{-s/2}\Gamma\!\left(\frac{s}{2}\right)}_{\Xi_\infty(s)}
 \underbrace{(s-1)\zeta(s)}_{\Xi_{\rm ar}(s)}.
 \label{eq:xi_arch_arith_factorization}
\end{equation}
Both factors are analytic and nonzero at $s=1$.  Choose the logarithms which are real there and define, for $k\ge1$,
\begin{align}
 A_k&:=k[z^k]\log \Xi_\infty(s(z)),
 \label{eq:Ak_definition}\\
 Z_k&:=k[z^k]\log \Xi_{\rm ar}(s(z)).
 \label{eq:Zk_definition}
\end{align}
Then
\begin{equation}
 \lambda_k=A_k+Z_k.
 \label{eq:lambda_arch_arith_split}
\end{equation}
The archimedean coefficient is explicit:
\begin{equation}
 A_k=
 1-\frac{k}{2}(\gamma_E+\log 4\pi)
 +\sum_{j=2}^{k}(-1)^j\binom{k}{j}(1-2^{-j})\zeta(j),
 \label{eq:Ak_closed_form}
\end{equation}
where $\gamma_E$ is Euler's constant.

\begin{proposition}[Regularized von Mangoldt--Laguerre formula]
\label{prop:Zk_regularized_laguerre}
For every fixed integer $k\ge1$,
\begin{equation}
 Z_k
 =-\lim_{X\to\infty}
 \left[
 \sum_{n\le X}\frac{\Lambda(n)}{n}
 L_{k-1}^{(1)}(\log n)
 +L_k(\log X)-1
 \right],
 \label{eq:Zk_regularized_laguerre}
\end{equation}
where $L_k=L_k^{(0)}$ and $L_{k-1}^{(1)}$ are Laguerre polynomials.  Equivalently, with
\begin{equation}
 E(t):=\psi(t)-t,
 \label{eq:PNT_error_E}
\end{equation}
one has the Stieltjes identity
\begin{equation}
 \boxed{
 Z_k=\int_{1^-}^{\infty}
 L_k'(\log t)\,\frac{\dd E(t)}{t}.}
 \label{eq:Zk_PNT_error_laguerre}
\end{equation}
\end{proposition}

\begin{proof}
For $\Re s>1$,
\begin{equation}
 \log\zeta(s)=\sum_{n\ge2}\frac{\Lambda(n)}{\log n}\,n^{-s}.
 \label{eq:logzeta_von_mangoldt}
\end{equation}
Since
\begin{equation}
 n^{-s(z)}
 =n^{-1}\exp\!\left(-\frac{(\log n)z}{1-z}\right)
 \label{eq:n_minus_s_cayley}
\end{equation}
and
\begin{equation}
 \sum_{r\ge0}L_r^{(1)}(u)z^r
 =\frac{1}{(1-z)^2}\exp\!\left(-\frac{uz}{1-z}\right),
 \label{eq:Laguerre_gen_alpha1}
\end{equation}
coefficient extraction gives
\begin{equation}
 k[z^k]n^{-s(z)}
 =-\frac{\log n}{n}L_{k-1}^{(1)}(\log n).
 \label{eq:Laguerre_coefficient_n_minus_s}
\end{equation}
The factor $s(z)-1=z/(1-z)$ supplies the canonical pole counterterm.  Using
\begin{equation}
 L_k'(u)=-L_{k-1}^{(1)}(u)
 \label{eq:Laguerre_derivative_identity}
\end{equation}
and hence
\begin{equation}
 \int_0^{\log X}L_{k-1}^{(1)}(u)\,\dd u
 =1-L_k(\log X),
 \label{eq:Laguerre_counterterm_integral}
\end{equation}
we obtain \eqref{eq:Zk_regularized_laguerre}.  The bracket in \eqref{eq:Zk_regularized_laguerre} is exactly
\begin{equation}
 \int_{1^-}^{X}
 L_{k-1}^{(1)}(\log t)\,\frac{\dd(\psi(t)-t)}{t}.
\end{equation}
Applying \eqref{eq:Laguerre_derivative_identity} gives \eqref{eq:Zk_PNT_error_laguerre}.  For fixed $k$, passage to infinity follows from the classical prime-number-theorem error bound, since the Laguerre factor has fixed polynomial degree.
\end{proof}

The first value is
\begin{equation}
 Z_1=\gamma_E,
 \label{eq:Z1_euler}
\end{equation}
so $A_1+Z_1=\lambda_1$ gives the familiar numerical value of the first Li coefficient.


\subsection{Arithmetic reconstruction of the Li coefficients without zeros}
\label{subsec:arithmetic_reconstruction}

The split \eqref{eq:lambda_arch_arith_split} is not only structural.  Together with the
closed form \eqref{eq:Ak_closed_form} for $A_k$ it yields a construction of
$\lambda_1,\dots,\lambda_K$ from the Stieltjes constants alone, with no zero of
$\zeta$ entering at any stage.  This matters here for a specific reason: every numerical
statement in this paper concerns positivity of matrices built from $\lambda$, and a
$\lambda$ computed by summing over a truncated list of ordinates would make those
statements circular.  The construction below removes that circularity.

\begin{proposition}[Stieltjes recursion for the arithmetic coefficients]
\label{prop:stieltjes_recursion}
Let $\gamma_m$ denote the Stieltjes constants, so that
\begin{equation}
 \zeta(s)=\frac{1}{s-1}+\sum_{m\ge0}\frac{(-1)^m}{m!}\gamma_m(s-1)^m .
 \label{eq:stieltjes_expansion}
\end{equation}
Put $b_0:=1$ and, for $k\ge1$,
\begin{equation}
 \boxed{\;
 b_k:=\sum_{m=0}^{k-1}(-1)^m\binom{k-1}{m}\frac{\gamma_m}{m!}\;}
 \label{eq:bk_stieltjes}
\end{equation}
Then $\Xi_{\rm ar}(s(z))=\sum_{k\ge0}b_kz^k$ for $|z|$ small, and the coefficients
$Z_k$ of \eqref{eq:Zk_definition} satisfy the recursion
\begin{equation}
 \boxed{\;
 Z_k=k\,b_k-\sum_{j=1}^{k-1}Z_j\,b_{k-j}\;}
 \qquad(k\ge1).
 \label{eq:Zk_recursion}
\end{equation}
Consequently $\lambda_k=A_k+Z_k$ is determined by $\gamma_0,\dots,\gamma_{k-1}$ and
$\zeta(2),\dots,\zeta(k)$.
\end{proposition}

\begin{proof}
Multiplying \eqref{eq:stieltjes_expansion} by $s-1$ gives
\begin{equation}
 \Xi_{\rm ar}(s)=(s-1)\zeta(s)
 =1+\sum_{m\ge0}\frac{(-1)^m}{m!}\gamma_m(s-1)^{m+1}.
\end{equation}
Under $s=s(z)=1/(1-z)$ one has $s-1=z/(1-z)$, and
\begin{equation}
 \left(\frac{z}{1-z}\right)^{m+1}
 =\sum_{k\ge m+1}\binom{k-1}{m}z^k ,
\end{equation}
since $[z^k]\,z^{m+1}(1-z)^{-(m+1)}=\binom{k-m-1+m}{m}=\binom{k-1}{m}$.  Collecting the
coefficient of $z^k$ gives \eqref{eq:bk_stieltjes}, and $b_0=1$ because
$\Xi_{\rm ar}(1)=1$.

For \eqref{eq:Zk_recursion}, write $F(z):=\log\Xi_{\rm ar}(s(z))=\sum_{k\ge1}(Z_k/k)z^k$,
which is legitimate because $\Xi_{\rm ar}$ is analytic and nonvanishing at $s=1$ and
$F(0)=0$.  Then $F'(z)\,\Xi_{\rm ar}(s(z))=\frac{\dd}{\dd z}\Xi_{\rm ar}(s(z))$, that is
\begin{equation}
 \Bigl(\sum_{k\ge1}Z_kz^{k-1}\Bigr)\Bigl(\sum_{j\ge0}b_jz^j\Bigr)
 =\sum_{k\ge1}k\,b_kz^{k-1}.
\end{equation}
Comparing the coefficient of $z^{k-1}$ and using $b_0=1$ gives
$\sum_{j=0}^{k-1}Z_{k-j}b_j=k\,b_k$, which is \eqref{eq:Zk_recursion}.
\end{proof}

\begin{remark}[Why this is the right normalization]
\label{rem:no_pole_in_recursion}
The pole of $\zeta$ has been absorbed once and for all by the factor $s-1$, so
\eqref{eq:bk_stieltjes}--\eqref{eq:Zk_recursion} involve no regularization and no
principal value.  This is the coefficient-side counterpart of the counterterm
cancellation in Proposition~\ref{prop:Zk_regularized_laguerre}: the same pole is removed
analytically there and combinatorially here.
\end{remark}

\begin{experiment}[Arithmetic $\lambda$ and its independent audit]
\label{exp:arithmetic_lambda}
Proposition~\ref{prop:stieltjes_recursion} and \eqref{eq:Ak_closed_form} were evaluated at
$300$ decimal working digits for $k\le50$.  Representative values:
\begin{center}
\small
\begin{tabular}{rrrr}
\toprule
$k$ & $Z_k$ & $A_k$ & $\lambda_k=A_k+Z_k$\\
\midrule
$1$ & $0.577215664901533$ & $-0.554119955935412$ & $0.0230957089661210$\\
$2$ & $0.966885096962700$ & $-0.874539361734654$ & $0.0923457352280467$\\
$3$ & $1.22069692821670$ & $-1.01305800766237$ & $0.207638920554325$\\
$4$ & $1.37558813187126$ & $-1.00679765237902$ & $0.368790479492242$\\
$5$ & $1.45826850019717$ & $-0.882725785735990$ & $0.575542714461177$\\
$10$ & $1.32380368369629$ & $\phantom{-}0.955535679496864$ & $2.27933936319316$\\
$20$ & $0.670652812798991$ & $\phantom{-}8.09862405929422$ & $8.76927687209322$\\
$40$ & $1.16405623139036$ & $\phantom{-}29.3133191923903$ & $30.4773754237807$\\
\bottomrule
\end{tabular}
\end{center}

Two independent audits were run.  First, $\lambda_k$ was recomputed as
$k\,[z^k]\log\xi(s(z))$ by a trapezoidal Cauchy integral on the circle $|z|=0.4$; for
$k\le20$ the worst relative discrepancy against the Stieltjes route is
$6.4\times10^{-54}$ at $60$ working digits.  Second, $Z_k$ alone was recomputed by the
same method applied to $\log\bigl((s-1)\zeta(s)\bigr)$ on $|z|=0.55$ at $120$ digits; the
discrepancies are $2.2\times10^{-110}$ at $k=40$ and $4.3\times10^{-107}$ at $k=50$.

For comparison, the zero-side representation
$\lambda_k=\sum_{\gamma>0}2\bigl(1-\cos k\theta_\gamma\bigr)$ truncated to the first $K$
ordinates has the relative deficits
\begin{center}
\small
\begin{tabular}{rcccc}
\toprule
$K$ & $k=1$ & $k=2$ & $k=4$ & $k=20$\\
\midrule
$300$  & $6.93\%$ & $6.93\%$ & $6.94\%$ & $7.30\%$\\
$600$  & $4.41\%$ & $4.41\%$ & $4.42\%$ & $4.64\%$\\
$900$  & $3.35\%$ & $3.35\%$ & $3.36\%$ & $3.53\%$\\
$1200$ & $2.75\%$ & $2.75\%$ & $2.75\%$ & $2.89\%$\\
$1500$ & $2.35\%$ & $2.35\%$ & $2.35\%$ & $2.47\%$\\
\bottomrule
\end{tabular}
\end{center}
The deficit is the truncated zero tail and decreases like $1/K$ up to logarithms, consistently with
the tail scale
\[
\sum_{\gamma>T}\frac{2}{4\gamma^2+1}.
\]
It is not an error of either computation.
\end{experiment}

\begin{remark}[Removal of circularity]
\label{rem:no_circularity}
Every numerical experiment in Sections~\ref{sec:two_generator},
\ref{sec:christoffel_bridge} and \ref{sec:computational_realization} that reports
positivity, pivots, determinants, $F_N$ or $\kappa_N$ uses the arithmetic $\lambda$ of
Experiment~\ref{exp:arithmetic_lambda}.  No zero of $\zeta$ is used as input there.  The
zero list is used only where a \emph{model} measure is deliberately constructed, namely in
the off-line detection experiments of
Subsection~\ref{subsec:detection}, where the atoms are the defining data of the model and
their arithmetic origin is irrelevant.
\end{remark}

\subsection{Exact Fourier--Laguerre transform of the Li test function}

For $r\in\mathbb R$, write
\begin{equation}
 \theta(r):=2\arctan\frac{1}{2r},
 \qquad
 H_k(r):=2\bigl(1-\cos(k\theta(r))\bigr).
 \label{eq:Hk_theta_definition}
\end{equation}
Use the Fourier convention
\begin{equation}
 \widehat f(u):=\frac{1}{2\pi}\int_{\mathbb R}f(r)\e^{-iru}\,\dd r.
 \label{eq:fourier_convention_laguerre}
\end{equation}

\begin{theorem}[Exact Fourier--Laguerre identity]
\label{thm:fourier_laguerre_identity}
For every $k\ge1$ and $u\in\mathbb R$,
\begin{equation}
 \boxed{
 \widehat H_k(u)
 =\e^{-|u|/2}L_{k-1}^{(1)}(|u|).}
 \label{eq:Hk_fourier_laguerre}
\end{equation}
\end{theorem}

\begin{proof}
For $|z|<1$, summing the elementary cosine series and using
\begin{equation}
 \cos\theta(r)=\frac{4r^2-1}{4r^2+1}
\end{equation}
gives
\begin{equation}
 \sum_{k\ge1}H_k(r)z^{k-1}
 =\frac{4(1+z)}{(1-z)\left(4r^2(1-z)^2+(1+z)^2\right)}.
 \label{eq:Hk_generating_r}
\end{equation}
Set
\begin{equation}
 a(z):=\frac{1+z}{2(1-z)}.
\end{equation}
Since
\begin{equation}
 \frac{1}{2\pi}\int_{\mathbb R}
 \frac{\e^{-iru}}{r^2+a^2}\,\dd r
 =\frac{1}{2a}\e^{-a|u|}
 \qquad(\Re a>0),
\end{equation}
Fourier transformation of \eqref{eq:Hk_generating_r} yields
\begin{equation}
 \sum_{k\ge1}\widehat H_k(u)z^{k-1}
 =\frac{\e^{-|u|/2}}{(1-z)^2}
 \exp\!\left(-\frac{|u|z}{1-z}\right).
 \label{eq:Hk_fourier_generating}
\end{equation}
Comparison with \eqref{eq:Laguerre_gen_alpha1} proves \eqref{eq:Hk_fourier_laguerre}.
\end{proof}

Thus the Cayley reflection system on the critical line and the associated Laguerre system on the logarithmic prime variable are Fourier dual.  If
\begin{equation}
 G_k(u):=\e^{-u/2}L_{k-1}^{(1)}(u),
 \qquad u>0,
 \label{eq:Laguerre_function_Gk}
\end{equation}
then the standard Laguerre orthogonality gives
\begin{equation}
 \int_0^\infty G_k(u)G_\ell(u)\,u\,\dd u
 =k\,\delta_{k\ell}.
 \label{eq:Laguerre_function_orthogonality}
\end{equation}

\subsection{Reflection cancellation of the prime main term}

Set $A_0=Z_0=0$ and define
\begin{align}
 (\mathsf A_N^E)_{mn}&:=A_{m+n}-A_{|m-n|},
 &
 (\mathsf Z_N^E)_{mn}&:=Z_{m+n}-Z_{|m-n|},
 \label{eq:E_arch_prime_matrices}\\
 (\mathsf A_N^O)_{mn}&:=A_{m+n-1}-A_{|m-n|},
 &
 (\mathsf Z_N^O)_{mn}&:=Z_{m+n-1}-Z_{|m-n|}.
 \label{eq:O_arch_prime_matrices}
\end{align}
Then, unconditionally,
\begin{equation}
 \E_N=\mathsf A_N^E+\mathsf Z_N^E,
 \qquad
 \Omat_N=\mathsf A_N^O+\mathsf Z_N^O.
 \label{eq:reflection_arch_prime_split}
\end{equation}

For the even family put
\begin{equation}
 \Phi_{mn}(u):=L_{m+n}(u)-L_{|m-n|}(u).
 \label{eq:Phi_mn_Laguerre_difference}
\end{equation}
Since $\Phi_{mn}(0)=0$, Proposition~\ref{prop:Zk_regularized_laguerre} gives the exact cancellation
\begin{equation}
 \boxed{
 (\mathsf Z_N^E)_{mn}
 =\int_{1^-}^{\infty}
 \Phi_{mn}'(\log t)\,\frac{\dd(\psi(t)-t)}{t}.}
 \label{eq:prime_matrix_PNT_error}
\end{equation}
No separate counterterm remains after reflection: the prime matrix pairs the Laguerre-difference kernel directly with the error $\psi(t)-t$ rather than with the main prime density.

The same construction applies to the odd family after replacing $m+n$ by $m+n-1$ in the first Laguerre index.  Moreover, because the displacement identities of Section~\ref{sec:displacement} are algebraic identities for reflection differences of a scalar sequence, each fixed-$u$ Laguerre kernel inherits the same rank-at-most-two displacement structure.  Likewise, the Chebyshev changes of basis of Section~\ref{sec:changes_of_basis} convert each such elementary reflection kernel into a Hankel matrix.  Hence the prime side is simultaneously a Laguerre superposition in the logarithmic variable and a low-displacement/Hankel object in the finite reflection coordinates.

There is an additional positivity property at the level of quadratic test functions.  For $c=(c_1,\dots,c_N)^T$ define
\begin{equation}
 \Phi_c(u):=\sum_{m,n=1}^Nc_mc_n\Phi_{mn}(u).
 \label{eq:Phi_c_definition}
\end{equation}
The trigonometric identity
\begin{equation}
 H_{m+n}(r)-H_{|m-n|}(r)
 =4\sin(m\theta(r))\sin(n\theta(r))
 \label{eq:H_reflection_square_identity}
\end{equation}
shows that
\begin{equation}
 h_c(r):=
 4\left|\sum_{m=1}^Nc_m\sin(m\theta(r))\right|^2\ge0.
 \label{eq:hc_nonnegative}
\end{equation}
By Theorem~\ref{thm:fourier_laguerre_identity}, its Fourier transform is
\begin{equation}
 q_c(u):=\widehat h_c(u)
 =-\e^{-|u|/2}\Phi_c'(|u|).
 \label{eq:qc_fourier_phi}
\end{equation}
Consequently $q_c$ is a continuous positive-definite function on $\mathbb R$.  In terms of this cone,
\begin{equation}
 \boxed{
 c^T\mathsf Z_N^Ec
 =-\int_{0^-}^{\infty}
 \e^{-u/2}q_c(u)\,\dd E(\e^u).}
 \label{eq:prime_quadratic_positive_definite_cone}
\end{equation}
Thus the arithmetic obstruction is not pairing the prime-number-theorem error against arbitrary Laguerre polynomials: it is pairing against the Fourier transforms of the nonnegative reflection squares \eqref{eq:hc_nonnegative}.

\subsection{First differences: a contractive matrix cosine family and an explicit archimedean lower decomposition}
\label{subsec:first_difference_prime_arch}

The even reflection sector admits a further normalization which makes the prime kernel uniformly bounded before the prime summation is performed.  Let
\begin{equation}
 D_N=
 \begin{pmatrix}
 1&&&\\
 -1&1&&\\
 &-1&1&\\
 &&\ddots&\ddots
 \end{pmatrix}
 \label{eq:DN_prime_arch}
\end{equation}
be the lower first-difference matrix.  For $u\ge0$ put
\begin{equation}
 K_0(u):=0,
 \qquad
 K_r(u):=L_{r-1}^{(1)}(u)\quad(r\ge1),
 \label{eq:Kr_prime_arch}
\end{equation}
and define the elementary Laguerre reflection matrix
\begin{equation}
 \mathsf M_N(u)
 :=\bigl(K_{m+n}(u)-K_{|m-n|}(u)\bigr)_{m,n=1}^N.
 \label{eq:MN_prime_arch}
\end{equation}
Finally set
\begin{equation}
 \mathsf C_N(u)
 :=\frac12\e^{-u/2}D_N\mathsf M_N(u)D_N^T.
 \label{eq:CN_prime_arch}
\end{equation}

\begin{proposition}[Contractive matrix-valued cosine representation]
\label{prop:CN_prime_arch_contraction}
For every $N$ there is a positive $N\times N$ matrix-valued measure $\Sigma_N$ on $\mathbb R$, with total mass $I_N$, such that
\begin{equation}
 \boxed{
 \mathsf C_N(u)=\int_{\mathbb R}\e^{-iru}\,\dd\Sigma_N(r).}
 \label{eq:CN_prime_arch_measure}
\end{equation}
Consequently
\begin{equation}
 \mathsf C_N(0)=I_N,
 \qquad
 \|\mathsf C_N(u)\|_2\le1
 \quad(u\in\mathbb R),
 \label{eq:CN_prime_arch_norm}
\end{equation}
and $u\mapsto\mathsf C_N(u)$ is matrix-valued positive definite.
\end{proposition}

\begin{proof}
Let
\begin{equation*}
 s(r):=(\sin\theta(r),\ldots,\sin(N\theta(r)))^T.
\end{equation*}
Theorem~\ref{thm:fourier_laguerre_identity} and
\eqref{eq:H_reflection_square_identity} imply
\begin{equation}
 \e^{-u/2}\mathsf M_N(u)
 =\frac1{2\pi}\int_{\mathbb R}
 4s(r)s(r)^T\e^{-iru}\,\dd r.
 \label{eq:MN_rankone_fourier}
\end{equation}
If
\begin{equation*}
 c(r):=
 \bigl(\cos(\theta(r)/2),\cos(3\theta(r)/2),\ldots,
 \cos((N-\tfrac12)\theta(r))\bigr)^T,
\end{equation*}
then
\begin{equation*}
 D_Ns(r)=2\sin(\theta(r)/2)c(r).
\end{equation*}
Hence \eqref{eq:CN_prime_arch_measure} holds with
\begin{equation}
 \dd\Sigma_N(r)
 =\frac4\pi\sin^2(\theta(r)/2)c(r)c(r)^T\,\dd r.
 \label{eq:SigmaN_prime_arch_density}
\end{equation}
This density is positive semidefinite.  Since
$L_{r-1}^{(1)}(0)=r$, one has
$\mathsf M_N(0)=(2\min(m,n))_{m,n}$, and the first-difference identity gives
$\mathsf C_N(0)=I_N$.  Thus $\Sigma_N(\mathbb R)=I_N$.
For a unit vector $v$, the scalar measure $v^T\Sigma_Nv$ is positive of mass $1$, so
$|v^T\mathsf C_N(u)v|\le1$.  Symmetry of $\mathsf C_N(u)$ yields
\eqref{eq:CN_prime_arch_norm}.
\end{proof}

\begin{corollary}[Prime term as a matrix prime--continuum discrepancy]
\label{cor:prime_continuum_matrix_discrepancy}
Put
\begin{equation}
 \widehat{\mathsf Z}_N^E
 :=\frac12D_N\mathsf Z_N^ED_N^T.
 \label{eq:Zhat_prime_arch}
\end{equation}
Then, entrywise and therefore in every norm on the fixed finite-dimensional matrix space,
\begin{equation}
 \boxed{
 \widehat{\mathsf Z}_N^E
 =\lim_{X\to\infty}
 \left[
 \int_0^{\log X}\e^{u/2}\mathsf C_N(u)\,\dd u
 -\sum_{n\le X}\frac{\Lambda(n)}{\sqrt n}
 \mathsf C_N(\log n)
 \right].}
 \label{eq:prime_continuum_matrix_discrepancy}
\end{equation}
Thus the first-difference prime matrix is a regularized discrepancy between the continuum measure and the prime powers, sampled against a single contractive matrix-valued cosine family.
\end{corollary}

\begin{proof}
Reflection in \eqref{eq:Zk_regularized_laguerre} gives
\begin{equation*}
 \mathsf Z_N^E
 =\lim_{X\to\infty}
 \left[
 \int_0^{\log X}\mathsf M_N(u)\,\dd u
 -\sum_{n\le X}\frac{\Lambda(n)}{n}\mathsf M_N(\log n)
 \right].
\end{equation*}
Apply $D_N/\sqrt2$ on the left and right and use
\eqref{eq:CN_prime_arch}.
\end{proof}

The archimedean term admits a complementary positive decomposition in the same coordinates.  Set $A_0:=0$ and
\begin{equation}
 \alpha_n:=A_{n+1}-2A_n+A_{n-1},
 \qquad n\ge1.
 \label{eq:alpha_prime_arch}
\end{equation}

\begin{proposition}[Positive moment formula for the archimedean second differences]
\label{prop:alpha_prime_arch}
For every $n\ge1$,
\begin{equation}
 \boxed{
 \alpha_n
 =\sum_{j=1}^{\infty}
 \frac1{(2j+1)^2}
 \left(\frac{2j}{2j+1}\right)^{n-1}>0.}
 \label{eq:alpha_prime_arch_series}
\end{equation}
\end{proposition}

\begin{proof}
From \eqref{eq:Ak_closed_form} and
$(1-2^{-r})\zeta(r)=\sum_{\substack{q\ge1\\ q\,\mathrm{odd}}}q^{-r}$, the centered second difference gives, for $n\ge2$,
\begin{align*}
 \alpha_n
 &=\sum_{\substack{q\ge1\\q\,\mathrm{odd}}}
 q^{-2}\sum_{r=2}^{n+1}
 \binom{n-1}{r-2}(-q^{-1})^{r-2}\\
 &=\sum_{\substack{q\ge1\\q\,\mathrm{odd}}}
 q^{-2}(1-q^{-1})^{n-1}.
\end{align*}
The $q=1$ term vanishes.  For $n=1$, the convention $A_0=0$ subtracts exactly that $q=1$ term, and the same formula results.  Writing $q=2j+1$ proves
\eqref{eq:alpha_prime_arch_series}.
\end{proof}

Define
\begin{equation}
 \widehat{\mathsf A}_N^E
 :=\frac12D_N\mathsf A_N^ED_N^T
 \label{eq:Ahat_prime_arch}
\end{equation}
and index its rows and columns by $0\le i,j\le N-1$.  A direct first-difference calculation gives
\begin{equation}
 (\widehat{\mathsf A}_N^E)_{ij}
 =A_1\delta_{ij}
 +\frac12(1-\delta_{ij})\alpha_{|i-j|}
 +\frac12\alpha_{i+j+1}.
 \label{eq:Ahat_alpha_toeplitz_hankel}
\end{equation}
Put
\begin{equation}
 \rho_j:=\frac{2j}{2j+1},
 \qquad
 w_j:=\frac1{(2j+1)^2},
 \label{eq:rho_w_prime_arch}
\end{equation}
and
\begin{equation}
 T_N(\rho):=(\rho^{|i-j|})_{i,j=0}^{N-1},
 \qquad
 v_N(\rho):=(1,\rho,\ldots,\rho^{N-1})^T.
 \label{eq:TN_vN_prime_arch}
\end{equation}

\begin{theorem}[Archimedean positive decomposition after first differences]
\label{thm:arch_prime_positive_decomposition}
For every $N$,
\begin{equation}
 \boxed{
 \widehat{\mathsf A}_N^E
 =c_\infty I_N
 +\frac12\sum_{j\ge1}\frac{w_j}{\rho_j}T_N(\rho_j)
 +\frac12\sum_{j\ge1}w_jv_N(\rho_j)v_N(\rho_j)^T,}
 \label{eq:arch_prime_positive_decomposition}
\end{equation}
where
\begin{equation}
 c_\infty
 :=\frac{1-\gamma_E-\log(2\pi)}2
 =-0.7075463656\ldots.
 \label{eq:c_infty_prime_arch}
\end{equation}
In particular
\begin{equation}
 \widehat{\mathsf A}_N^E\succeq c_\infty I_N.
 \label{eq:Ahat_prime_arch_lower}
\end{equation}
\end{theorem}

\begin{proof}
Proposition~\ref{prop:alpha_prime_arch} gives
\begin{equation*}
\alpha_n=\sum_{j\ge1}w_j\rho_j^{n-1}.
\end{equation*}
Thus the Toeplitz part of \eqref{eq:Ahat_alpha_toeplitz_hankel} is
\begin{equation*}
\frac12\sum_j\frac{w_j}{\rho_j}T_N(\rho_j),
\end{equation*}
and the Hankel part is
\begin{equation*}
\frac12\sum_jw_jv_N(\rho_j)v_N(\rho_j)^T.
\end{equation*}
The diagonal correction is
\begin{align*}
 A_1-\frac12\sum_{j\ge1}\frac{w_j}{\rho_j}
 &=1-\frac12(\gamma_E+\log4\pi)
 -\frac12\sum_{j\ge1}\frac1{2j(2j+1)}\\
 &=1-\frac12(\gamma_E+\log4\pi)-\frac12(1-\log2)
 =c_\infty.
\end{align*}
For $0<\rho<1$, $T_N(\rho)$ is positive definite: it is the Toeplitz moment matrix of the positive Poisson kernel
$(1-\rho^2)/(1-2\rho\cos\vartheta+\rho^2)$.  The remaining summands are rank-one positive semidefinite matrices, proving
\eqref{eq:Ahat_prime_arch_lower}.
\end{proof}

Define the explicit positive matrix
\begin{equation}
 \mathsf S_N
 :=\frac12\sum_{j\ge1}\frac{w_j}{\rho_j}T_N(\rho_j)
 +\frac12\sum_{j\ge1}w_jv_N(\rho_j)v_N(\rho_j)^T\succeq0.
 \label{eq:SN_prime_arch}
\end{equation}
Combining the two sides gives the exact normal form
\begin{equation}
 \boxed{
 \frac12D_N\E_ND_N^T
 =c_\infty I_N+\mathsf S_N+\widehat{\mathsf Z}_N^E.}
 \label{eq:even_prime_arch_normal_form}
\end{equation}
Since $D_N$ is invertible, positivity of the even reflection block is equivalent to
\begin{equation}
 \boxed{
 \mathsf S_N+\widehat{\mathsf Z}_N^E
 \succeq
 \frac{\gamma_E+\log(2\pi)-1}{2}\,I_N.}
 \label{eq:even_prime_arch_matrix_target}
\end{equation}
The archimedean matrix itself is not positive definite: already
$\mathsf A_1^E=(A_2)=(-0.8745393617\ldots)$.  Thus a normalization based on a positive square root of $\mathsf A_N^E$ is unavailable; \eqref{eq:even_prime_arch_normal_form} instead isolates its negative scalar part and an explicit positive semidefinite remainder.


\subsection{A scalar prime-only coefficient criterion}

The arithmetic factor also gives a direct coefficient formulation.  Define
\begin{equation}
 \mathcal F_{\rm ar}(z)
 :=\log\!\left[
 \left(\frac{1}{1-z}-1\right)
 \zeta\!\left(\frac{1}{1-z}\right)
 \right]
 =\mathcal F_{\rm ar}(0)
 +\sum_{k\ge1}\frac{Z_k}{k}z^k,
 \label{eq:Far_generating_Zk}
\end{equation}
where the logarithm denotes the branch analytic near $z=0$.  The Cayley relation is
\begin{equation}
 z=1-\frac1s=\frac{s-1}{s}.
 \label{eq:z_Cayley_again}
\end{equation}

\begin{theorem}[Prime--Laguerre root-growth criterion]
\label{thm:prime_laguerre_RH_criterion}
The following statements are equivalent:
\begin{enumerate}[label=(\roman*)]
\item $\RH$ holds;
\item $\mathcal F_{\rm ar}$ has an analytic logarithmic branch throughout $|z|<1$;
\item
\begin{equation}
 \boxed{
 \limsup_{k\to\infty}|Z_k|^{1/k}\le1;}
 \label{eq:Zk_root_growth_RH}
\end{equation}
\item for every $\varepsilon>0$ there is a constant $C_\varepsilon$ such that
\begin{equation}
 |Z_k|\le C_\varepsilon(1+\varepsilon)^k
 \qquad(k\ge1).
 \label{eq:Zk_subexponential_bound}
\end{equation}
\end{enumerate}
Equivalently, by \eqref{eq:Zk_PNT_error_laguerre}, $\RH$ is equivalent to subexponential root growth of the Laguerre transform of $\psi(t)-t$.
\end{theorem}

\begin{proof}
The map $z=(s-1)/s$ satisfies
\begin{equation}
 |z|<1
 \quad\Longleftrightarrow\quad
 |s-1|<|s|
 \quad\Longleftrightarrow\quad
 \Re s>\frac12.
 \label{eq:Cayley_disk_halfplane}
\end{equation}
The function $(s-1)\zeta(s)$ is analytic and nonzero at $s=1$.  On the simply connected disk $|z|<1$, it therefore admits an analytic logarithm exactly when it has no zero there.  By \eqref{eq:Cayley_disk_halfplane}, this is the absence of nontrivial zeros with $\Re s>1/2$; the functional equation then makes this equivalent to $\RH$.  Trivial zeros map outside the unit disk.  This proves (i)$\Longleftrightarrow$(ii).

By Cauchy--Hadamard, the radius \(R\) of the Taylor series in
\eqref{eq:Far_generating_Zk} satisfies
\begin{equation}
 R^{-1}
 =\limsup_{k\to\infty}
 \left|\frac{Z_k}{k}\right|^{1/k}
 =\limsup_{k\to\infty}|Z_k|^{1/k}.
 \label{eq:Zk_Cauchy_Hadamard}
\end{equation}
If (ii) holds, its Taylor series has radius at least one, so (iii) follows.  Conversely,
if (iii) holds, the series defines an analytic function \(G(z)\) on \(|z|<1\) which
agrees near \(0\) with the chosen logarithm of the arithmetic factor.  Hence
\(\exp G(z)\) agrees near \(0\), and therefore throughout the disk by the identity theorem,
with the analytic arithmetic factor.  In particular that factor has no zero in \(|z|<1\),
so \(G\) is its analytic logarithmic branch there and (ii) follows.  The equivalence of
(iii) and (iv) is the elementary root-test formulation of subexponential growth.
\end{proof}

The criterion admits a quantitative form.  If $R_Z$ denotes the radius of convergence in \eqref{eq:Far_generating_Zk}, then
\begin{equation}
 \boxed{
 R_Z
 =\min\left\{1,
 \inf_{\rho}\left|1-\frac1\rho\right|\right\},}
 \label{eq:Zk_radius_zero_geometry}
\end{equation}
where the infimum runs over the nontrivial zeros.  Indeed, their Cayley images are precisely the zeros of the analytic arithmetic factor, trivial zeros map outside $|z|=1$, and $z=1$ is itself a singular point because $s(z)\to\infty$ and $\log((s-1)\zeta(s))\sim\log s$.  In particular, a zero $\rho=\beta+i\gamma$ with $\beta>1/2$ satisfies
\begin{equation}
 \left|1-\frac1\rho\right|^2
 =\frac{(\beta-1)^2+\gamma^2}{\beta^2+\gamma^2}<1
 \label{eq:off_line_zero_inside_disk}
\end{equation}
and forces exponential coefficient growth at the reciprocal Cayley scale.

\begin{proposition}[Dominant off-line Cayley modes]
\label{prop:dominant_offline_modes}
Suppose $R_Z<1$, and let $z_1,\dots,z_M$ be the Cayley images of the nontrivial zeros lying on $|z|=R_Z$, with multiplicities $m_1,\dots,m_M$.  Then there exists $R_2>R_Z$ such that
\begin{equation}
 Z_k
 =-\sum_{j=1}^{M}m_jz_j^{-k}
 +O(R_2^{-k}).
 \label{eq:dominant_offline_mode_asymptotic}
\end{equation}
In particular,
\begin{equation}
 \limsup_{k\to\infty}|Z_k|^{1/k}=R_Z^{-1}>1.
 \label{eq:dominant_offline_root_growth}
\end{equation}
\end{proposition}

\begin{proof}
The closed disk $|z|\le r<1$ corresponds under $s=1/(1-z)$ to a compact subset of $\Re s>1/2$, so it contains only finitely many zero images.  Near a zero image $z_j$, the arithmetic generating function has the local form
\begin{equation}
 \mathcal F_{\rm ar}(z)
 =m_j\log(1-z/z_j)+\text{a function analytic near }z_j.
\end{equation}
Choose $R_2>R_Z$ smaller than the modulus of every other singularity.  Subtracting the finitely many logarithmic singularities on $|z|=R_Z$ leaves a function analytic in $|z|<R_2$.  Since
\begin{equation}
 k[z^k]\log(1-z/z_j)=-z_j^{-k},
\end{equation}
Cauchy estimates for the analytic remainder give \eqref{eq:dominant_offline_mode_asymptotic}.  The root-growth statement follows from the radius formula \eqref{eq:Zk_radius_zero_geometry}.
\end{proof}


\subsection{\texorpdfstring{The scalar $L^2$ route, its exact Bessel weight, and the return to the matrix cone}{The scalar L2 route, its exact Bessel weight, and the return to the matrix cone}}
\label{subsec:scalar_L2_barrier_matrix_return}

The scalar criterion above is a useful diagnostic reduction, but by itself it discards the reflection structure.  It is therefore important to record exactly what can and cannot be obtained from Laguerre orthogonality alone.

\begin{proposition}[Second cancellation and the $L^{(2)}$ representation]
\label{prop:Zk_second_cancellation_L2}
Let $E(t)=\psi(t)-t$.  For every $k\ge1$,
\begin{equation}
 \boxed{
 Z_k=-\int_1^\infty
 \frac{\psi(t)-t+1}{t^2}
 L_{k-1}^{(2)}(\log t)\,\dd t.}
 \label{eq:Zk_Laguerre_alpha2_ordinary}
\end{equation}
Equivalently, with $u=\log t$,
\begin{equation}
 Z_k=-\int_0^\infty
 \e^{-u}\bigl(E(\e^u)+1\bigr)L_{k-1}^{(2)}(u)\,\dd u.
 \label{eq:Zk_Laguerre_alpha2_u}
\end{equation}
\end{proposition}

\begin{proof}
Using \eqref{eq:Zk_PNT_error_laguerre} and \eqref{eq:Laguerre_derivative_identity}, first write
\begin{equation}
 Z_k=-\int_{0^-}^{\infty}\e^{-u}L_{k-1}^{(1)}(u)\,\dd E(\e^u).
 \label{eq:Zk_alpha1_before_second_cancellation}
\end{equation}
Then put
\begin{equation}
 R(u):=\e^{-u}E(\e^u).
\end{equation}
Then
\begin{equation}
 \e^{-u}\,\dd E(\e^u)=\dd R(u)+R(u)\,\dd u.
\end{equation}
The prime number theorem gives $R(u)\to0$ as $u\to\infty$, while $R(0)=E(1^-)=-1$.  Writing $n=k-1$ and integrating the $\dd R$ term by parts gives
\begin{equation}
 Z_k=-k-\int_0^\infty R(u)
 \left[L_n^{(1)}(u)-(L_n^{(1)})'(u)\right]\dd u.
\end{equation}
The Laguerre identities
\begin{equation}
 (L_n^{(1)})'=-L_{n-1}^{(2)},
 \qquad
 L_n^{(1)}=L_n^{(2)}-L_{n-1}^{(2)}
\end{equation}
give the exact cancellation
\begin{equation}
 L_n^{(1)}-(L_n^{(1)})'=L_n^{(2)}.
\end{equation}
Finally,
\begin{equation}
 \int_0^\infty \e^{-u}L_n^{(2)}(u)\,\dd u=n+1=k,
\end{equation}
which absorbs the boundary term $-k$ and proves \eqref{eq:Zk_Laguerre_alpha2_u}; the change of variables $t=\e^u$ gives \eqref{eq:Zk_Laguerre_alpha2_ordinary}.
\end{proof}

The ordinary integral \eqref{eq:Zk_Laguerre_alpha2_ordinary} removes the Stieltjes atoms, but it does not remove the critical-line difficulty.  Indeed, the $\alpha=2$ orthogonality relation is
\begin{equation}
 \int_0^\infty
 L_n^{(2)}(u)L_m^{(2)}(u)u^2\e^{-u}\,\dd u
 =(n+1)(n+2)\delta_{nm}.
 \label{eq:Laguerre_alpha2_orthogonality}
\end{equation}
Set
\begin{equation}
 f(u):=\frac{E(\e^u)+1}{u^2}.
 \label{eq:f_alpha2_definition}
\end{equation}
Then
\begin{equation}
 Z_{n+1}
 =-\left\langle f,L_n^{(2)}\right\rangle_{L^2(u^2\e^{-u}\dd u)}.
\end{equation}
Therefore Bessel's inequality has the direction and normalization
\begin{equation}
 \boxed{
 \sum_{n\ge0}
 \frac{|Z_{n+1}|^2}{(n+1)(n+2)}
 \le
 \int_0^\infty
 \frac{|E(\e^u)+1|^2}{u^2}\e^{-u}\,\dd u,}
 \label{eq:Bessel_correct_Zk}
\end{equation}
whenever the right-hand side is finite.  Equivalently,
\begin{equation}
 \boxed{
 \sum_{k\ge1}\frac{|Z_k|^2}{k(k+1)}
 \le
 \int_1^\infty
 \frac{|\psi(t)-t+1|^2}{t^2(\log t)^2}\,\dd t.}
 \label{eq:Bessel_correct_Zk_t}
\end{equation}
The polynomial factor is in the denominator, not the numerator.  Thus no estimate of the form
\begin{equation}
 \sum k(k+1)|Z_k|^2\le \|f\|_2^2
\end{equation}
follows from Bessel.  If the right-hand side of \eqref{eq:Bessel_correct_Zk_t} were finite, then one would obtain the much stronger individual estimate
\begin{equation}
 |Z_k|\le \sqrt{k(k+1)}\,\|f\|_2,
 \label{eq:Zk_polynomial_from_global_L2}
\end{equation}
and hence \eqref{eq:Zk_subexponential_bound}.

\begin{remark}[What the standard RH pointwise bound gives in the global $L^2$ norm]
\label{rem:global_L2_standard_RH_bound}
The standard pointwise consequence of $\RH$,
\begin{equation}
 E(t)=O(t^{1/2}\log^2 t),
\end{equation}
inserted into the right-hand side of \eqref{eq:Bessel_correct_Zk_t} gives only
\begin{equation}
 \frac{|E(t)+1|^2}{t^2(\log t)^2}
 =O\!\left(\frac{\log^2 t}{t}\right),
\end{equation}
whose integral diverges.  Thus the standard pointwise RH estimate does not establish the global $L^2$ integrability required on the right-hand side.  This calculation does not prove that RH itself forces divergence of that norm; it shows only that the standard pointwise bound is insufficient.  Any successful Laguerre argument along this line would therefore need additional cancellation, for example inside finite projections.
\end{remark}

\begin{warning}[Abel deformation does not bound the undeformed coefficients]
\label{warn:delta_deformation_not_bound}
For $\delta>0$ one may insert an extra factor $\e^{-\delta u}$ in \eqref{eq:Zk_Laguerre_alpha2_u} and obtain deformed coefficients $Z_k(\delta)$ for which global Bessel estimates are easier.  These coefficients form a different sequence.  Bounds for fixed $\delta>0$ do not imply bounds for $Z_k=Z_k(0)$ unless one also proves uniform control as $\delta\downarrow0$ on the $k$-dependent scale under consideration.  Thus Abel regularization alone does not cross the critical barrier.
\end{warning}

There is also an elementary one-sided sign implication.  If the coefficients $Z_k$ are eventually all nonnegative, or eventually all nonpositive, then Pringsheim's theorem applied after removing finitely many initial terms forces the radius of convergence of \eqref{eq:Far_generating_Zk} to be at least one, and hence $\RH$ follows from Theorem~\ref{thm:prime_laguerre_RH_criterion}.  This implication is logically correct but is not used as a proposed proof route.  High-precision numerical exploration through $k=80$ gives positive values but a pronounced oscillation about a baseline near $1$; representative values are
\begin{equation}
 Z_5\approx1.4583,\qquad
 Z_{20}\approx0.6707,\qquad
 Z_{50}\approx1.4979,\qquad
 Z_{75}\approx0.5882,\qquad
 Z_{80}\approx0.3576.
 \label{eq:Zk_numerical_oscillation_80}
\end{equation}
Two independent evaluations were used, the Stieltjes recursion of Proposition~\ref{prop:stieltjes_recursion} for $k\le50$ and a Cauchy integral on $\log\bigl((s-1)\zeta(s)\bigr)$ for $k\le80$, agreeing to $4.3\times10^{-107}$ at $k=50$.  All eighty values are strictly positive and lie in $[0.3576,1.5128]$, with the minimum attained at the right endpoint and the sequence still falling there.  The observed excursion amplitude does not show stabilization.  These data neither prove nor disprove eventual one-sidedness, but they provide no numerical basis for treating it as the main analytic mechanism.

\begin{remark}[Classical scalar mechanism versus the reflection compression]
\label{rem:scalar_classical_matrix_distinction}
The root-growth criterion \eqref{eq:Zk_root_growth_RH} is a factorwise version of the classical Keiper--Li singularity mechanism \cite{Keiper1992,Li1997,BombieriLagarias1999}: after the Cayley map, an off-line zero becomes a singularity inside the unit disk and therefore produces exponentially growing coefficients.  The role of the scalar sequence $Z_k$ in the present paper is consequently diagnostic.  The structure specific to the reflection construction is lost upon scalarization.

Before scalarization, the even arithmetic matrix satisfies \eqref{eq:prime_matrix_PNT_error}, and every quadratic test has the form
\begin{equation}
 c^T\mathsf Z_N^Ec
 =-\int_{0^-}^\infty \e^{-u/2}q_c(u)\,\dd E(\e^u),
\end{equation}
where $q_c$ is the Fourier transform of the nonnegative sine square \eqref{eq:hc_nonnegative} and is therefore positive-definite.  At the same time each fixed-$u$ prime atom inherits rank-at-most-two displacement and becomes Hankel after the Chebyshev change of basis.  The first-difference normal form further represents the prime side by the contractive matrix-valued cosine family \eqref{eq:CN_prime_arch}, while the archimedean side has the positive decomposition \eqref{eq:even_prime_arch_normal_form}.  These simultaneous constraints are absent from the scalar criterion.  They do not, however, weaken the logical content of the problem: positivity of all even reflection truncations is itself equivalent to $\RH$ by Theorem~\ref{thm:even_family_alone}.  The fact that the admissible quadratic tests form a smaller structured cone is therefore not by itself an analytic gain.  A matrix proof requires an additional mechanism that converts displacement, nesting, or the prime--continuum representation into a sign estimate; without such a mechanism the cone description is only a structural reformulation.
\end{remark}

\begin{remark}[Strategic status of the matrix route]
\label{rem:matrix_route_needs_mechanism}
The matrix formulation is retained not because it is logically weaker, but because it carries operations unavailable to a scalar sequence.  The most concrete non-circular target is to combine the nesting of the principal reflection blocks with their rank-two displacement structure.  For a positive definite truncation write the last Schur complement
\begin{equation}
 p_N:=\frac{\det \E_N}{\det \E_{N-1}}
 =e_{NN}-b_N^T\E_{N-1}^{-1}b_N,
 \label{eq:reflection_last_pivot}
\end{equation}
where $b_N=(e_{1N},\ldots,e_{N-1,N})^T$.  Positivity of all truncations is equivalent, away from singular boundary cases, to positivity of the successive pivots.  The research question is therefore whether the rank-two displacement identities can turn \eqref{eq:reflection_last_pivot} into a recurrence or comparison principle whose sign is visible directly on the prime--archimedean side.  This would be an actual mechanism; merely restating positivity on the reflection cone would not.
\end{remark}

The preceding discussion identifies the current boundary of this scalar estimate: the standard pointwise RH error bound does not by itself establish the global Laguerre $L^2$ norm required by Bessel.  This leaves finite-projection cancellation open and motivates the matrix route, where Schur complements, displacement generators, and critical-vector geometry provide additional exact structure on which a sign mechanism could act.

\subsection{Exact action of the Laguerre transform on exponential modes}

The preceding singularity statement has a direct integral counterpart.

\begin{lemma}[Laguerre--Laplace transform]
\label{lem:laguerre_laplace_transform}
For $\Re a>0$ and $k\ge1$,
\begin{equation}
 \boxed{
 \int_0^\infty \e^{-ax}L_{k-1}^{(1)}(x)\,\dd x
 =1-\left(1-\frac1a\right)^k.}
 \label{eq:Laguerre_Laplace_exact}
\end{equation}
\end{lemma}

\begin{proof}
Integrating the generating function \eqref{eq:Laguerre_gen_alpha1} termwise gives
\begin{align}
 \sum_{k\ge1}z^{k-1}
 \int_0^\infty \e^{-ax}L_{k-1}^{(1)}(x)\,\dd x
 &=\frac{1}{(1-z)^2}
 \frac{1}{a+z/(1-z)}\\
 &=\frac{1}{(1-z)(a-(a-1)z)}.
\end{align}
Coefficient extraction yields \eqref{eq:Laguerre_Laplace_exact}.
\end{proof}

For a nontrivial zero $\rho$ one may take $a=1-\rho$, since $0<\Re\rho<1$.  With
\begin{equation}
 z_\rho:=1-\frac1\rho
 \label{eq:zrho_laguerre}
\end{equation}
one obtains
\begin{equation}
 \int_0^\infty
 \e^{-(1-\rho)x}L_{k-1}^{(1)}(x)\,\dd x
 =1-z_\rho^{-k}.
 \label{eq:zero_mode_geometric_sequence}
\end{equation}
Thus the Laguerre transform sends the exponential mode generated by a zero directly to a geometric sequence in the Cayley coordinate.  The critical line is exactly the unit-modulus threshold $|z_\rho|=1$.

\begin{remark}[Correct large-$k$ formulation]
\label{rem:laguerre_asymptotic_route}
A large-$k$ proof of \eqref{eq:Zk_subexponential_bound} must be applied to the weighted Stieltjes form
\begin{equation}
 Z_k
 =-\int_{0^-}^{\infty}
 \e^{-u}L_{k-1}^{(1)}(u)\,\dd E(\e^u)
 =-\int_{0^-}^{\infty}
 \e^{-u/2}G_k(u)\,\dd E(\e^u),
 \label{eq:Zk_weighted_u_form}
\end{equation}
not to an unweighted integral in the variable $u$.  The standard turning-point structure also belongs to the weighted Laguerre functions rather than to a bare assertion that $L_k(k\xi)$ decays for $\xi<1$.  Finally, the prime-side Stieltjes representation and the zero-side explicit formula are alternative representations of the same functional; they cannot be added as independent contributions.  Equation \eqref{eq:zero_mode_geometric_sequence} identifies the exact phenomenon that any uniform asymptotic argument must suppress: a zero with $\Re\rho>1/2$ produces a geometric mode with modulus greater than one.
\end{remark}

The scalar identity therefore remains a precise diagnostic of exponentially growing Cayley modes, but the Bessel analysis in Subsection~\ref{subsec:scalar_L2_barrier_matrix_return} shows why scalar orthogonality alone does not supply the missing estimate.  The principal analytic target retained here is the reflection-matrix inequality \eqref{eq:even_prime_arch_matrix_target}, where the arithmetic discrepancy is tested only against the structured positive-definite cone \eqref{eq:prime_quadratic_positive_definite_cone} and the contractive cosine family \eqref{eq:CN_prime_arch}.



\section{Route A: the even reflection matrix and the exact RH target}
\label{sec:route_A_even}

The proof problem can now be stated without the odd matrix.  By
Theorem~\ref{thm:even_family_alone},
\begin{equation}
\boxed{
\RH
\Longleftrightarrow
\E_N\succeq0\ \forall N
\Longleftrightarrow
\E_N\succ0\ \forall N.}
\label{eq:routeA_even_equivalence}
\end{equation}
The role of \(\Omat_N\) in the preceding structural sections is comparative; it is not needed in
the implication to RH.

\subsection{A1. The exact Schur induction}
\label{subsec:routeA_schur}

Write
\begin{equation}
\E_N=
\begin{pmatrix}
\E_{N-1}&b_N\\
b_N^T&\lambda_{2N}
\end{pmatrix},
\qquad
p_N:=\lambda_{2N}-b_N^T\E_{N-1}^{-1}b_N.
\label{eq:routeA_block_pivot}
\end{equation}
Whenever \(\E_{N-1}\succ0\),
\begin{equation}
\boxed{
\E_N\succ0
\Longleftrightarrow
p_N>0.}
\label{eq:routeA_schur_equivalence}
\end{equation}
Therefore the single inductive implication
\begin{equation}
\boxed{
\E_{N-1}\succ0
\quad\Longrightarrow\quad
\lambda_{2N}>b_N^T\E_{N-1}^{-1}b_N
\qquad(N\ge2)}
\label{eq:routeA_inverse_schur_target}
\end{equation}
would prove \(\E_N\succ0\) for every \(N\), and hence RH.

\subsection{\texorpdfstring{A2. The local $LDL^T$ form}{A2. The local LDLT form}}
\label{subsec:routeA_local}

The recurrence part is already exact.  For the even boundary condition \(\ell_0=0\),
Theorem~\ref{thm:local_ldl_pivot_ratio} and Corollary~\ref{cor:defect_energy_pivot} give
\begin{equation}
\frac{p_{n+1}}{p_n}
=
1+(\tau_n-\tau_{n-1})
-\ell_n(\ell_n-\ell_{n-1})
=
V_n-V_{n-1}-\frac12(d_n-d_{n-1})^2.
\label{eq:routeA_local_identity}
\end{equation}
Hence
\begin{equation}
\boxed{
V_n-V_{n-1}>
\frac12(d_n-d_{n-1})^2
\quad\forall n}
\label{eq:routeA_local_target}
\end{equation}
would imply \(p_{n+1}/p_n>0\) for every \(n\).  Since \(p_1^E=\lambda_2>0\), induction would
give every even pivot positive and prove RH.  No odd pivot is required.

\subsection{A3. Exact prime--continuum scalarization of the matrix target}
\label{subsec:routeA_prime_scalar}

The first-difference normal form
\begin{equation}
\frac12D_N\E_ND_N^T
=
c_\infty I_N+\mathsf S_N+\widehat{\mathsf Z}_N^E
\label{eq:routeA_normal_form}
\end{equation}
gives a second exact formulation.  Fix a unit vector \(v\in\mathbb R^N\), put
\begin{equation}
c:=D_N^Tv,
\qquad
p_v(u):=\frac12c^T\mathsf M_N(u)c,
\qquad
g_v(u):=p_v'(u)-p_v(u),
\label{eq:routeA_pv_gv}
\end{equation}
and set
\begin{equation}
s_v:=v^T\mathsf S_Nv,
\qquad
P_v:=v^T\widehat{\mathsf Z}_N^Ev.
\label{eq:routeA_sv_Pv}
\end{equation}

\begin{proposition}[Exact signed PNT discrepancy]
\label{prop:routeA_signed_discrepancy}
For every fixed \(N\) and every unit \(v\),
\begin{equation}
p_v(0)=1
\label{eq:routeA_pv0}
\end{equation}
and
\begin{equation}
\boxed{
P_v
=
-1+
\int_1^\infty
(\psi(t)-t)\,
g_v(\log t)\,
\frac{\dd t}{t^2}.}
\label{eq:routeA_Pv_integrated}
\end{equation}
Consequently,
\begin{equation}
\boxed{
\E_N\succ0
\Longleftrightarrow
\int_1^\infty
(\psi(t)-t)\,
g_v(\log t)\,
\frac{\dd t}{t^2}
>
1-c_\infty-s_v
\quad\text{for every unit }v.}
\label{eq:routeA_signed_RH_target}
\end{equation}
\end{proposition}

\begin{proof}
The identity \(\mathsf M_N(0)=(2\min(m,n))_{m,n}\), together with
\(D_N(\min(m,n))D_N^T=I_N\), gives
\[
p_v(0)=c^T(\min(m,n))c=v^Tv=1.
\]
Since \(\Phi_{mn}'=-\mathsf M_{N,mn}\), the prime matrix formula
\eqref{eq:prime_matrix_PNT_error} yields
\[
P_v=-\int_{1^-}^{\infty}p_v(\log t)\,\frac{\dd(\psi(t)-t)}{t}.
\]
Integration by parts is legitimate for fixed \(N\): \(p_v\) is a polynomial in \(u\), while the
classical zero-free region for \(\zeta\) gives, for some \(c>0\),
\[
\psi(t)-t=O\!\left(t\exp(-c\sqrt{\log t})\right),
\]
which is more than sufficient to make the boundary term vanish and the resulting integral converge;
see, for example, \cite{Titchmarsh1986}.  The lower boundary is
\((\psi(1^-)-1)p_v(0)=-1\), and differentiation of \(p_v(\log t)/t\) gives
\(g_v(\log t)\dd t/t^2\).  This proves \eqref{eq:routeA_Pv_integrated}.  Substitution into
\[
v^T\left(\frac12D_N\E_ND_N^T\right)v
=c_\infty+s_v+P_v
\]
gives \eqref{eq:routeA_signed_RH_target}.
\end{proof}

\begin{remark}[Why magnitude bounds do not close the sign]
\label{rem:routeA_size_wrong_direction}
Equation~\eqref{eq:routeA_signed_RH_target} is a \emph{lower} bound on a signed oscillatory
integral.  A pointwise estimate for \(|\psi(t)-t|\), including the classical Schoenfeld bound under
RH \cite{Schoenfeld1976}, yields only an upper bound for the absolute value of that integral.
Accordingly, no argument that discards the sign correlation between \(\psi(t)-t\) and the structured
test \(g_v\) can prove \eqref{eq:routeA_signed_RH_target}.
\end{remark}

\begin{experiment}[Finite evaluation of the signed discrepancy and the magnitude barrier]
\label{exp:routeA_signed_discrepancy}
For the finite critical directions used in the multiprecision computation, the target values
\(1-c_\infty-s_v\) were
\[
1.67748,\ 1.68054,\ 1.68200,\ 1.6834,\ 1.6843,\ 1.6848
\]
for \(N=2,3,4,6,9,12\), approaching the observed level \(1.6860917\).
For \(N=2,3\), direct prime summation and piecewise integration up to \(X=10^7\) gave
\begin{center}
\small
\begin{tabular}{rcc}
\toprule
\(N\)&\(P_v+1\)&finite-\(X\) integral at \(X=10^7\)\\
\midrule
2&\(1.6774829902\)&\(1.67851\)\\
3&\(1.6805385806\)&\(1.68211\)\\
\bottomrule
\end{tabular}
\end{center}
with the finite-\(X\) values decreasing over the tested truncations.  No monotonicity or
asymptotic convergence claim is inferred from those finite values.

Using the Schoenfeld-type RH pointwise majorant only as an absolute-value diagnostic gives the
following much wider bounds:
\begin{center}
\small
\begin{tabular}{rccc}
\toprule
\(N\)&required lower level&absolute-value majorant&ratio\\
\midrule
2&1.67748&16.72&9.97\\
4&1.68200&39.54&23.5\\
9&1.68432&116.68&69.3\\
12&1.68477&173.66&103.1\\
\bottomrule
\end{tabular}
\end{center}
These are finite numerical diagnostics, not theorem statements.  Their analytic meaning is already
contained in Remark~\ref{rem:routeA_size_wrong_direction}: the missing information is signed
prime cancellation against the reflection-selected test.
\end{experiment}

\section{Route B: inverse dynamics and the adaptive Schur direction}
\label{sec:route_B_inverse}

This route keeps the same target \(p_N>0\) but treats \(\E_N^{-1}\) as the object that selects the
only new direction at each bordering step.  The purpose is not to estimate
\(\|\E_N^{-1}\|_2\).  Under RH the norm necessarily diverges at least at the rate proved in
Corollary~\ref{cor:inverse_norm_divergence}; by itself this divergence contains no sign information.  The useful information is the structured
\emph{directional} action of the inverse.

\subsection{B1. Two generators determine the inverse}
\label{subsec:routeB_generators}

Let \(\mathcal G_N:=\E_N^{-1}\), whenever the inverse exists, and write
\[
g_N:=\mathcal G_Ne_N,
\qquad
h_N:=\mathcal G_Nr_E^{(N)}.
\]
Corollary~\ref{cor:inverse_displacement} gives
\begin{equation}
\boxed{
\J_E^{(N)}\mathcal G_N-\mathcal G_N\J_E^{(N)}
=
\frac12(g_Nh_N^T-h_Ng_N^T).}
\label{eq:routeB_inverse_displacement}
\end{equation}
Thus the inverse has the same exact displacement rank two as the original reflection matrix.

Let
\[
\J_E^{(N)}
=S_N^{(J)}\diag(x_1,\ldots,x_N)(S_N^{(J)})^T,
\qquad
x_k=\cos\frac{k\pi}{N+1},
\]
be the Dirichlet sine diagonalization, and put hats on vectors and matrices in that basis.  For
\(i\ne j\), \eqref{eq:routeB_inverse_displacement} gives the explicit off-diagonal formula
\begin{equation}
\boxed{
\widehat{\mathcal G}_{ij}
=
\frac{\widehat g_i\widehat h_j-\widehat h_i\widehat g_j}
{2(x_i-x_j)}.}
\label{eq:routeB_inverse_offdiag}
\end{equation}
The diagonal is then fixed by \(g_N=\mathcal G_Ne_N\), because every component of
\((S_N^{(J)})^Te_N\) is nonzero.  Hence the pair \((g_N,h_N)\), together with the known
Dirichlet spectrum, determines the entire inverse.

\subsection{B2. Nesting is a rank-one inverse update}
\label{subsec:routeB_rankone}

With
\[
G:=\E_{N-1}^{-1},
\qquad
q_N:=Gb_N,
\]
the exact block inverse of Proposition~\ref{prop:closed_bordering} can be written
\begin{equation}
\boxed{
\E_N^{-1}
=
\begin{pmatrix}G&0\\0&0\end{pmatrix}
+
\frac1{p_N}
\binom{q_N}{-1}
\binom{q_N}{-1}^{\!T}.}
\label{eq:routeB_rankone_update}
\end{equation}
In particular,
\begin{equation}
\boxed{
(\E_N^{-1})_{NN}=\frac1{p_N}.}
\label{eq:routeB_corner_inverse}
\end{equation}
The size of this entry measures the inverse Schur margin, but its sign is what matters:
under the inductive hypothesis \(\E_{N-1}\succ0\),
\[
(\E_N^{-1})_{NN}>0
\Longleftrightarrow
p_N>0
\Longleftrightarrow
\E_N\succ0.
\]
This is an exact induction criterion, not a norm criterion.

\begin{experiment}[Arithmetic growth of a fixed inverse entry]
\label{exp:routeB_inverse11}
Using the arithmetic reconstruction of the Li coefficients and no zero ordinates as input, direct
multiprecision inversion gives
\begin{center}
\small
\begin{tabular}{rcc}
\toprule
\(N\)&\((\E_N^{-1})_{11}\)&\((\E_N^{-1})_{NN}=1/p_N\)\\
\midrule
2&\(1.1884763894\times10^{7}\)&\(2.9759641880\times10^{6}\)\\
3&\(5.6137194399\times10^{13}\)&\(2.2546793098\times10^{12}\)\\
4&\(9.0498470627\times10^{20}\)&\(4.6400962986\times10^{18}\)\\
5&\(2.6501319296\times10^{28}\)&\(1.5104983683\times10^{25}\)\\
6&\(1.0568460170\times10^{36}\)&\(6.1010537580\times10^{31}\)\\
7&\(8.9652480957\times10^{43}\)&\(4.9015837096\times10^{38}\)\\
8&\(1.3441381914\times10^{52}\)&\(6.6152967546\times10^{45}\)\\
9&\(1.8902274295\times10^{60}\)&\(8.0490977053\times10^{52}\)\\
\bottomrule
\end{tabular}
\end{center}
The first column gives the fixed-entry values.  No separate
``deficit'' notation is introduced here, because the symbol $d_n$ is already reserved for the local
$LDL^T$ defect variables.  The computation verifies rapid finite growth; the asymptotic divergence
itself is supplied independently by Corollary~\ref{cor:inverse_norm_divergence} and is not inferred
from this table.
\end{experiment}

\subsection{B3. The inverse selects the residual polynomial}
\label{subsec:routeB_residual}

Define
\begin{equation}
\chi_N:=
\binom{-q_N}{1}.
\label{eq:routeB_chiN}
\end{equation}
Then
\begin{equation}
\boxed{
\E_N\chi_N=p_Ne_N,
\qquad
\chi_N^T\E_N\chi_N=p_N,
\qquad
\E_N^{-1}e_N=\frac{\chi_N}{p_N}.}
\label{eq:routeB_residual_identities}
\end{equation}
Thus \(\chi_N\) is the unique Schur direction orthogonal, in the Gram geometry, to all preceding
modes.

In the Chebyshev basis
\(u_N(x)=(U_0(x),\ldots,U_{N-1}(x))^T\), set
\begin{equation}
R_{N-1}(x):=\chi_N^Tu_N(x).
\label{eq:routeB_residual_polynomial}
\end{equation}
Under RH this is the residual obtained by projecting \(U_{N-1}\) onto
\(\Span\{U_0,\ldots,U_{N-2}\}\), and
\begin{equation}
\boxed{
\|R_{N-1}\|_{L^2(\mu_E)}^2=p_N.}
\label{eq:routeB_residual_energy}
\end{equation}
Moreover,
\[
u_N(x)^T\E_N^{-1}e_N
=
\frac{R_{N-1}(x)}{p_N}.
\]

\subsection{B4. Christoffel--Darboux kernel and leverage}
\label{subsec:routeB_CD}

Theorem~\ref{thm:inverse_CD} identifies
\begin{equation}
\mathcal K_N^E(x,y)
=
u_N(x)^T\E_N^{-1}u_N(y)
\label{eq:routeB_CD_kernel}
\end{equation}
with the degree-\(<N\) reproducing kernel.  Equation~\eqref{eq:routeB_rankone_update} becomes
\begin{equation}
\boxed{
\mathcal K_N^E(x,y)
=
\mathcal K_{N-1}^E(x,y)
+
\frac{R_{N-1}(x)R_{N-1}(y)}{p_N}.}
\label{eq:routeB_CD_update}
\end{equation}
This is the kernel form of the same Schur step.

For a critical-line zero,
\[
a_{N,\gamma}^E
=
2\sin\theta_\gamma\,u_N(x_\gamma),
\]
hence
\begin{equation}
\boxed{
\ell_{N,\gamma}^E
=
4\sin^2\theta_\gamma\,
\mathcal K_N^E(x_\gamma,x_\gamma).}
\label{eq:routeB_leverage_CD}
\end{equation}
At the accumulation point \(x=1\),
\[
u_N(1)=(1,2,\ldots,N)^T,
\]
so the high-zero constant in Theorem~\ref{thm:high_zero_leverage_decay} is exactly
\begin{equation}
c_N^E=\mathcal K_N^E(1,1).
\label{eq:routeB_endpoint_kernel}
\end{equation}

The same formula complexifies.  For \(q\in\mathbb C\setminus\{0,\pm1\}\), put
\(x(q)=(q+q^{-1})/2\).  Since
\[
q^m-q^{-m}=(q-q^{-1})U_{m-1}(x(q)),
\]
the zero-orbit vector
\[
v_N(q):=(q-q^{-1},q^2-q^{-2},\ldots,q^N-q^{-N})^T
\]
satisfies
\begin{equation}
\boxed{
v_N(q)^T\E_N^{-1}v_N(r)
=
(q-q^{-1})(r-r^{-1})
\mathcal K_N^E(x(q),x(r)).}
\label{eq:routeB_complex_kernel}
\end{equation}
Thus every rank-two secular determinant associated with a hypothetical off-line orbit can be
written using only complex evaluations of the same inverse kernel.

\subsection{B5. The adaptive prime--continuum inequality}
\label{subsec:routeB_adaptive_prime}

The inverse becomes potentially useful for RH because it selects one structured test rather than all
vectors.  Define \(\chi_N\) by \eqref{eq:routeB_chiN} and
\begin{equation}
y_N:=D_N^{-T}\chi_N.
\label{eq:routeB_yN}
\end{equation}
The first-difference normal form gives exactly
\begin{equation}
\boxed{
\frac{p_N}{2}
=
y_N^T
\left(
c_\infty I_N+\mathsf S_N+\widehat{\mathsf Z}_N^E
\right)y_N.}
\label{eq:routeB_pivot_normal_form}
\end{equation}
Set
\[
a:=-c_\infty
=\frac{\gamma_E+\log(2\pi)-1}{2}.
\]
Then the single inequality
\begin{equation}
\boxed{
y_N^T(\mathsf S_N+\widehat{\mathsf Z}_N^E)y_N
>
a\|y_N\|^2
\qquad(N\ge2)}
\label{eq:routeB_adaptive_target}
\end{equation}
under the inductive hypothesis \(\E_{N-1}\succ0\) implies \(p_N>0\), hence propagates
positivity and proves RH.

The arithmetic term is tested against
\begin{equation}
\phi_N(u):=y_N^T\mathsf C_N(u)y_N.
\label{eq:routeB_phiN}
\end{equation}
Since \(\mathsf C_N\) is a contractive matrix-valued positive-definite cosine family,
\begin{equation}
\phi_N(0)=\|y_N\|^2,
\qquad
|\phi_N(u)|\le\|y_N\|^2.
\label{eq:routeB_phi_constraints}
\end{equation}
The exact prime--continuum formula yields
\begin{equation}
y_N^T\widehat{\mathsf Z}_N^Ey_N
=
\lim_{X\to\infty}
\left[
\int_0^{\log X}\e^{u/2}\phi_N(u)\,\dd u
-
\sum_{n\le X}\frac{\Lambda(n)}{\sqrt n}\phi_N(\log n)
\right].
\label{eq:routeB_adaptive_discrepancy}
\end{equation}
Unlike the full matrix inequality
\(\mathsf S_N+\widehat{\mathsf Z}_N^E\succeq aI\), the induction needs
\eqref{eq:routeB_adaptive_target} only on the recursively generated sequence \(y_N\).
This is the precise open target of the inverse route.

\subsection{B6. Finite numerical status and shift nonmonotonicity}
\label{subsec:routeB_numerics}

The following numerical statements are recorded at their verified finite level:
\begin{itemize}
\item arithmetic Li reconstruction gives positive even reflection pivots through \(N=22\); the
reported values include
\[
p_2=3.360255490\times10^{-7},\quad
p_{12}=1.214721178\times10^{-75},\quad
p_{22}=2.854130423\times10^{-154};
\]
\item for the first-difference congruence
\(\widehat E_N=\tfrac12D_N\E_ND_N^T\), the computed smallest eigenvalues at
\(N=4,8,12,16,20\) are respectively
\begin{align*}
N=4:&\quad 8.19930569\times10^{-22},&
N=8:&\quad 2.84288630\times10^{-53},\\
N=12:&\quad 6.68407776\times10^{-87},&
N=16:&\quad 3.22962551\times10^{-122},\\
N=20:&\quad 1.30149539\times10^{-158}.&&
\end{align*}
all positive in this finite computation; these are eigenvalues of \(\widehat E_N\), not of
\(\E_N\) itself;
\item the Christoffel pivot ratio \(\kappa_N=p_N^E/p_N^O\) is numerically inside the interval
\([4-\lambda_1,4]\) through the tested orders; it begins near \(3.99650\) and moves toward \(4\) over that finite range;
\item zero leverage satisfies the exact identities of Section~\ref{sec:zero_leverage}.  At \(N=10\),
the first five positive zeros carry about \(49.6\%\) of total leverage, the first ten about
\(67.3\%\), and the first two hundred about \(92.7\%\);
\item the arithmetic coefficients \(Z_k\), computed through \(k=80\), are all positive in the tested
range.  Their observed range is \([0.3576,1.5128]\); see
Section~\ref{sec:arch_prime_laguerre}.
\end{itemize}

The shift-index sequence is not monotone even at fixed matrix order.  For \(k\ge0\), define the shifted even block
\begin{equation}
\E_N^{[k]}:=
\bigl(\lambda_{2k+m+n}-\lambda_{|m-n|}\bigr)_{m,n=1}^N.
\label{eq:shifted_even_block}
\end{equation}
Under RH it has the exact Gram representation
\begin{equation}
(\E_N^{[k]})_{mn}
=4\sum_{\gamma>0}\sin((k+m)\theta_\gamma)\sin((k+n)\theta_\gamma).
\label{eq:shifted_even_gram}
\end{equation}
Whenever the block is invertible, set
\begin{equation}
\delta_k^{(N)}:=\frac{1}{((\E_N^{[k]})^{-1})_{11}},
\qquad
R_k^{(N)}:=\lambda_{2k+2}-\delta_k^{(N)}.
\label{eq:shifted_delta_R}
\end{equation}
Under RH, \(\delta_k^{(N)}\) is the squared residual norm of the mode \(k+1\) after projection
onto the modes \(k+2,\ldots,k+N\).  Therefore, for each fixed \(k\),
\begin{equation}
\delta_k^{(N+1)}\le\delta_k^{(N)},
\qquad
R_k^{(N+1)}\ge R_k^{(N)},
\qquad
0\le R_k^{(N)}\le\lambda_{2k+2}.
\label{eq:shifted_N_monotonicity}
\end{equation}
These finite projection inequalities do \emph{not} by themselves identify the limit with
\(\lambda_{2k+2}\); that equality would require the corresponding tail-completeness statement.

By contrast, monotonicity in the shift \(k\) is false even at \(N=2\).  Here
\begin{equation}
\delta_k^{(2)}
=\lambda_{2k+2}
-\frac{(\lambda_{2k+3}-\lambda_1)^2}{\lambda_{2k+4}}.
\label{eq:shifted_delta_N2_closed}
\end{equation}
Arithmetic reconstruction of the Li coefficients gives
\begin{center}
\small
\begin{tabular}{c@{\qquad}c@{\qquad}c}
\toprule
$k$ & $\delta_k^{(2)}$ & $\delta_k^{(2)}-\delta_{k-1}^{(2)}$\\
\midrule
40 & $0.02985633140654687$ & $+5.782\times10^{-4}$\\
41 & $0.03028625109019011$ & $+4.299\times10^{-4}$\\
42 & $0.03056677743905413$ & $+2.805\times10^{-4}$\\
43 & $0.03069982127912338$ & $+1.330\times10^{-4}$\\
44 & $0.03069030817860707$ & $-9.513\times10^{-6}$\\
45 & $0.03054607923982695$ & $-1.442\times10^{-4}$\\
46 & $0.03027771991886325$ & $-2.684\times10^{-4}$\\
47 & $0.02989832437134847$ & $-3.794\times10^{-4}$\\
48 & $0.02942320442607452$ & $-4.751\times10^{-4}$\\
\bottomrule
\end{tabular}
\end{center}
Thus the finite sequence rises through \(k=43\) and decreases from \(k=44\) onward.  No
proof route in this paper uses monotonicity in the shift index.

\section{What would prove RH in the present framework}
\label{sec:what_proves_RH}

For clarity, the proof obligations supplied by the preceding sections are listed below; in the local line, \(V_n,d_n\) are the even-family quantities (boundary parameter \(\sigma=0\)):
\begin{align}
\text{local route:}\qquad&
V_n-V_{n-1}>\frac12(d_n-d_{n-1})^2
\quad\forall n,
\label{eq:final_RH_local}\\
\text{Schur/inverse route:}\qquad&
\lambda_{2N}>
b_N^T\E_{N-1}^{-1}b_N
\quad\forall N\ge2,
\label{eq:final_RH_schur}\\
\text{adaptive prime route:}\qquad&
y_N^T(\mathsf S_N+\widehat{\mathsf Z}_N^E)y_N
>
a\|y_N\|^2
\quad\forall N\ge2.
\label{eq:final_RH_prime}
\end{align}
Each one, proved under the inductive positivity of the preceding even block, implies all
\(p_N^E>0\), hence all \(\E_N\succ0\), hence RH by
Theorem~\ref{thm:even_family_alone}.  The odd family does not occur in this implication chain.

The de Bruijn--Newman filter of Section~\ref{sec:newman_filter} further constrains the form of such
a proof.  Reflection, displacement, Schur, and inverse identities alone cannot provide the sign,
because the same algebra survives negative heat time while positivity fails.  Within the present
coordinate system, the zeta-specific ingredient that does not survive is the Euler-product/prime
representation \eqref{eq:prime_continuum_matrix_discrepancy}.  Hence the natural remaining task is
to combine the adaptive inverse direction \(y_N\) with prime cancellation strongly enough to prove
\eqref{eq:final_RH_prime}.

\section{Conclusion}
\label{sec:conclusion}

The input kernel is Lagarias's published identity
\[
\langle G_n,G_m\rangle_W=\lambda_n+\lambda_{-m}-\lambda_{n-m},
\]
taken from Lagarias~\cite[Section~3, Theorem~3.1]{Lagarias2007}; see the exact source ledger in
Remark~\ref{rem:status_provenance}.  The matrices \(\E_N\) and \(\Omat_N\) are obtained
from that kernel by the explicit antisymmetric basis changes recorded in Remark~\ref{rem:status_provenance}.
The finite reflection transformation and the identities derived from it are the subject of this paper.  The even matrix
alone is sufficient for RH:
\[
\RH\Longleftrightarrow \E_N\succeq0\ \forall N
\Longleftrightarrow \E_N\succ0\ \forall N.
\]
The odd family remains structurally useful but is not required in the proof chain.

For \(\E_N\), the algebraic part of the matrix route is now closed to a single sign step.  Nesting,
rank-two displacement, the column wave equation, and the closed two-generator scheme give an exact
Schur recurrence; the local \(LDL^T\) identity reduces the next pivot to a \(2\times2\) cell.
Independently, the prime--archimedean decomposition converts the same positivity question into a
signed prime-number-theorem discrepancy against explicit reflection-selected tests.

The inverse route identifies a sharper directional formulation.  The inverse retains displacement
rank two, is updated by rank one under nesting, and is the coefficient matrix of the
Christoffel--Darboux kernel.  Its last column selects the Schur residual polynomial and therefore
the unique new direction that needs to be controlled.  In first-difference coordinates this
produces the adaptive test \(y_N\) and the sufficient inequality
\[
y_N^T(\mathsf S_N+\widehat{\mathsf Z}_N^E)y_N>a\|y_N\|^2.
\]
Proving that inequality for every \(N\), or equivalently proving every even Schur pivot positive,
would prove RH.

The de Bruijn--Newman deformation shows why the missing mechanism cannot be purely structural.
Negative heat time preserves the coefficient/reflection algebra while leaving the all-real-zero
region.  Therefore a successful proof must use information not invariant under that deformation.
In the present paper that information is carried by the Euler-product/prime side of the
prime--continuum discrepancy.  This isolates the next analytic task without replacing it by a
finite numerical observation.

\appendix
\section{Lateral questions}
\label{sec:lateral_questions}

The following problems are retained because they may reveal additional structure or useful
asymptotics, but they are not steps in the logical chain \eqref{eq:final_RH_local}--\eqref{eq:final_RH_prime}.  In
particular, failure to solve any of them does not block the local pivot route described in
Section~\ref{sec:route_A_even}.

\subsection{Rate of the Christoffel factor}
\label{ques:F_rate}
By Theorem~\ref{thm:margin_band} the sequence $F_N$ decreases to
$\det(I-\tfrac12\mathcal D)=\prod_\gamma 4\gamma^2/(4\gamma^2+1)$.  Determine the rate.
The numerical evidence of Experiment~\ref{exp:christoffel_margin} suggests that each
Jacobi level resolves approximately one further Cayley angle, which would give
$F_N-\det(I-\tfrac12\mathcal D)\asymp\sum_{j>N}(4\gamma_j^2+1)^{-1}\asymp N^{-1}\log N$
after the Riemann--von Mangoldt count.  A proof would convert the band
\eqref{eq:margin_band} into a quantitative statement and, by
Corollary~\ref{cor:kappa_le_4}, into an effective test.

\subsection{Sharpness of the detection threshold}
\label{ques:detection_threshold}
Experiment~\ref{exp:detection_index} finds $N^*\theta_\gamma$ between approximately
$0.48$ and $0.56$ over the tested range, with $N^*$ increasing as $\beta$ approaches
$1/2$ in the reported data.  Determine whether this product has a limit as
$\beta\uparrow1/2$ and $\gamma\to\infty$, and, if so, identify it.  A natural guess is that
such a limit is determined by the first zero of a limiting sine-polynomial extremal
problem associated with \eqref{eq:critical_sine_energy}.

\subsection{Prime--archimedean reflection inequality}
\label{ques:prime_laguerre_subexponential}
Establish the matrix inequality
\begin{equation}
 \mathsf S_N+\widehat{\mathsf Z}_N^E
 \succeq
 \frac{\gamma_E+\log(2\pi)-1}{2}\,I_N
 \qquad(N\ge1)
\end{equation}
directly from the prime--continuum representation
\eqref{eq:prime_continuum_matrix_discrepancy}.  This is an alternative equivalent form of
$\RH$, not a prerequisite for the local route.  The scalar estimate
\eqref{eq:Zk_subexponential_bound} remains an equivalent diagnostic, while
Remark~\ref{rem:global_L2_standard_RH_bound} records only the insufficiency of the standard
RH pointwise estimate for the global Bessel norm.

\subsection{Critical sine-polynomial flatness}
\label{ques:critical_vector_scaling}
Let $v^{(N)}$ be the normalized smallest-eigenvalue direction for the even reflection
problem, with a consistent sign convention, and let $P_N=P_{v^{(N)}}$ be defined by
\eqref{eq:critical_sine_polynomial}.  Determine the largest scale $r=r(N)$ for which the
odd moments
\begin{equation}
 \sum_{m=1}^N m^{2j+1}v_m^{(N)},
 \qquad 0\le j<r,
 \label{eq:critical_flatness_open_moments}
\end{equation}
are asymptotically negligible in a normalization intrinsic to the eigenvalue problem.
Equivalently, determine the growth of the order of flatness of $P_N$ at $\theta=0$.  The
present computations show low-index concentration but do not establish saturation at the
scale $\theta_1^{-1}$; that scale remains a numerical comparison until larger-$N$ moment
data are available.

\subsection{Joint leverage profile}
\label{ques:joint_leverage_profile}
For fixed $N$, Theorem~\ref{thm:high_zero_leverage_decay} gives a summable high-zero tail.
Determine the scaling limit of
\begin{equation}
 \ell_{N,\gamma_j}^{E,O}
 \label{eq:joint_leverage_question}
\end{equation}
when $N\to\infty$ and $j=j(N)$ grows simultaneously.  Determine whether the canonical
Cauchy--Binet ensemble has a macroscopic selection profile in the zero index and whether
its transition region is governed only by the Riemann--von Mangoldt density or also by
finer zero statistics.

\subsection{Asymptotics of the smallest eigenvalue}
\label{ques:min_eigenvalue}
Determine the asymptotic behavior of
\begin{equation}
 \sigma_{\min}(\E_N),
 \qquad
 \sigma_{\min}(\Omat_N)
 \label{eq:min_eigenvalue_question}
\end{equation}
under $\RH$, and relate it to approximation of the zero-angle support by finite sine
polynomials.

\subsection{Asymptotics of determinant ratios}
\label{ques:determinant_ratios}
Determine asymptotics for
\begin{equation}
 \frac{\Delta_N^E}{\Delta_{N-1}^E},
 \qquad
 \frac{\Delta_N^O}{\Delta_{N-1}^O}.
 \label{eq:determinant_ratio_question}
\end{equation}
These are residual Gram energies and squared leading Cholesky coefficients.

\subsection{Jacobi parameters}
\label{ques:jacobi_parameters}
Find direct formulas or asymptotics for
\begin{equation}
 \{a_n^E,b_n^E\}_{n\ge0},
 \qquad
 \{a_n^O,b_n^O\}_{n\ge0}
 \label{eq:jacobi_parameter_question}
\end{equation}
in terms of Li coefficients or classical zeta data, and quantify the
Christoffel/Darboux relation induced by \eqref{eq:christoffel_measure_relation}.

\subsection{Inverse-entry structure}
\label{ques:inverse_entries}
Determine whether the inverse entries possess asymptotic sign, decay, or localization
patterns that are not forced by generic positive Gram matrices.  The rank-at-most-two
inverse displacement identity \eqref{eq:E_inverse_displacement}--\eqref{eq:O_inverse_displacement}
provides a starting point.

\subsection{Automorphic reflection systems below the universal exponent}
\label{ques:automorphic_extension}
Lagarias defines generalized Li coefficients for principal automorphic $L$-functions
\cite{Lagarias2007}.  Determine which functional-equation symmetries are sufficient to
reproduce the pair of reflection matrices, their sine/half-sine Gram systems, and the
Christoffel relation between the two associated measures.  Once such a positive Cayley
reflection system is available under GRH, Corollary~\ref{cor:universal_cayley_exponent}
fixes its leading determinant coefficient at $-2$; the genuinely automorphic information
must therefore be sought in lower-order terms, leverage profiles, or other matrix
invariants.

\begin{thebibliography}{99}

\bibitem{BergChristensenRessel1984}
C. Berg, J. P. R. Christensen and P. Ressel,
\emph{Harmonic Analysis on Semigroups: Theory of Positive Definite and Related Functions},
Graduate Texts in Mathematics \textbf{100}, Springer, New York, 1984.
DOI: 10.1007/978-1-4612-1128-0.

\bibitem{Bombieri2010}
E. Bombieri,
\emph{The classical theory of zeta and $L$-functions},
Milan J. Math. \textbf{78} (2010), 11--59.
DOI: 10.1007/s00032-010-0121-8.

\bibitem{BombieriLagarias1999}
E. Bombieri and J. C. Lagarias,
\emph{Complements to Li's criterion for the Riemann hypothesis},
J. Number Theory \textbf{77} (1999), no. 2, 274--287.
DOI: 10.1006/jnth.1999.2392.


\bibitem{Bombieri2000}
E. Bombieri,
\emph{Remarks on Weil's quadratic functional in the theory of prime numbers, I},
Rend. Lincei Mat. Appl. (9) \textbf{11} (2000), no. 3, 183--233.


\bibitem{BucurErnvallOdzakSmajlovic2016}
A. Bucur, A.-M. Ernvall-Hyt\"onen, A. Od\v{z}ak and L. Smajlovi\'c,
\emph{On a Li-type criterion for zero-free regions of certain Dirichlet series with real coefficients},
LMS J. Comput. Math. \textbf{19} (2016), no. 1, 259--280.
DOI: 10.1112/S1461157016000115.

\bibitem{Droll2012}
A. Droll,
\emph{Variations of Li's criterion for an extension of the Selberg class},
Ph.D. thesis, Queen's University, Kingston, Ontario, 2012.

\bibitem{Freitas2006}
P. Freitas,
\emph{A Li-type criterion for zero-free half-planes of Riemann's zeta function},
J. London Math. Soc. (2) \textbf{73} (2006), no. 2, 399--414.
DOI: 10.1112/S0024610706022599.

\bibitem{OmarOuniMazhouda2014}
S. Omar, R. Ouni and K. Mazhouda,
\emph{On the Li coefficients for the Hecke $L$-functions},
Math. Phys. Anal. Geom. \textbf{17} (2014), nos. 1--2, 67--81.
DOI: 10.1007/s11040-014-9141-x.

\bibitem{OmarOuni2017}
S. Omar and R. Ouni,
\emph{A note on Li's coefficients for Hecke $L$-functions},
Int. J. Number Theory \textbf{13} (2017), no. 7, 1747--1753.
DOI: 10.1142/S1793042117501007.

\bibitem{Palojarvi2018}
N. Paloj\"arvi,
\emph{Explicit zero-free regions and a $\tau$-Li-type criterion},
Albanian J. Math. \textbf{14} (2020), no. 1, 47--77.
DOI: 10.51286/albjm/1608313766; arXiv:1807.01506.

\bibitem{Sekatskii2014}
S. K. Sekatskii,
\emph{Generalized Bombieri--Lagarias' theorem and generalized Li's criterion with its arithmetic interpretation},
Ukrainian Math. J. \textbf{66} (2014), no. 3, 415--431.
DOI: 10.1007/s11253-014-0940-9.

\bibitem{Chihara1978}
T. S. Chihara,
\emph{An Introduction to Orthogonal Polynomials},
Gordon and Breach, New York, 1978.

\bibitem{Coffey2007ThetaLaguerre}
M. W. Coffey,
\emph{The theta--Laguerre calculus formulation of the Li/Keiper constants},
J. Approx. Theory \textbf{146} (2007), no. 2, 267--275.
DOI: 10.1016/j.jat.2006.10.006.

\bibitem{Coffey2010Polygamma}
M. W. Coffey,
\emph{Polygamma theory, the Li/Keiper constants and the Li criterion for the Riemann hypothesis},
Rocky Mountain J. Math. \textbf{40} (2010), no. 6, 1841--1862.
DOI: 10.1216/RMJ-2010-40-6-1841.

\bibitem{Coffey2010ZeroSums}
M. W. Coffey,
\emph{On certain sums over the non-trivial zeta zeroes},
Proc. R. Soc. A \textbf{466} (2010), no. 2124, 3679--3692.
DOI: 10.1098/rspa.2010.0064.

\bibitem{Conrey2003}
J. B. Conrey,
\emph{The Riemann hypothesis},
Notices Amer. Math. Soc. \textbf{50} (2003), 341--353.

\bibitem{Edwards1974}
H. M. Edwards,
\emph{Riemann's Zeta Function},
Academic Press, New York, 1974.

\bibitem{HoughKrishnapurPeresVirag2009}
J. B. Hough, M. Krishnapur, Y. Peres and B. Vir\'ag,
\emph{Zeros of Gaussian Analytic Functions and Determinantal Point Processes},
University Lecture Series, vol. 51, American Mathematical Society, Providence, RI, 2009.

\bibitem{Hilberdink2006}
T. Hilberdink,
\emph{Determinants of multiplicative Toeplitz matrices},
Acta Arith. \textbf{125} (2006), 265--284.
DOI: 10.4064/aa125-3-4.

\bibitem{Hilberdink2015}
T. Hilberdink,
\emph{Multiplicative Toeplitz matrices and the Riemann zeta function},
in \emph{Four Faces of Number Theory}, European Mathematical Society, 2015, 77--122.
DOI: 10.4171/142.

\bibitem{KannanBachem1979}
R. Kannan and A. Bachem,
\emph{Polynomial algorithms for computing the Smith and Hermite normal forms of an integer matrix},
SIAM J. Comput. \textbf{8} (1979), 499--507.
DOI: 10.1137/0208040.

\bibitem{Keiper1992}
J. B. Keiper,
\emph{Power series expansions of Riemann's $\xi$ function},
Math. Comp. \textbf{58} (1992), no. 198, 765--773.
DOI: 10.1090/S0025-5718-1992-1122072-5.

\bibitem{HaynesZudilin2015}
A. Haynes and W. Zudilin,
\emph{Hankel determinants of zeta values},
SIGMA \textbf{11} (2015), Paper 101, 5 pp.
DOI: 10.3842/SIGMA.2015.101.

\bibitem{Johansson2001}
K. Johansson,
\emph{Discrete orthogonal polynomial ensembles and the Plancherel measure},
Ann. of Math. (2) \textbf{153} (2001), no.~1, 259--296.
DOI: 10.2307/2661375.

\bibitem{Lagarias2007}
J. C. Lagarias,
\emph{Li coefficients for automorphic $L$-functions},
Ann. Inst. Fourier (Grenoble) \textbf{57} (2007), no. 5, 1689--1740.
DOI: 10.5802/aif.2311.

\bibitem{Li1997}
X.-J. Li,
\emph{The positivity of a sequence of numbers and the Riemann hypothesis},
J. Number Theory \textbf{65} (1997), no. 2, 325--333.
DOI: 10.1006/jnth.1997.2137.

\bibitem{Monien2009}
H. Monien,
\emph{Hankel determinants of Dirichlet series},
arXiv:0901.1883 [math.NT] (2009).

\bibitem{Riemann1859}
B. Riemann,
\emph{Ueber die Anzahl der Primzahlen unter einer gegebenen Gr\"osse},
Monatsberichte der Berliner Akademie (1859), 671--680.

\bibitem{Roesler1986}
F. Roesler,
\emph{Riemann's hypothesis as an eigenvalue problem},
Linear Algebra Appl. \textbf{81} (1986), 153--198.
DOI: 10.1016/0024-3795(86)90255-7.

\bibitem{Roesler1987}
F. Roesler,
\emph{Riemann's hypothesis as an eigenvalue problem. II},
Linear Algebra Appl. \textbf{92} (1987), 45--73.
DOI: 10.1016/0024-3795(87)90250-3.

\bibitem{Schoenberg1938}
I. J. Schoenberg,
\emph{Metric spaces and positive definite functions},
Trans. Amer. Math. Soc. \textbf{44} (1938), no. 3, 522--536.
DOI: 10.1090/S0002-9947-1938-1501980-0.

\bibitem{Suzuki2023Norms}
M. Suzuki,
\emph{Li coefficients as norms of functions in a model space},
J. Number Theory \textbf{252} (2023), 177--194.
DOI: 10.1016/j.jnt.2023.05.007.

\bibitem{Suzuki2023Screw}
M. Suzuki,
\emph{Aspects of the screw function corresponding to the Riemann zeta-function},
J. London Math. Soc. \textbf{108} (2023), no. 4, 1448--1487.
DOI: 10.1112/jlms.12785.

\bibitem{SuzukiScrewLine}
M. Suzuki,
\emph{The screw line of the Riemann zeta-function and its applications},
preprint, arXiv:2209.04658.

\bibitem{Suzuki2025Hilbert}
M. Suzuki,
\emph{On the Hilbert space derived from the Weil distribution},
Canad. J. Math. (published online 2025).
DOI: 10.4153/S0008414X25101739.

\bibitem{StahlTotik1992}
H. Stahl and V. Totik,
\emph{General Orthogonal Polynomials},
Encyclopedia of Mathematics and its Applications, vol. 43,
Cambridge University Press, Cambridge, 1992.
DOI: 10.1017/CBO9780511759420.

\bibitem{Totik2005Orthogonal}
V. Totik,
\emph{Orthogonal polynomials},
Surveys in Approximation Theory \textbf{1} (2005), 70--125.

\bibitem{WuLinXuZhao2016}
X.-B. Wu, Y. Lin, S.-X. Xu and Y.-Q. Zhao,
\emph{Uniform asymptotics for discrete orthogonal polynomials on infinite nodes with an accumulation point},
Anal. Appl. (Singap.) \textbf{14} (2016), no. 5, 705--737.
DOI: 10.1142/S0219530515500177.

\bibitem{Titchmarsh1986}
E. C. Titchmarsh,
\emph{The Theory of the Riemann Zeta-Function},
2nd ed., revised by D. R. Heath-Brown,
Oxford University Press, Oxford, 1986.

\bibitem{Vaughan1996}
R. C. Vaughan,
\emph{On the eigenvalues of Redheffer's matrix, II},
J. Austral. Math. Soc. Ser. A \textbf{60} (1996), 260--273.
DOI: 10.1017/S1446788700037654.


\bibitem{deBruijn1950}
N.~G. de Bruijn,
\emph{The roots of trigonometric integrals},
Duke Math. J. \textbf{17} (1950), no.~3, 197--226.
DOI: 10.1215/S0012-7094-50-01720-0.

\bibitem{Matiyasevich2007a}
Y.~Matiyasevich,
\emph{Hidden Life of Riemann's Zeta Function 1. Arrow, Bow, and Targets},
arXiv:0707.1983 [math.NT], 2007.

\bibitem{Matiyasevich2007b}
Y.~Matiyasevich,
\emph{Hidden Life of Riemann's Zeta Function 2. Electrons and Trains},
arXiv:0709.0028 [math.NT], 2007.

\bibitem{Newman1976}
C.~M. Newman,
\emph{Fourier transforms with only real zeros},
Proc. Amer. Math. Soc. \textbf{61} (1976), no.~2, 245--251.
DOI: 10.1090/S0002-9939-1976-0434982-5.

\bibitem{RodgersTao2020}
B.~Rodgers and T.~Tao,
\emph{The de Bruijn--Newman constant is non-negative},
Forum Math. Pi \textbf{8} (2020), e6.
arXiv:1801.05914.

\bibitem{Schoenfeld1976}
L.~Schoenfeld,
\emph{Sharper bounds for the Chebyshev functions \(\theta(x)\) and \(\psi(x)\). II},
Math. Comp. \textbf{30} (1976), no.~134, 337--360.

\bibitem{Weil1952}
A. Weil,
\emph{Sur les ``formules explicites'' de la th\'eorie des nombres premiers},
Medd. Lunds Univ. Mat. Sem. (1952), 252--265.

\end{thebibliography}

\end{document}
